% concatenated from main.tex, Sections/Intro.tex, Sections/1-preliminaries_0_categorical_background.tex, Sections/1-preliminaries_1_graphs.tex, Sections/1-preliminaries_2_modular_operads.tex, Sections/1-preliminaries_3_model_cat.tex, Sections/1-preliminaries_4_cyclic_operads.tex, Sections/1-preliminaries_mapping_class_group.tex, Sections/2-the-groupoid-bit.tex, Sections/2.1-modular_operad.tex, Sections/1-preliminaries_PaRB.tex, Sections/2.2-maps_out_of_genus0.tex, Sections/2.3-two_level.tex, Sections/2.4-maps_out_higher_genus_setup.tex, Sections/2.5-proof_of_maps_out_of_higher_genus.tex, Sections/2.6-proofs_of_two_level.tex, Sections/2.7-NS.tex, Sections/3-infinity_modular.tex, Sections/3-profinite_completion.tex, Sections/main_thms.tex, Sections/appendix_nerve.tex, Sections/appendix-complex_of_markings_over_pants.tex
% \pdfoutput=1 
\documentclass{amsart}
\usepackage{parskip}
\usepackage{lipsum}
\usepackage{tikz-cd}
\usepackage{hyperref}
\usepackage{graphicx,color}
\usepackage{stix}
\usepackage{amssymb,amsmath,bm,mathtools}
\usepackage{bm}
%\usepackage{mathrsfs}
\usepackage{enumerate}
\usepackage{tkz-graph}
\usepackage{ragged2e}
\usepackage{enumitem}
\usepackage{adjustbox}
\usepackage{tipa}
\usepackage{comment}
%========================================= Use tikz library ===========================================================


\usepackage{tikz}
\usepackage[all]{xy}

\usetikzlibrary{arrows,decorations.pathmorphing}
\usetikzlibrary{backgrounds,positioning}
\usetikzlibrary{fit,petri,shapes.misc}
\usetikzlibrary{shapes.geometric}
\usetikzlibrary{decorations.markings}
\usetikzlibrary{calc}

\newlength{\defaultpgflinewidth}
\setlength{\defaultpgflinewidth}{\pgflinewidth}

%=== edge label on top of edge
\tikzset{auto}
%=== node[swap]{?} to put label under edge

%=== empty vertex
\tikzset{empty/.style={circle,inner sep=0pt,minimum size=6mm}}
\tikzset{emptyvt/.style={circle,inner sep=0pt,minimum size=0mm}}

%=== circular vertex no color
\tikzset{plain/.style={circle,draw,very thick,
inner sep=0pt,minimum size=6mm}}

\tikzset{xplain/.style={circle,draw,very thick,
inner sep=0pt,minimum size=8mm}}

\tikzset{smallplain/.style={circle,draw,very thick,
inner sep=0pt,minimum size=4mm}}

\tikzset{xsplain/.style={circle,draw, thick,
inner sep=0pt,minimum size=1.5mm}}

\tikzset{tinyplain/.style={circle,draw, thick,
inner sep=0pt,minimum size=1mm}}

%====== filled vertices =========
\tikzset{smalldotted/.style={circle,draw,very thick, densely dotted,
inner sep=0pt,minimum size=3mm}}

\tikzset{dottedplain/.style={circle,draw,very thick, densely dotted,
inner sep=0pt,minimum size=6mm}}

\tikzset{normaldot/.style={circle,black,fill=black,inner sep=0pt,minimum size=2mm}}

%=== rectanglar vertex no color
\tikzset{rectplain/.style={rectangle,draw,very thick,minimum size=6mm}}
\tikzset{bigplain/.style={rectangle,draw,very thick,minimum size=1cm}}

%=== triangular vertex
\tikzset{triangular/.style={regular polygon, regular polygon sides=3, draw,very thick,
inner sep=0pt,minimum size=1.2cm}}

%===== edges 
\tikzset{arrow/.style={->,thick}}
\tikzset{dashedarrow/.style={->,dashed,thick}}
\tikzset{dottedarrow/.style={->,dotted,thick}}
\tikzset{mapto/.style={|->,thick}}
\tikzset{->-/.style={decoration={markings, mark=at position #1 with {\arrow{>}}},postaction={decorate}}}

%===== double arrow, like =>
\tikzset{implies/.style={thick,double,double equal sign distance,-implies}} 

%===== lines
\tikzset{line/.style={thick}}
\tikzset{dottedline/.style={dotted,thick}}
\tikzset{dashedline/.style={dashed,thick}}

%===== in/output leg
\tikzset{inputleg/.style={<-,thick}}
\tikzset{outputleg/.style={->,thick}}
\tikzset{dottedinput/.style={<-,dotted,thick}}

%===== end tikz macros =====





\usepackage{transparent}
\usepackage{graphicx,import}
\usepackage{caption}
\usepackage{subcaption}
\usepackage{enumitem}


\addtolength{\textwidth}{1.5 cm}
\addtolength{\hoffset}{-.75cm}
\addtolength{\textheight}{2cm}
\addtolength{\voffset}{-1cm}





\newcommand{\mrnote}[1]{\marginpar{\RaggedRight\tiny{\color{red}#1}}}
\newcommand{\lbnote}[1]{\marginpar{\RaggedRight\tiny{\color{blue}#1}}}



%================
%Notation
%================


%==============
%Fonts
%====================

\newcommand{\R}{\mathbb{R}}
\newcommand{\Z}{\mathbb{Z}}
\newcommand{\N}{\mathbb{N}}

\newcommand{\ba}{\mathbf{a}}
\newcommand{\bb}{\mathbf{b}}
\newcommand{\bc}{\mathbf{c}}

\newcommand{\B}{\mathcal{B}}
\newcommand{\F}{\mathsf{F}}
\newcommand{\bG}{\mathbb{G}}
\newcommand{\GG}{\mathcal{G}}

\newcommand{\CC}{\mathcal{C}}
\newcommand{\bC}{\mathbf{C}}

\newcommand{\OO}{\mathcal{O}}
\newcommand{\bO}{\mathbf{O}}
%\newcommand{\bE}{\mathbf{E}}
\newcommand{\bE}{\mathscr{E}}

\newcommand{\bM}{\mathscr{M}}
\newcommand{\bfC}{\mathbf{C}}

%\newcommand{\calE}{\mathcal{E}}
\newcommand{\calE}{\mathscr{E}}
\newcommand{\inftyE}{\mathtt{E}}
\newcommand{\calO}{\mathcal{O}}
\newcommand{\calQ}{\mathcal{Q}}
\newcommand{\calC}{\mathcal{C}}
\newcommand{\calS}{\mathcal{S}}
\newcommand{\bS}{\mathbf{S}}


\newcommand{\fC}{\mathfrak{C}}
\newcommand{\calM}{\mathcal{M}}
\newcommand{\pantscpx}{\mathcal{C}}
\newcommand{\xra}[1]{\overset{#1}{\rightsquigarrow}}
\newcommand{\surf}{S}%% Notation for a surface. S is bad since we already have S-moves and he modular operad S. F is bad because of F-moves. Sigma is bad because of the symmetric group.

\newcommand{\sG}{\mathsf{G}} % arbitrary groupoid 


\newcommand{\tG}{\text{G}} % arbitrary functor 
\newcommand{\tF}{\text{F}} % arbitrary functor 






%=====================
% Categories and Operads
%======================

%\newcommand{\cact}{\mathsf{cact}}
%\newcommand{\sSet}{\mathbf{sSet}}
\newcommand{\sSet}{\mathscr{S}}

\newcommand{\isSet}{\mathtt{S}}

\newcommand{\sS}{S}
\newcommand{\set}{\mathrm{Set}}
\newcommand{\Set}{\mathrm{Set}}
%\newcommand{\Grpd}{\mathbf{G}}
\newcommand{\Grpd}{\mathscr{G}}

\newcommand{\tree}{\mathsf{tree}}
\newcommand{\leg}{\mathsf{L}}
\newcommand{\edge}{\mathsf{E}}
\newcommand{\vertex}{\mathsf{V}}


\newcommand{\Gr}{\mathsf{Grp}}
\newcommand{\M}{\mathbb{M}}
\newcommand{\G}{\mathbf{G}}
\newcommand{\hbS}{\widehat{\mathbf{S}}}
\newcommand{\E}{\mathsf{E}}
\newcommand{\PaRB}{\mathsf{PaRB}}
\newcommand{\PaB}{\mathsf{PaB}}
\newcommand{\hPP}{\widehat{\PP}}
\newcommand{\PP}{\mathcal{P}}
\newcommand{\markS}{\mathcal{M}(S)}
\newcommand{\calP}{\mathcal{P}}
\newcommand{\bP}{\mathbb{P}}
\newcommand{\hGamma}{\hat{\Gamma}}
\newcommand{\Top}{\mathrm{Top}}
\newcommand{\Inn}{\mathrm{Inn}}
\newcommand{\onen}[1][]{{\id_{#1}}}
\newcommand{\twist}{\tau}
\newcommand{\dehn}{\mathbf{\mathrm{D}}^{1/2}}
\newcommand{\galQ}{\mathbf{\mathrm{Gal}}(\overline{\mathbb{Q}}/\mathbb{Q})}
\newcommand{\pants}{\mathbf{P}}



\newcommand{\Mod}{\mathrm{Mod}}
\newcommand{\dMod}{\mathrm{dMod}}

\newcommand{\dCyc}{\mathrm{dCyc}}
\newcommand{\Cyc}{\mathrm{Cyc}}
\newcommand{\Op}{\mathrm{Op}}
\newcommand{\dOp}{\mathrm{dOp}}
\newcommand{\op}{\mathrm{op}}

\newcommand{\Sc}{\mathrm{Sc}}

\newcommand{\widesthat}[1]{(#1)^{\wedge}}
%=======================
%Functors
%++++++++++++++++++++++
\newcommand{\forget}{\mathbf{r}}
\newcommand{\Pre}{\operatorname{Pre}}


%====================
%Operations
%====================

\newcommand{\id}{\operatorname{id}}
\newcommand{\Id}{\operatorname{Id}}
\newcommand{\Iso}{\operatorname{Iso}}
\DeclareMathOperator*{\colim}{colim}
\newcommand{\Pro}{\operatorname{Pro}}
%for the classifying space functor to look nicer
\newcommand{\cs}{\mathrm{B}}
\newcommand{\nerve}{\mathrm{N}}


\newcommand{\End}{\operatorname{End}}
\newcommand{\Hom}{\operatorname{Hom}}
\newcommand{\Ho}{\operatorname{Ho}}
\newcommand{\Mor}{\operatorname{Mor}}
\newcommand{\ob}{\operatorname{ob}}
\newcommand{\Aut}{\operatorname{Aut}}
\newcommand{\Emb}{\operatorname{Emb}}
\newcommand{\Out}{\operatorname{Out}}
\newcommand{\out}{\operatorname{out}}
\newcommand{\inpt}{\operatorname{in}}
\newcommand{\Map}{\operatorname{Map}}
\newcommand{\Alg}{\operatorname{Alg}}
\newcommand{\RMap}{\mathbb{R}\operatorname{Map}}
\newcommand{\HoEnd}{\operatorname{HoEnd}}
\newcommand{\HoAut}{\operatorname{HoAut}}

\newcommand{\Diff}{\operatorname{Diff}}
\newcommand{\Homeo}{\operatorname{Homeo}}
\newcommand{\del}{\partial}
\newcommand{\diff}{\mathrm{Diff}}
\newcommand{\tr}{\mathrm{tr}}
\newcommand{\trun}[1]{\mathrm{tr}_{\leq{#1}}}

\newcommand{\freegraphmonad}{\top}



%-===========Really Wide Hat =====

\usepackage{scalerel,stackengine}
\stackMath
\newcommand\reallywidehat[1]{%
\savestack{\tmpbox}{\stretchto{%
  \scaleto{%
    \scalerel*[\widthof{\ensuremath{#1}}]{\kern-.6pt\bigwedge\kern-.6pt}%
    {\rule[-\textheight/2]{1ex}{\textheight}}%WIDTH-LIMITED BIG WEDGE
  }{\textheight}% 
}{0.5ex}}%
\stackon[1pt]{#1}{\tmpbox}%
}
\parskip 1ex
%=============

%============
%Groups
%================
\newcommand{\GT}{\widehat{\mathsf{GT}}}
\newcommand{\NS}{\widehat{\mathsf{\Gamma}}}
\newcommand{\uNS}{\underline{\widehat{\mathsf{\Gamma}}}}
\newcommand{\uGT}{\underline{\widehat{\mathsf{GT}}}}
\newcommand{\LL}{\Lambda_{\mathrm{NS}}}
\newcommand{\mLL}{\underline{\LL}}
\newcommand{\Br}{\mathsf{B}} %braid groups
\newcommand{\PB}{\mathsf{PB}} %braid groups

\newcommand{\RB}{\mathsf{RB}} %ribbon braid groups
\newcommand{\PRB}{\mathsf{PRB}} %pure ribbon braid groups

\newcommand{\Free}{\mathsf{F}} %free groups
\newcommand{\Mmax}{\mathcal{M}^{max}}




%====================
%Trees and Graphs
%==============

\newcommand{\ie}{iE}%{Ed_{i}}
\newcommand{\ed}{Ed}
\newcommand{\nb}{nb}

\newcommand{\om}{\Omega}
\newcommand{\bU}{\mathbf{U}}

\newcommand{\corolla}{C}

\renewcommand{\S}{\bm{\mathrm{S}}}
\newcommand{\pNG}{\widehat{\N\G}}
\renewcommand{\NG}{\N\G}
\newcommand{\pG}{\widehat{\bm{\mathrm{G}}}}
\newcommand{\pPP}{\widehat{\PP}}

\newcommand{\gentwoone}{\raisebox{-11pt}{\includegraphics[scale=0.9]{Figures/gen2-1.pdf}}}
\newcommand{\gentwotwo}{\raisebox{-15pt}{\includegraphics[scale=0.9]{Figures/gen2-2.pdf}}}
\newcommand{\gentwothree}{\raisebox{-15pt}{\includegraphics[scale=0.5]{Figures/gen2-3.pdf}}}
\newcommand{\gentwofour}{\raisebox{-15pt}{\includegraphics[scale=0.5]{Figures/gen2-4.pdf}}}
\newcommand{\triv}{\raisebox{-2pt}{\includegraphics[height=14pt]{Figures/triv.pdf}}}
\newcommand{\forkone}{\raisebox{-4pt}{\includegraphics[height=8pt]{Figures/fork1.pdf}}}
\newcommand{\forktwo}{\raisebox{-4pt}{\includegraphics[height=8pt]{Figures/fork2.pdf}}}
\newcommand{\contraction}{\raisebox{-4pt}{\includegraphics[height=8pt]{Figures/contraction.pdf}}}
\newcommand{\graphgtwo}{\raisebox{-4pt}{\includegraphics[height=8pt]{Figures/graphg2.pdf}}}


%================
%Environments
%===============
%\newcounter{dummy} 
%\numberwithin{dummy}{section}
\newtheorem{theorem}{Theorem}[section]
\newtheorem{thm}[theorem]{Theorem}%[section]
\newtheorem{prop}[theorem]{Proposition}%[section]
\newtheorem*{prop*}{Proposition}%[section]
\newtheorem{lemma}[theorem]{Lemma}%[section]
\newtheorem{cor}[theorem]{Corollary}%[section]
\newtheorem*{cor*}{Corollary}%[section]
\newtheorem{claim}[theorem]{Claim}%[section]
\newtheorem{Idea}[theorem]{Idea}%[section]
\newtheorem*{thm*}{Theorem}
\newtheorem*{construction}{Construction}



\newtheorem{Th}{Theorem}\renewcommand{\theTh}{\Alph{Th}}

\theoremstyle{definition}
\newtheorem{definition}[theorem]{Definition}%[section]
\newtheorem{cons}[theorem]{Construction}%[section]
\newtheorem{example}[theorem]{Example}%[section]
\newtheorem{remark}[theorem]{Remark}%[section]
\newtheorem{notation}[theorem]{Notation}
%\newtheorem{notation}[dummy]{Notation}
%\newtheorem{que}[dummy]{Question}

\numberwithin{equation}{section}

\setcounter{tocdepth}{2}% set 3 to get subsubsections in toc
\let\oldtocsection=\tocsection
\let\oldtocsubsection=\tocsubsection
\renewcommand{\tocsection}[2]{\hspace{0em}\oldtocsection{#1}{#2}}
\renewcommand{\tocsubsection}[2]{\hspace{1em}\oldtocsubsection{#1}{#2}}
%\renewcommand{\tocsubsubsection}[2]{\hspace{2em}\oldtocsubsubsection{#1}{#2}}
\DeclareRobustCommand{\SkipTocEntry}[5]{}
\renewcommand{\paragraph}[1]{\textbf{{#1}.}\hspace{5pt}}

\newenvironment{todo}
    {\color{red}}
    
\newenvironment{idea}
    {\color{blue} \begin{Idea}}
    {\end{Idea}}

\title[Galois actions on surfaces]{Galois actions on surfaces and a higher genus Grothendieck-Teichm\"{u}ller group}


\author[L. Basualdo Bonatto]{Luciana Basualdo Bonatto}
\address{Mathematical Institute, University of Oxford \\ Oxford, UK}
\email{luciana.basualdobonatto@maths.ox.ac.uk}

\author[M. Robertson]{Marcy Robertson}
\address{School of Mathematics and Statistics \\ The University of Melbourne \\ Melbourne, Victoria, Australia}
\email{marcy.robertson@unimelb.edu.au}

\date{\today}

\begin{document}


\begin{abstract}
We construct an operadic model for the higher-genus Teichm\"uller tower. More precisely, we define a modular operad $\bS$ in groupoids built from mapping class groups, with compositions and contractions encoding gluing operations on surfaces. We prove a presentation theorem for maps out of $\bS$, showing that they are determined by a small number of genus-zero and genus-one generators and relations. Using this presentation and the work of Nakamura--Schneps, we construct a faithful action of the Nakamura--Schneps subgroup $\NS\subseteq\GT$ on the profinite completion $\widehat{\bS}$, and hence an action of $\operatorname{Gal}(\overline{\mathbb Q}/\mathbb Q)$. The genus-zero truncation of $\bS$ recovers the cyclic operad of parenthesized ribbon braids, and its group of object-fixing profinite automorphisms recovers $\GT$. Finally, the profinite completion of the classifying spaces of $\bS$ assemble into a modular $\infty$-operad in profinite spaces whose values identify with the \'etale homotopy types of moduli stacks of curves with marked tangent vectors, and the $\NS$-action extends to this homotopy-coherent Teichm\"uller tower.
\end{abstract}

%\input{Sections/outline}\newpage

\maketitle
\section{Introduction}

In the essay \emph{Esquisse d'un Programme}, Grothendieck proposed that the absolute Galois group of $\mathbb Q$, $\galQ$, could be studied through its action on a ``tower'' of geometric objects, namely \'etale fundamental groups of the moduli spaces $\mathcal M_{g}^{n}$ , together with the geometric morphisms relating these spaces \cite{grothendieck_esquisse}. These morphisms record that the standard action of $\galQ$ is compatible with maps such as forgetting marked points, gluing curves along marked points or tangent directions, and passing to boundary strata.  Already in genus zero this action is highly nontrivial:  Belyi's theorem implies that the action of $\galQ$ on $\pi_1^{\acute et}(\mathcal M_0^4)$ is faithful \cite{belyi}. Grothendieck's guiding idea was that there should be an ``ideal Teichm\"uller tower'' containing enough of these compatibility maps such that the automorphisms of the tower would recover precisely $\galQ$. The main purpose of this paper is to construct a topological model for a Teichm\"uller tower as a modular operad.

For $2g-2+n>0$, the \'etale fundamental group of $\mathcal M_g^n$ identifies with the profinite completion of the mapping class group $\Gamma_g^n$ of a genus $g$ surface with $n$ marked points fixed pointwise \cite{matsumoto2000arithmetic,boggi2009fundamental}. The action of $\galQ$ then gives an outer action on the profinite completion of $\Gamma_g^n$. Passing from marked points to boundary components corresponds topologically to remembering tangent directions at the marked points, and this gives corresponding actions on the profinite completions of the mapping class groups $\Gamma_{g,n}$ of genus $g$ surfaces with $n$ boundary components fixed pointwise (Definition~\ref{def: mapping class group}). Thus the tower of \'etale fundamental groups has a topological counterpart built from mapping class groups and compatibility maps coming from topological operations such as cutting along simple closed curves, gluing boundary components, and forgetting marked points. 

This approach via mapping class groups provides the topological framework for the modular operad (Definition~\ref{def: modular operad}) constructed in this paper. Namely, we define a family of groupoids $\bS(g,n+1)$, for $2g+n\geq 1$, satisfying
    \[\cs\mathbf S(g,n+1)\simeq \cs\Gamma_{g,n+1}.\]
The groupoids $\mathbf S(g,n+1)$ assemble into a modular operad in groupoids, denoted $\mathbf S$ (Definition~\ref{def: surface modular operad}).  Its compositions and contractions encode the maps induced by gluing boundary components. After profinite completion, we obtain a modular operad $\widehat{\bS}$ in profinite groupoids (see Section~\ref{sec: profinite completion of modular operads in groupoids}). This construction supports the Galois action envisioned by Grothendieck:

\begin{thm*}[Theorem~\ref{main theorem NS action}]
The absolute Galois group $\galQ$ acts faithfully on the modular operad
$\widehat{\mathbf S}$, via an explicit homomorphism
\[
\galQ\longrightarrow \operatorname{Aut}(\widehat{\mathbf S}).
\]
\end{thm*}

A substantial part of the paper is devoted to constructing this action and relating it to the classical genus-zero Grothendieck--Teichm\"uller theory.

The genus-zero part of the Teichm\"uller tower was first studied in depth by Drinfeld \cite[Section~4]{Drin}. He considered a tower of profinite braid groups whose structure maps include the standard operations of forgetting strands and replacing one strand by several parallel strands. Drinfeld introduced the profinite Grothendieck--Teichm\"uller group $\GT$ via the symmetries of this tower (Definition~\ref{defn:GT}). Ihara showed that there is an embedding
    \[\galQ\hookrightarrow \GT,\]
and that the induced action on Drinfeld's tower of profinite braid groups agrees with the Galois actions arising from the geometric \'etale fundamental groups
of the genus-zero moduli spaces \cite{ihara}. 

Operadic models for the genus-zero Teichm\"uller tower have been studied in, for instance, \cite{FresseBook1,BN,willwacher2015m,Horel_profinite_groupoids,Boavida-Horel-Robertson,robertson2025grothendieck}. The version most directly related to the surfaces used in this paper is the cyclic operad of parenthesized ribbon braids $\PaRB^{\mathrm{cyc}}$, introduced in \cite{campos2019configuration}. Its cyclic structure reflects the fact that, for genus-zero surfaces, all boundary components can be treated symmetrically rather than singling out one as an output.

Our modular operad $\bS$ extends this genus-zero surface model to higher genus. More precisely, its genus-zero truncation $\trun{0}\bS$, defined in Appendix~\ref{subsec: truncations}, is naturally a cyclic operad and satisfies the following comparison.


\begin{prop*}[Proposition~\ref{prop:PaRB-cyclic-is-S0}]
There is an isomorphism of cyclic operads
    \[\operatorname{tr}_{\leq 0}\mathbf S \cong \PaRB^{\mathrm{cyc}}. \]
\end{prop*}

\medskip

The key technical input for constructing actions on $\widehat{\bS}$ is a presentation theorem for maps out of $\bS$. This is the higher-genus analogue of the presentation of maps out of genus-zero operadic models such as parenthesized ribbon braids. In the genus-zero case, an operadic map out of $\PaRB$ is determined by an object together with morphisms representing a twist, a braiding, and an associator satisfying certain relations \cite{FresseBook1,Boavida-Horel-Robertson}. In the higher-genus setting we obtain the following analogue.

\begin{thm*}[Theorem\ref{thm: presentation for maps out of bS}]
    Let $\calQ$ be a modular operad in groupoids. A map of modular operads $f:\bS\longrightarrow \calQ$ is uniquely determined by a tuple
        \[ (\mathbf{m},\mathbf{t},\mathbf{b},\mathbf{a},\mathbf{s}) \]
consisting of an object $\mathbf{m}\in \ob(\calQ(0,3))$ 
and morphisms $\mathbf{t}$ (twist), $\mathbf{b}$ (braiding), $\mathbf{a}$ (associator), and $\mathbf{s}$ (genus-one switch), satisfying the standard genus-zero relations and three additional genus-one relations.
\end{thm*}

The proof of this presentation theorem uses complexes of marked pants decompositions on surfaces. More precisely, we compare the groupoids $\bS(g,n+1)$ with the fundamental groupoids of CW complexes pants decompositions equipped with additional marking data (Definition~\ref{def: complex of markings}). These complexes refine the Hatcher--Thurston pants-decomposition complex \cite{Hatcher_Thurston} by adding marking data similar to that of \cite{bk_marked_surfaces}. Their connectedness and simple connectedness, proved in Appendix~\ref{app: complex of markings}, allow us to read off generators and relations for maps out of $\bS$.

Using the presentation theorem, we construct the action on $\widehat{\bS}$ from the compatible actions of Nakamura--Schneps on profinite mapping class groups. In \cite{Nakamura-Schneps}, Nakamura and Schneps introduced a subgroup $\NS$ of $\GT$ by imposing additional genus-one relations on elements of $\GT$. They proved that $\NS$ contains the image of $\galQ$, giving inclusions
    \[\galQ\hookrightarrow \NS\hookrightarrow \GT,\]
and it remains open whether either inclusion is strict \cite[Theorem~1.2]{Nakamura-Schneps}. They also constructed compatible actions of $\NS$ on profinite mapping class groups of surfaces equipped with additional structure, extending the genus-zero action of $\GT$ \cite[Theorem~1.4]{Nakamura-Schneps}. Our construction packages these actions operadically: the presentation theorem reduces the verification to the defining genus-zero and genus-one relations, and hence produces a faithful action of $\NS$ on $\widehat{\bS}$.


Restricting to genus zero recovers the usual Grothendieck--Teichm\"uller symmetry. More precisely, we obtain the following compatibility with the cyclic genus-zero model.

\begin{prop*}[Proposition~\ref{prop:genus-zero-GT}]
There is an isomorphism of profinite groups
    \[\GT\cong \operatorname{Aut}_0(\operatorname{tr}_{\leq 0}\widehat{\mathbf S})\]
where $\operatorname{Aut}_0(-)$ denotes automorphisms fixing the objects of the underlying groupoids.
\end{prop*}

Thus the Grothendieck--Teichm\"uller action respects the operation of treating all boundary components of a genus-zero surface symmetrically, rather than distinguishing one boundary component as an output. This is compatible with Kontsevich's discussion of the framed genus-zero Teichm\"uller tower and its higher-genus modular-operadic analogue \cite{KontBourbaki}, and is closely related to the work of the second author and Singh on cyclic parenthesized ribbon braids and framed tangles \cite{robertson2025grothendieck}.


The presentation theorem also gives an operadic form of Grothendieck's \emph{two-level principle}: the full tower should be generated from its genus-zero and genus-one parts. In the present setting, the generators and defining relations occur already in the levels corresponding to
    \[\Gamma_{0,4},\qquad \Gamma_{0,5},\qquad \Gamma_{1,1},\qquad \Gamma_{1,2}.\]
Equivalently, if $\trun{1}\bS$ denotes the genus-one truncation of $\bS$, defined in Appendix~\ref{subsec: truncations}, then $\bS$ is freely recovered from this truncation.

\begin{cor*}[Corollary~\ref{cor: maps out of bS}]
There is an isomorphism of modular operads
    \[\bS\cong (\trun{1})_!(\trun{1}\bS),\]
where $(\trun{1})_!$ denotes the left adjoint to the genus-one truncation functor.
\end{cor*}

\medskip

In the final part of the paper, we pass from profinite groupoids to profinite spaces in order to compare our construction with the \'etale homotopy types of moduli spaces. This gives a homotopical refinement of the Teichm\"uller tower: it retains the \'etale fundamental groups while also remembering higher homotopy information. We use the modern $\infty$-categorical formulation of \'etale homotopy types as shapes of \'etale $\infty$-topoi, together with their profinite refinements. This perspective is compatible with the classical constructions of Artin--Mazur \cite{Artin_Mazur} and Friedlander \cite{Friedlander}; see also \cite[Definition~7.1.6.1 and Example~7.1.6.9]{Lurie}, \cite{Carchedi}, and \cite{haine2024profinite}.

Oda \cite[Theorem 1]{Oda} gives a concrete description of the profinite \'etale homotopy types of the moduli spaces $\mathcal M_g^n$. As recalled in Example~\ref{example: Oda}, for $2g-2+n>0$ there is an equivalence in profinite spaces
\[
\Pi^{\acute et}_\infty(\mathcal M_g^n\otimes_{\mathbb Z}\overline{\mathbb Q})
\simeq
\widehat{\cs\Gamma_g^n}.
\]
Thus, in this case, the \'etale homotopy type is identified with the profinite completion of the classifying space of a mapping class group.

For our topological model of the Teichm\"uller tower, we work with boundary mapping class groups $\Gamma_{g,n}$. On the algebro-geometric side, this corresponds to replacing marked points by marked points equipped with non-zero tangent vectors. We denote by $\mathcal M_{g,n}^{\tan}$ the moduli stack of smooth curves with $n$ ordered marked points equipped with non-zero tangent vectors. The corresponding form of Oda's comparison identifies its profinite \'etale homotopy type with the classifying space of the completed boundary mapping class group:
\[
\Pi^{\acute et}_\infty(\mathcal M_{g,n}^{\tan}\otimes_{\mathbb Z}\overline{\mathbb Q})
\simeq
\widehat{\cs\Gamma_{g,n}}.
\]
We recall this comparison in Lemma~\ref{lemma:tangent-moduli-comparison}. Using the equivalence $\cs\mathbf S(g,n+1)\simeq \cs\Gamma_{g,n+1}$, this identifies $\widehat{\cs\mathbf S}(g,n+1)$ with the corresponding \'etale homotopy type of $\mathcal M_{g,n+1}^{\tan}$.

It is not formal, however, that these profinite spaces assemble into a modular operad. The classifying space functor preserves products, so the collection of spaces $\cs\bS(g,n+1)$ inherit a strict modular operad structure from $\bS$. After profinite completion, this argument breaks down because profinite completion of spaces rarely preserves products \cite[Proposition~3.9]{Boavida-Horel-Robertson}. This is a problem because the composition laws in a modular operad are encoded by products over the vertices of graphs (Section~\ref{sec: modular operads}).

To address this, we use modular $\infty$-operads, a homotopy-coherent version of modular operads. More precisely, we use a genus-graded version of the Segal model developed by Hackney, the second author, and Yau in \cite{hry1}. In this model, a modular $\infty$-operad is a contravariant functor on the modular dendroidal category of genus-graded graphs satisfying Segal equivalences over the vertices of each graph. The required compatibility with products is thus imposed up to homotopy, rather than as a strict product-preservation statement.


Applying this formalism to the profinite classifying spaces of $\bS$, we obtain a modular $\infty$-operad in profinite spaces whose values on corollas identify with the \'etale homotopy types of the moduli stacks $\mathcal M_{g,n+1}^{\tan}$. Its values on more general graphs encode the corresponding gluing data through the modular Segal conditions. After applying classifying spaces and passing to derived automorphisms, the strict $\NS$-action on $\widehat{\bS}$ induces an action in the homotopy category of modular $\infty$-operads in profinite spaces.


\begin{thm*}[Theorem~\ref{NS action on topological operad}]
The profinite classifying spaces $\widehat{\cs(\bS(g,n+1))}$, for $2g+n\geq1$, assemble into a modular $\infty$-operad $\widehat{\cs\nerve\bS}$ in profinite spaces. The action of $\NS$ on $\widehat{\bS}$ induces a faithful action of $\NS$
on $\widehat{\cs\bS}$ in the homotopy category, via an explicit injective homomorphism
    \[\NS\longrightarrow \pi_0\mathbb R\operatorname{Aut}(\widehat{\cs\bS}).\]
\end{thm*}

Restricting to genus zero recovers the isomorphism:
\[
\GT\cong \pi_0\mathbb R\operatorname{Aut}
(\trun{0}\widehat{\cs\nerve\bS}).
\]


\subsection{Acknowledgments}
The authors would like to thank Pedro Boavida de Brito, Philip Hackney, Peter Hain, Geoffroy Horel, and Jan Steinebrunner for many useful conversations and comments.  M.R. recognizes the support of Australian Research Council Future Fellowship FT210100256. L.B.B. is grateful to the Max Planck Institute for Mathematics in Bonn for its hospitality and financial support.

\begin{comment}
Let $\mathcal M_{g}^{n}:=\mathcal M_{g}^{n}\otimes_{\mathbb Q}\overline{\mathbb Q}$ denote the moduli stack of smooth genus $g$ curves with $n$ marked points, defined over $\mathbb Q$, and let $\galQ$ denote the absolute Galois group of $\mathbb Q$.  The geometric \'etale fundamental group $\pi_1^{\acute et}(\mathcal M_{g}^{n})$ carries an outer action of $\galQ$. Already in genus zero this action is highly nontrivial.  Belyi's theorem implies that the action of $\galQ$ on $\pi_1^{\acute et}(\mathcal M_0^4)\cong \widehat{\F}_2$ is faithful \cite{belyi}.  Motivated by this example, Grothendieck proposed in his \emph{Esquisse d'un Programme} that $\galQ$ should be studied through its action on a ``tower'' of \'etale fundamental groupoids of the moduli spaces $\mathcal M_{g}^{n}$, together with the geometric morphisms relating these spaces \cite{grothendieck_esquisse}.  These morphisms include forgetting marked points, gluing curves along marked points or tangent directions, and passing to boundary strata.  Grothendieck's guiding idea was that by forcing compatibility with the geometric operations which make up the ``Teichm\"uller tower'' should impose strong constraints on the Galois action.

The main purpose of this paper is to construct a topological model for the Teichm\"uller tower which retains both the mapping class groups and the geometric operations relating them.  For $2g-2+n>0$, the \'etale fundamental group of $\mathcal M_g^n$ identifies with the profinite completion of the mapping class group $\Gamma_g^n$ of a genus $g$ surface with $n$ marked points fixed pointwise \cite{matsumoto2000arithmetic,boggi2009fundamental}.  Thus the tower of \'etale fundamental groups has a topological counterpart built from mapping class groups.  The geometric morphisms between moduli spaces correspond topologically to operations on surfaces, such as cutting along simple closed curves, gluing boundary components, and forgetting marked points.  These surface operations are the input for the modular-operadic model constructed below.

To model the cutting and gluing operations topologically, we first pass from curves with marked points to surfaces with parametrized boundary components. In practice, this means replacing the mapping class groups $\Gamma_g^n$ by the mapping class groups $\Gamma_{g,n}$.  Here, $\Gamma_{g,n}$ denotes the mapping class group of a genus $g$ surface with $n$ boundary components fixed pointwise (Definition~\ref{def: mapping class group}), then filling in the boundary components gives a central extension
\[
1\longrightarrow \mathbb Z^n\longrightarrow \Gamma_{g,n}\longrightarrow
\Gamma_g^n\longrightarrow 1,
\]
where the kernel is generated by the boundary Dehn twists.  Thus replacing marked points by boundary components amounts to retaining tangent directions at the marked points.  In this boundary setting, gluing parametrized boundary components induces homomorphisms between mapping class groups.  On moduli spaces, these operations correspond to clutching morphisms whose images lie in boundary strata of the Deligne--Mumford compactification.  This point of view is closely related to the work of Hatcher, Lochak and Schneps, who construct a topological tower using homomorphisms associated to inclusions of subsurfaces obtained by cutting along disjoint simple closed curves \cite{hls}.


In Section~\ref{sec: modular operad of mapping class groups}, we introduce a
family of groupoids $\mathbf S(g,n+1)$.  The objects are connected, labelled,
trivalent graphs of genus $g$ with $n+1$ legs.  Such a graph $G$
determines a labelled parametrized surface $S_G$ by replacing each vertex with
a pair of pants and gluing boundary components according to the internal edges
of $G$.  The morphisms $G\to H$ are isotopy classes of
boundary-label-preserving diffeomorphisms $S_G\to S_H$.  Thus the automorphism
group of any object is naturally isomorphic to $\Gamma_{g,n+1}$, and we have
an equivalence of classifying spaces:
\[
\cs\mathbf S(g,n+1)\simeq \cs\Gamma_{g,n+1}.
\]

The groupoids $\mathbf S(g,n+1)$ assemble into a modular operad in groupoids,
denoted $\mathbf S$.  Its partial compositions and contractions are induced
by gluing parametrized boundary components.  To obtain a finite presentation of
this modular operad, we compare the groupoids $\mathbf S(g,n+1)$ with
two-dimensional complexes of markings over fixed surfaces.  These complexes are
modelled on the Hatcher--Thurston pants-decomposition complex
\cite{Hatcher_Thurston}, with additional marking data similar to that of
\cite{bk_marked_surfaces}.  Their vertices are marked pants decompositions,
their edges are twists, braidings, and $A$- and $S$-moves, and their
two-cells encode the relations among these moves.  We prove in
Appendix~\ref{app: complex of markings} that the complexes of markings are
connected and simply connected.  This gives a finite presentation compatible
with categorical composition, modular operadic composition, contractions, and
symmetric group actions.  The resulting presentation theorem for maps out of
$\mathbf S$ is Theorem~\ref{thm: presentation for maps out of bS}.


We can profinitely complete the modular operad $\mathbf S$, as explained in Section~\ref{sec: profinite completion of modular operads in groupoids}.  We then compare $\widehat{\mathbf S}$ with the arithmetic symmetries already known in low genus.  The genus-zero comparison involves the profinite Grothendieck--Teichm\"uller group, while the higher-genus action will be constructed using the Nakamura--Schneps subgroup introduced below.

In \cite[Section~4]{Drin}, Drinfeld considered a tower of profinite braid groups whose structure maps include the standard operations of forgetting strands and replacing one strand by several parallel strands.  He introduced the profinite Grothendieck--Teichm\"uller group $\GT$ as a group of symmetries of this tower (Definition~\ref{defn:GT}).  Lochak and Schneps later proved in \cite{LS_braid} that $\GT$ is isomorphic to the group of outer automorphisms of this profinite braid tower compatible with these strand operations.  Horel interpreted this point of view operadically in \cite{Horel_profinite_groupoids}, building on operadic descriptions of prounipotent Grothendieck--Teichm\"uller groups such as \cite{FresseBook1,BN,willwacher2015m}.

Ihara showed that there is an embedding
\[
\galQ\hookrightarrow \GT,
\]
and that the induced action on Drinfeld's tower of profinite braid groups agrees
with the Galois actions arising from the geometric \'etale fundamental groups
of the genus-zero moduli spaces \cite{ihara}.  Since braid groups are closely
related to genus-zero mapping class groups \cite{birman1969mapping}, Drinfeld's
braid tower gives an explicit genus-zero model for the Teichm\"uller tower
proposed by Grothendieck.  The comparison becomes even more direct after
passing to framed, or ribbon, braid groups.  The genus-zero mapping class groups
with boundary are identified with pure ribbon braid groups,
\[
\Gamma_{0,n+1}\cong \PRB_n,
\]
as recalled in Section~\ref{sec: operad of ribbon braids}.

We use this comparison to identify the genus-zero part of our modular operad.
The genus-zero truncation $\operatorname{tr}_{\leq 0}\mathbf S$, defined in
Appendix~\ref{subsec: truncations}, discards all positive-genus operations and
is naturally a cyclic operad.  The relationship between genus-zero mapping class
groups and ribbon braid groups extends to an isomorphism with the cyclic operad
of parenthesized ribbon braids $\PaRB^{\mathrm{cyc}}$, introduced in
\cite{campos2019configuration}.

\begin{prop*}[Proposition~\ref{prop:PaRB-cyclic-is-S0}]
There is an isomorphism of cyclic operads
\[
\operatorname{tr}_{\leq 0}\mathbf S \cong \PaRB^{\mathrm{cyc}}.
\]
\end{prop*}

After profinite completion, this comparison identifies the automorphisms of the
genus-zero part of $\widehat{\mathbf S}$ with the profinite
Grothendieck--Teichm\"uller group.

\begin{prop*}[Proposition~\ref{prop:genus-zero-GT}]
There is an isomorphism of profinite groups
\[
\GT\cong \operatorname{Aut}_0(\operatorname{tr}_{\leq 0}\widehat{\mathbf S}).
\]
\end{prop*}

This shows that the description of $\GT$ as the group of object-fixing automorphisms of the profinite operad of parenthesized ribbon braids \cite{Boavida-Horel-Robertson} is compatible with the cyclic structure on the genus-zero truncation of $\bS$.  In particular, the Grothendieck--Teichm\"uller action respects the operation of treating all boundary components of a genus-zero surface symmetrically, rather than distinguishing one boundary component as an output.  This fits with Kontsevich's discussion of the framed genus-zero Teichm\"uller tower and its higher-genus modular-operadic analogue \cite{KontBourbaki}.  It is also closely related to the work of the second author and Singh on cyclic parenthesized ribbon braids and framed tangles \cite{robertson2025grothendieck}.


In \cite{Nakamura-Schneps}, Nakamura and Schneps introduced a subgroup $\NS\subseteq \GT$ by imposing additional genus-one relations on elements of $\GT$. These relations reflect the behaviour of the Galois action on Dehn twist generators in genus one.  The subgroup $\NS$
contains the image of $\galQ$, giving inclusions \cite[Theorem 1.2]{Nakamura-Schneps}
\[
\galQ\hookrightarrow \NS\hookrightarrow \GT.
\]

\begin{thm*}[Theorem~\ref{main theorem NS action}]
The Nakamura--Schneps group $\NS$ acts by object-fixing automorphisms on the
profinite modular operad $\widehat{\mathbf S}$.  Equivalently, there is a
homomorphism of profinite groups
\[
\NS\longrightarrow \operatorname{Aut}_0(\widehat{\mathbf S}).
\]
Consequently, the absolute Galois group $\galQ$ acts on
$\widehat{\mathbf S}$.
\end{thm*}

The proof of this result uses an operadic form of Grothendieck's two-level principle.  

In this paper, the two-level principle arises as a statement about maps out of the modular operad $\mathbf S$.  Theorem~\ref{thm: presentation for maps out of bS} shows that such maps are determined by their values on the genus-zero and genus-one generators, subject to relations already occurring in
\[
\mathbf S(0,4),\qquad \mathbf S(0,5),\qquad
\mathbf S(1,1),\qquad \mathbf S(1,2).
\]
After profinite completion, this gives the genus-one truncation statement of Corollary~\ref{cor: genus one truncation}.  The construction of the $\NS$-action is therefore reduced to checking precisely the relations in the Nakamura--Schneps definition.  Since $\galQ$ maps into $\NS$, this gives a Galois action on the profinite modular operad $\widehat{\mathbf S}$.


In the second half of the paper, we pass from profinite groupoids to profinite spaces in order to compare this construction with the \'etale homotopy types of the moduli spaces.  This is the homotopical refinement of the Teichm\"uller tower. It retains the \'etale fundamental groups, but also remembers the higher homotopy information of the \'etale site.  The \'etale homotopy type was introduced by Artin--Mazur as a pro-object in the homotopy category of spaces \cite{Artin_Mazur} and refined by Friedlander using pro-simplicial sets \cite{Friedlander}.  We use the $\infty$-categorical formulation in which the \'etale homotopy type is the shape of the \'etale $\infty$-topos.  In this formulation it is a pro-object in the  $\infty$-category of spaces, using Lurie's language of pro-objects in an $\infty$-category \cite[Definition~7.1.6.1]{Lurie}.  The interpretation of the \'etale homotopy type as the shape of the \'etale $\infty$-topos is discussed in \cite[Example~7.1.6.9]{Lurie}; see also \cite{Carchedi} for the extension to higher stacks and \cite{haine2024profinite} for the profinite formulation used below.

Oda's theorem gives a concrete description of the \'etale homotopy types of the moduli spaces $\mathcal M_g^{n}$.  As recalled in Example~\ref{example: Oda}, for $2g-2+n>0$ there is an equivalence in the
$\infty$-category of prospaces
\[
\Pi^{\acute et}_\infty(\mathcal M_g^n\otimes_{\mathbb Z}\overline{\mathbb Q})
\simeq \widehat{\cs\Gamma_g^n}.
\]
Thus the \'etale homotopy type, which is generally difficult to compute, is
identified in this case with the profinite completion of the classifying space of 
of a mapping class group.  

For our topological model of the Teichm\"uller tower, we replace $\Gamma_g^n$ by the boundary mapping class group $\Gamma_{g,n}$.  On the algebro-geometric side, this corresponds to replacing marked points by marked points equipped with non-zero tangent vectors.  In Section~\ref{sec: tower of etale homotopy types}, we let $\mathcal M_{g,n}^{\tan}$ denote the moduli stack of smooth curves with $n$ ordered marked points equipped with non-zero tangent vectors.  We extend Oda's comparison in Lemma~\ref{lemma:tangent-moduli-comparison}, proving that
\[
\Pi^{\acute et}_\infty(\mathcal M_{g,n}^{\tan}\otimes_{\mathbb Z}
\overline{\mathbb Q})
\simeq \widehat{\cs\Gamma_{g,n}}.
\]
Using the equivalence $\cs\mathbf S(g,n+1)\simeq \cs\Gamma_{g,n+1},$ this identifies $\widehat{\cs\mathbf S}(g,n+1)$ with the corresponding \'etale homotopy type of $\mathcal M_{g,n+1}^{\tan}$.


It is not formal, however, that these profinite spaces assemble into a modular operad.  The classifying space functor preserves products, so the spaces $\{\cs\mathbf S(g,n+1)\}$ inherit a strict modular operad structure from $\mathbf S$.  After profinite completion, this argument no longer applies: profinite completion of spaces rarely preserves products (\cite[Proposition 3.9]{Boavida-Horel-Robertson}).  This is a problem because the composition laws in a modular operad are encoded by products over the vertices of graphs (Section~\ref{sec: modular operads}).

To address this, we introduce a homotopy-coherent version of the relevant structure.  In Section~\ref{sec: homotopy action}, we work with modular $\infty$-operads, using a genus-graded version of the Segal model
developed by Hackney, the second author and Yau in \cite{hry1}.  In this model, a modular $\infty$-operad is a contravariant functor on the modular dendroidal category of genus-graded graphs satisfying Segal equivalences over the vertices of each graph.  The required compatibility with products is thus imposed up to homotopy, rather than as a strict product-preservation statement.

Applying this formalism to the profinite classifying spaces of $\bS$, we
obtain a modular $\infty$-operad in profinite spaces whose values on corollas
identify with the \'etale homotopy types of the moduli stacks
$\mathcal M_{g,n+1}^{\tan}$.  Its values on more general graphs encode the
corresponding gluing data through the modular Segal conditions.  After applying
classifying spaces and passing to derived automorphisms, the strict
$\NS$-action on $\widehat{\bS}$ induces an action in the homotopy category
of modular $\infty$-operads in profinite spaces.

\begin{thm*}[Theorem~\ref{NS action on topological operad}]
The profinite classifying spaces $\widehat{\cs\bS}(g,n+1)$ assemble into
a modular $\infty$-operad $\widehat{\cs\bS}$ in profinite spaces.
The action of $\NS$ on $\widehat{\bS}$ induces an action of $\NS$
on $\widehat{\cs\bS}$ in the homotopy category.  Equivalently, there is
a homomorphism
\[
\NS\longrightarrow \pi_0\mathbb R\operatorname{Aut}(\widehat{\cs\bS}).
\]
After restricting to genus zero, this gives an isomorphism
\[
\GT\cong \pi_0\mathbb R\operatorname{Aut}
   (\operatorname{tr}_{\leq 0}\widehat{\cs\bS}).
\]
\end{thm*}



\subsection{Acknowledgments}
The authors would like to thank Pedro Boavida de Brito, Philip Hackney, Peter Hain, Geoffroy Horel, and Jan Steinebrunner for many useful conversations and comments.  M.R. recognizes the support of Australian Research Council Future Fellowship FT210100256. L.B.B. is grateful to the Max Planck Institute for Mathematics in Bonn for its hospitality and financial support.

\end{comment}
















\tableofcontents
\section{Preliminaries}
Throughout, we will write $\bE=(\bE, \otimes, 1_{\bE})$ to denote a cocomplete, cartesian, symmetric monoidal category such that the monoidal product $\otimes$ commutes with small colimits in each variable. In this paper, $\bE$ will be (completed versions of) the category of sets, groupoids, or simplicial sets. All of our categories are \emph{simplicial} categories, meaning that they are tensored and co-tensored over the category of simplicial sets. In particular, this means that one can replace the set of maps $\Hom_{\bE}(X,Y)$ by a space\footnote{We follow a common abuse of terminology and often use the term \emph{space} when referring to simplicial sets.} of maps $\Map_{\bE}(X,Y)$, for any $X,Y$ in $\bE$. To reduce clutter, we will often suppress the category $\bE$ from our notation and simply write $\Hom(X,Y)$ or $\Map(X,Y)$ when no confusion can arise. 

\medskip 

In order to discuss homotopy classes of maps between (modular and cyclic) operads, we will make use of the language of Quillen model categories and $\infty$-categories. For Quillen model categories, we take \cite{hirsch} or \cite{balchin} as comprehensive references.  We will use the formalism of relative categories from \cite{barwick_kan} or \cite[Chapter 13]{balchin} as our model for $\infty$-categories. 


\begin{definition}\label{def: relative category}
A \emph{relative category} is a pair $(\bE, \mathscr{W})$ where $\bE$ is a small category and $\mathscr{W}$ is a wide subcategory of $\bE$ whose morphisms are a class of designated \emph{weak equivalences} in $\bE$.   A functor of relative categories $F:(\bE,\mathscr{W})\rightarrow(\bE',\mathscr{W}')$ consists of a functor $F:\bE\rightarrow\bE'$ such that $F(\mathscr{W})\subset \mathscr{W}'$.
\end{definition}

\begin{example}
Any model category $\bE$ is a relative category in which $\mathscr{W}$ is the class of weak equivalences in the model structure on $\bE$.    
\end{example}

The \emph{homotopy category} of a relative category $\bE$, denoted $\Ho(\bE)$, is obtained by formally inverting the morphisms in $\mathscr{W}$.  To describe maps in the homotopy category $\Ho(\bE)$ we use the Dwyer-Kan enhancement of the homotopy category. This is a simplicially enriched category, $\mathcal{L}(\bE,\mathscr{W})$, which has the same objects as $\bE$ and, for any two objects $X,Y\in\ob(\bE)$, a space of maps $\mathbb{R}\Map(X,Y)$.\footnote{Here, the $\mathbb{R}$ is to indicate that we are taking the ``right derived'' version of the $\Map(-,-)$ functor. This is also commonly denoted $\Map^h(X,Y)$ in the literature, e.g. \cite{hirsch}.} By considering $\bE$ to be a discrete simplicial category, one can define a natural inclusion $\bE\rightarrow\mathcal{L}(\bE, \mathscr{W})$, with the property that, for every $X,Y\in\ob(\bE)$, \[\Hom_{\Ho(\bE)}(X,Y) =\pi_0\mathbb{R}\Map_{\bE}(X,Y).\] The fact that the connected components of $\mathbb{R}\Map_{\bE}(X,Y)$ correspond to morphisms in $\bE$ reflects that the simplicial category $\mathcal{L}(\bE,\mathscr{W})$ captures the higher homotopical data of morphisms in $\bE$. In particular, for any weak equivalence $X\rightarrow X'$ and any $Y\in\ob(\bE)$ in $\bE$ the induced maps
\[\begin{tikzcd} \mathbb{R}\Map_{\bE}(X',Y)\arrow[r]& \mathbb{R}\Map_{\bE}(X,Y)\end{tikzcd} \quad \text{and} \quad \begin{tikzcd} \mathbb{R}\Map_{\bE}(Y,X)\arrow[r]& \mathbb{R}\Map_{\bE}(Y,X')\end{tikzcd}\] are weak equivalences of spaces.  

In general, it is difficult to compute the derived mapping spaces $\mathbb{R}\Map_{\bE}(X,Y)$. However, if $\bE$ admits the structure of a simplicial model category (\cite[Chapter 9]{hirsch}), we have the following useful fact (\cite[Corollary 4.7]{dk}):

\begin{thm}\label{thm: identifying mapping spaces in a simplicial model category}
Let $\bE$ be a simplicial model category. If $X$ is cofibrant and $Y$ is fibrant in $\bE$, then there exists a natural zigzag of weak equivalences \[\Map_{\bE}(X,Y)\simeq \mathbb{R}\Map_{\bE}(X,Y).\]
\end{thm}




Let $\mathscr{C}$ and $\mathscr{D}$ be two simplicial model categories. A simplicial Quillen adjunction between two simplicial model categories $
\begin{tikzcd} \mathscr{C}  \arrow[r, shift left =1, "\text{F}"]& \arrow[l, shift left =1, "\text{G}"]\mathscr{D} \end{tikzcd}$ gives rise to a homotopy adjunction \begin{equation}
\begin{tikzcd} \mathcal{L}(\mathscr{C},\mathscr{W}_{\mathscr{C}})  \arrow[r, shift left =1, "\mathbb{L}\text{F}"]& \arrow[l, shift left =1, "\mathbb{R}\text{G}"]\mathcal{L}(\mathscr{D},\mathscr{W}_{\mathscr{D}}). \end{tikzcd}
\end{equation} This means that there exists a weak equivalence of spaces \[\mathbb{R}\Map_{\mathscr{D}}(\mathbb{L}\text{F}(X),Y)\simeq \mathbb{R}\Map_{\mathscr{C}}(X,\mathbb{R}\text{G}(Y)),\] for all $X\in\mathscr{C}$ and $Y\in\mathscr{D}$. Here, the notion $\mathbb{L}$ and $\mathbb{R}$ denotes that we are using the left and right derived functors of $\text{F}$ and $\text{G}$, respectively. 

\subsection{Adjunctions between spaces and groupoids}\label{sec: classifying space functor}  Let $\sSet$ denote the category of simplicial sets with its usual Kan-Quillen simplicial model structure (e.g. \cite[Chapter 6.1]{balchin}) and let $\Grpd$ denote the category of groupoids equipped with the model structure in which weak equivalences are equivalences of categories, cofibrations are morphisms that are injective on objects, and fibrations are isofibrations (e.g. \cite{And}, \cite{Bous_model_cat_groupoid}).  Note that in this model structure, every object in $\Grpd$ is both fibrant and cofibrant.  

The functor which sends a space $X$ to its fundamental groupoid $\Pi_1(X)$, is the left adjoint in an adjunction \begin{equation}\label{classifying space adjunction}\begin{tikzcd} \Pi_{1}:\sSet  \arrow[r, shift left =1]& \arrow[l, shift left =1]\Grpd: \cs. \end{tikzcd}\end{equation} The right adjoint is called the \emph{classifying space} functor. The category of groupoids becomes a simplicial category via this adjunction. In particular, one defines the space of maps between two groupoids as \[\Map_{\Grpd}(\mathscr{C}, \mathscr{D}):=\Map_{\sSet}(\cs\mathscr{C}, \cs\mathscr{D}).\]  The classifying space functor preserves and reflects weak equivalences and fibrations which makes \eqref{classifying space adjunction} into a simplicial Quillen adjunction in which the right adjoint is \emph{homotopically fully faithful} in the sense that the induced map \[\begin{tikzcd}\mathbb{R}\Map_{\Grpd}(\mathscr{C},\mathscr{D})\arrow[r]& \mathbb{R}\Map_{\sSet}(\cs \mathscr{C},\cs \mathscr{D})\end{tikzcd}\] is a weak equivalence of spaces for all $\mathscr{C},\mathscr{D}\in\Grpd$. 

\subsection{Profinite completion of spaces and groupoids}
For any small category $\bE$, the category of pro-objects in $\bE$, denoted $\Pro(\bE)$, can be thought of as the category obtained by freely adding all co-filtered limits to $\bE$. More precisely, one can describe the \emph{category of pro-objects} in $\bE$, $\Pro(\bE)$, as the category whose objects are pairs $(I,X)$ where $I$ is a cofiltered category and $X=\{X_{i}\}_{i\in I}$ is a diagram $X:I\rightarrow \bE$. If $X:I\to\bE$ and $Y:J\to\bE$ are cofiltered diagrams, then morphisms in $\Pro(\bE)$ are defined by
\[
\Hom_{\Pro(\bE)}(\{X\}_{i}, \{Y\}_{j})
=\lim_{j\in J}\,\colim_{i\in I}\,\Hom_{\bE}(X_i,Y_j).
\]
Similarly, given a full subcategory $\mathscr{C}\subseteq \bE$, we write $\Pro_{\bE}(\mathscr{C})\subseteq \Pro(\bE)$ for the full subcategory spanned by those pro-objects $(I,X)$ such that $X_i\in \mathscr{C}$ for all $i\in I$; equivalently, $\Pro_{\bE}(\mathscr{C})\cong \Pro(\mathscr{C})$.

If the category $\bE$ is complete and $\mathscr{C}\subseteq\bE$ is a full subcategory closed under finite limits, we can define an adjunction \begin{equation}\label{profinite compeletion adjunction}
\begin{tikzcd} \widehat{(-)}:\bE\arrow[r, shift left =1] & \arrow[l, shift left =1] \Pro_{\bE}(\mathscr{C}): |-|. \end{tikzcd}
\end{equation} The left adjoint, $\widehat{(-)}$, is called the \emph{pro-$\mathscr{C}$-completion} functor. The right adjoint, $|-|$, sends a pro-object $\{X_i\}_{i\in I}$ to its limit in $\mathscr{E}$.  If, moreover, $\bE$ is a cofibrantly generated simplicial model category and $\mathscr{C}\subseteq \bE$ satisfies a specific list of conditions (e.g., the conditions of a \emph{fibration test category} in \cite[Definition 5.1]{BM20}) then $\Pro_{\bE}(\mathscr{C})$ admits a fibrantly generated model category structure and the adjunction \eqref{profinite compeletion adjunction} becomes a simplicial Quillen pair. 

\subsubsection{Profinite groupoids}
Let $\Gr$ denote the category of discrete groups and let $\Gr_{\mathrm{fin}}\subseteq \Gr$ be the full subcategory of finite groups. The pro-category
\[
\widehat{\Gr}:=\Pro_{\Gr}(\Gr_{\mathrm{fin}})
\] is the category of \emph{profinite groups}. There is an adjunction
\[
\begin{tikzcd}
\widehat{(-)}:\Gr \arrow[r, shift left=1] & \arrow[l, shift left=1] \widehat{\Gr}:|-|
\end{tikzcd}
\]
in which the right adjoint sends a pro-object $X=\{X_j\}_{j\in J}$ to its limit in $\Gr$.
The left adjoint is the classical profinite completion functor which sends a group $\mathsf G$ to the pro-object $\{\mathsf G/\mathsf N\}_{\mathsf N\triangleleft_f \mathsf G}$ indexed by the cofiltered poset of finite-index normal subgroups $\mathsf N\triangleleft_f \mathsf G$. Equivalently, one may identify $\widehat{\mathsf G}$ with the inverse limit
\[
\widehat{\mathsf G}:= \lim_{\mathsf N\triangleleft_f \mathsf G}\,\mathsf G/\mathsf N
\]
in $\Gr$.

\begin{example}\label{example: etale fundamental groups}
Let $X$ be a connected, locally Noetherian scheme and let
$\bar{x}\to X$ be a geometric point.  The \emph{$\acute etale$
fundamental group} $\pi_1^{\acute et}(X,\bar{x})$ may be defined as the
automorphism group of the fibre functor from finite $\acute etale$ covers of
$X$ to finite sets.  Equivalently, it is the inverse limit
\[
\pi_1^{\acute et}(X,\bar{x})
\cong
\lim_{(Y,y)\to (X,\bar{x})}\Aut_X(Y),
\]
where the limit ranges over connected finite $\acute etale$ Galois covers
$Y\to X$ equipped with a lift $y:\bar{x}\to Y$ of $\bar{x}$.  In
particular, $\pi_1^{\acute et}(X,\bar{x})$ is a profinite group. See \cite[\href{https://stacks.math.columbia.edu/tag/0BQ8}{0BQ8}]{stacks-project} or \cite[Expos\'e V]{grothendieck2003revetements} for more details. 
\end{example}

\begin{example}\label{example: mapping class groups}
Let $\calM_g^n$ denote the moduli stack of smooth genus $g$ curves with
$n$ marked points, defined over $\mathbb Q$.  For $2g-2+n>0$, the
geometric $\acute etale$ fundamental group of $\calM_g^n$ identifies,
after choosing a geometric base point, with the profinite completion of the
mapping class group of a genus $g$ surface with $n$ marked points fixed
pointwise:
\[
\pi_1^{\acute et}\bigl(\calM_g^n\otimes_{\mathbb Q}\overline{\mathbb Q}\bigr)
\cong \widehat{\Gamma_g^n}.
\]
See \cite{matsumoto2000arithmetic,boggi2009fundamental} for more details. 
%Later in this paper we will also work with mapping class groups of genus $g$ surfaces with $n$ boundary components, with each boundary component fixed pointwise. We denote this group by $\Gamma_{g,n}$ (see Definition~\ref{def: mapping class group}). Filling in each boundary component by a disc determines a natural surjection $\Gamma_{g,n}\twoheadrightarrow \Gamma_g^{n}$ whose kernel is generated by the Dehn twists about the boundary components. More precisely, there is a central short exact sequence \begin{equation}\label{eq:boundary_twist_extension} 1 \longrightarrow \mathbb{Z}^{n} \longrightarrow \Gamma_{g,n} \longrightarrow \Gamma_g^{n} \longrightarrow 1, \end{equation} where the map $\mathbb{Z}^{n}\to \Gamma_{g,n}$ sends the $i$th standard basis vector to the Dehn twist about the $i$th boundary component.
\end{example}

The profinite completion functor for groups admits a direct analogue for groupoids. We say that a
groupoid $\mathscr{E}$ is \emph{finite} if it has finitely many morphisms (hence finitely many objects),
and write $\Grpd_{\mathrm{fin}}\subseteq \Grpd$ for the full subcategory of finite groupoids. Since
$\Grpd_{\mathrm{fin}}$ is closed under finite limits, we define the category of \emph{profinite groupoids} by
\[
\widehat{\Grpd}:=\Pro_{\Grpd}(\Grpd_{\mathrm{fin}}).
\]
%Equivalently, the profinite completion of a groupoid $\mathscr{E}$ may be described as the cofiltered
%diagram of its finite quotients (obtained by quotienting by congruence relations so that the resulting
%groupoid has finitely many objects and morphisms). In particular, every finite groupoid is profinite,
%viewed as a constant pro-object.



The category of profinite groupoids $\widehat{\Grpd}$ admits a cocombinatorial model structure, and
the profinite completion functor fits into a Quillen adjunction
\[
\begin{tikzcd}
\widehat{(-)}:\Grpd \arrow[r, shift left=1] &
\arrow[l, shift left=1] \widehat{\Grpd}:|-|
\end{tikzcd}
\]
(see \cite[Theorem~4.12 and Proposition~4.22]{Horel_profinite_groupoids}), where $|-|$ denotes the
limit functor $\widehat{\Grpd}\to \Grpd$. 
The following proposition shows that, for groupoids with a finite set of objects, profinite completion
preserves binary products.

\begin{prop}\cite[Proposition~4.23]{Horel_profinite_groupoids}\label{prop: profinite completion preserves products of groupoids}
Let $\mathscr{E}$ and $\mathscr{D}$ be groupoids with a finite set of objects. Then the canonical map
\[
\widehat{\mathscr{E}\times \mathscr{D}}\longrightarrow
\widehat{\mathscr{E}}\times \widehat{\mathscr{D}}
\]
is an isomorphism in $\widehat{\Grpd}$.
\end{prop}


\subsubsection{Profinite spaces}\label{sec: profinite spaces}
A simplicial set $X$ is called \emph{$\pi$-finite} if $\pi_0(X)$ is finite and, for every basepoint
$x_0$, the groups $\pi_n(X,x_0)$ are finite for $n\ge 1$ and trivial for $n\gg 0$.
Let $\sSet_{\pi}\subseteq \sSet$ denote the full subcategory of $\pi$-finite spaces; it is closed under
finite limits. Using \cite[Theorem~5.2]{BM20}, we define the category of \emph{profinite spaces} to be
\[
\widehat{\sSet}:=\Pro_{\sSet}(\sSet_{\pi}).
\]
Quick \cite[Theorem~2.12]{Quick} equips $\widehat{\sSet}$ with a fibrantly generated model structure in
which cofibrations are monomorphisms and weak equivalences are detected by $\pi_0$, $\pi_1$, and
(twisted) cohomology with finite coefficients. The induced adjunction
\[
\begin{tikzcd}
\widehat{(-)}:\sSet \arrow[r, shift left=1] &
\arrow[l, shift left=1] \widehat{\sSet}:|-|
\end{tikzcd}
\]
is a Quillen adjunction \cite[Proposition~2.28]{Quick}, where $|-|$ denotes the limit functor.


The classifying space functor sends finite groupoids to $\pi$-finite simplicial sets. %Indeed, any finite groupoid is equivalent to a finite disjoint union of classifying groupoids of finite groups, so its classifying space is a finite disjoint union of spaces of the form $\cs\mathsf{G}$ with $\mathsf{G}$ finite and, in particular, it is $\pi$-finite. 
It follows that we have a simplicial Quillen adjunction
\[
\begin{tikzcd}
\widehat{\Pi}_1:\widehat{\sSet} \arrow[r, shift left=1] &
\arrow[l, shift left=1] \widehat{\Grpd} : \widehat{\cs}.
\end{tikzcd}
\]
Moreover, these functors fit into a commutative square of simplicial Quillen adjunctions
(e.g.\ \cite[Example~7.14]{BM20}):
\[
\begin{tikzcd}
\Grpd \arrow[d, shift right =1, "\widehat{(-)}"'] \arrow[r, shift right=1, "\cs"'] & \sSet \arrow[l, shift right=1, "\Pi_1"'] \arrow[d, shift left=1, "\widehat{(-)}"] \\
\widehat{\Grpd} \arrow[u, shift right =1, "|-|"'] \arrow[r, shift right=1, "\widehat{\cs}"'] & \widehat{\sSet} \arrow[l, shift right=1, "\widehat{\Pi}_1"'] \arrow[u, shift left=1, "|-|"]
\end{tikzcd}
\]


\begin{example}\label{example: Oda}
Artin--Mazur \cite{Artin_Mazur} and Friedlander \cite{Friedlander}
construct the $\acute etale$ homotopy type of a scheme as a pro-simplicial
set.  We write $\Pi^{\acute et}(X)$
for this pro-simplicial set, and write $\widehat{\Pi}^{\acute et}(X)$  for its profinite completion, viewed as a profinite space.  The same notation
will be used for stacks, using $\acute etale$ hypercovers.  In Section~\ref{sec: tower of etale homotopy types} we use the corresponding $\infty$-categorical formulation, where the $\acute etale$ homotopy type is viewed as a pro-object in the $\infty$-category of spaces.

Let $\calM_g^n$ denote the moduli stack of smooth genus $g$ curves with $n$ distinct ordered marked points, defined over $\mathbb Q$.  Oda considers the $\acute etale$ homotopy type of the geometric stack
$\calM_g^n\otimes_{\mathbb Q}\overline{\mathbb Q}$.  If $\Gamma_g^n$ denotes the mapping class group of a genus $g$ surface with $n$ marked points fixed pointwise, then for $2g-2+n>0$,  \cite[Theorem 1]{Oda} shows that there is a weak equivalence of profinite spaces
\[
\Pi^{\acute et}
  \bigl(\calM_g^n\otimes_{\mathbb Q}\overline{\mathbb Q}\bigr)
\simeq
\widehat{\cs\Gamma_g^n}.
\]
Equivalently, the profinite $\acute etale$ homotopy type is the profinite completion of the classifying space of the mapping class group $\Gamma_g^n$.
\end{example}


\subsection{Graphs} We model the operations of operads, cyclic operads, modular operads and their $\infty$-versions with decorated graphs. 

\begin{notation}
Let $\Sigma_n=\Aut(\{1,\ldots,n\})$ and $\Sigma_n^+=\Aut(\{0,1,\ldots,n\})$; then $\Sigma_n$ identifies naturally with the subgroup of $\Sigma_n^+$ consisting of permutations that fix $0$.
\end{notation}


\begin{definition}\label{def: graph}
A \emph{graph} $G$ consists of a (possibly empty) finite set of \emph{vertices} $V_G$, a finite set of \emph{half-edges} $A_G$,\footnote{Half-edges are also called ``arcs'' in the literature, hence the use of the letter $A$ for the set of half-edges.} together with gluing data encoded by a diagram of finite sets
\[
\begin{tikzcd}
\arrow[loop left]{l}{\dagger} A_G & D_G \arrow[l,hook'] \arrow[r,"t"] & V_G .
\end{tikzcd}
\] Here $\dagger:A_G\to A_G$ is a fixed-point-free involution, and $D_G\subseteq A_G$ is the subset of half-edges adjacent to vertices, with $t:D_G\to V_G$ recording the incident vertex.
\end{definition}

An \emph{edge} of $G$ is an orbit of $\dagger$, written $e=[a,a^\dagger]$, and we write $E_G=\{[a,a^\dagger]\mid a\in A_G\}$ for the set of edges. For $v\in V_G$, the \emph{neighborhood} of $v$ is the set of incident half-edges $\nb(v):=t^{-1}(v)\subseteq D_G$; its elements are said to be \emph{adjacent} to $v$. An edge $e=[a,a^\dagger]$ is \emph{internal} if $a,a^\dagger\in D_G$; we write $iE_G\subseteq E_G$ for the set of internal edges. The remaining half-edges
\[
\partial(G):=A_G\setminus D_G
\] are called \emph{legs}, and $\partial(G)$ is called the \emph{boundary} of $G$.

\begin{example}
\begin{enumerate}
    \item The graph consisting of a single \emph{edge} and no vertices is denoted $\updownarrow$. Concretely, it has $V_{\updownarrow}=\varnothing$, $D_{\updownarrow}=\varnothing$, and $A_{\updownarrow}=\{a,a^{\dagger}\}$ with $\dagger(a)=a^{\dagger}$.

    \item The $(n+1)$st \emph{corolla}, denoted $C_{n+1}$, is the graph with a single vertex $v$ and no internal edges. It has $n+1$ half-edges adjacent to $v$, and each of these is paired under $\dagger$ with a boundary half-edge. These boundary half-edges are the $n+1$ legs of the corolla.
\end{enumerate}
Many additional examples of graphs can be found in \cite{hry1}, \cite{hry_modular_nerve}, and in the survey paper \cite{hackney2022segal}.
\end{example}

A graph $H$ is a \emph{subgraph} of a graph $G$ if there is a morphism of the defining diagrams
\[
\begin{tikzcd}
\arrow[loop left]{l}{\dagger}A_{H}\arrow[d] & D_{H}\arrow[l,hook']\arrow[d] \arrow[r, "t"] & V_{H}\arrow[d]\\
\arrow[loop left]{l}{\dagger}A_{G}          & D_{G}\arrow[l,hook']          \arrow[r, "t"] & V_{G}
\end{tikzcd}
\]
such that the square
\[
\begin{tikzcd}
D_H \arrow[r]\arrow[d,"t"'] & D_G \arrow[d,"t"]\\
V_H \arrow[r] & V_G
\end{tikzcd}
\]
is a pullback and the induced map $V_H\to V_G$ is injective (cf.\ \cite[Definition~1.13]{hry1}).

\begin{example}\label{example: corolla at a vertex}
If $G$ has at least one vertex and $v\in V_G$, there is a canonical subgraph $C_v\hookrightarrow G$, called the \emph{corolla at $v$} (cf.\ \cite[Definition~1.5]{hry1}). It has a single vertex $\{v\}$, and its adjacent half-edges are exactly the elements of $\nb(v)$. Each such half-edge is paired with a boundary half-edge in $C_v$, so the legs of $C_v$ are naturally in bijection with $\nb(v)$.
\end{example}


For graphs with at least one vertex, the gluing data of $G$ can be assembled into a coequalizer diagram (cf.\ \cite[Construction~1.18]{hry1}).

\begin{lemma}\label{lem:coequalizer-graph}
Let $\mathrm{Fin}$ denote the category of finite sets, and let $\mathscr{I}$ be the category with three objects and three generating morphisms of the shape
\[
\begin{tikzcd}
\arrow[loop left] \bullet & \bullet \arrow[l, hook'] \arrow[r] & \bullet .
\end{tikzcd}
\]
A graph $G$ with at least one vertex, viewed as an object of $\mathrm{Fin}^{\mathscr I}$, is the coequalizer
\[
\begin{tikzcd}
\coprod_{e\in iE_G}\updownarrow \arrow[r, shift left=1.2] \arrow[r, shift right=1.2] & \coprod_{v\in V_G} C_v \arrow[r] & G .
\end{tikzcd}
\]
The two arrows send each internal edge to the two corresponding boundary legs of the corollas at its incident vertices.
\end{lemma}

\subsubsection{Additional structure on graphs}
Our combinatorial model admits a geometric realization $|G|$ obtained by gluing a point for each vertex and a closed interval for each edge, identifying the endpoints of each interval with the incident vertices via the incidence map $t:D_G\to V_G$. We say that $G$ is \emph{connected} (respectively, \emph{simply connected}) if $|G|$ is connected (respectively, simply connected) (see also \cite[Definition~1.7]{hry1}).

\begin{definition}\label{def: tree}
A \emph{tree} is a simply connected graph. A tree $T$ is \emph{rooted} if it is equipped with a distinguished leg $r\in \partial(T)$, called the \emph{root}.
\end{definition}

\begin{remark}
Any rooted tree admits a canonical orientation by directing every edge toward the root. Since $T$ is simply connected, each vertex then has a unique \emph{outgoing} incident edge (the one on the path to the root), and all remaining incident edges are \emph{incoming} (cf.\ \cite[Definition~3.4]{hackney2022segal}).
\end{remark}

In order to use graphs to model operations in modular and cyclic operads we will equip them with additional labeling data.

\begin{definition}\label{def: labelled graph}
A graph $G$ is said to be \emph{labelled} if:
\begin{enumerate}
\item it comes with a total ordering of its vertices, i.e.\ a bijection
      $\lambda:\{1,\ldots,k\}\xrightarrow{\cong} V_G$;
\item for each vertex $v\in V_G$, it comes with a bijection 
      \[
      \ell_v:\{0,1,\ldots,|\nb(v)|-1\}\xrightarrow{\cong}\nb(v);
      \]
\item it comes with a bijection of the boundary (legs)
      \[
      \ell_G:\{0,1,\ldots,n\}\xrightarrow{\cong} \partial(G),
      \qquad\text{where } n=|\partial(G)|-1.
      \] 
\end{enumerate}
We write $(G,\lambda,\ell)$ for this data. 
\end{definition}


A \emph{genus grading} on a graph $G$ is a function $\epsilon:V_G\to \mathbb{N}$. The \emph{genus}
of a genus-graded graph $(G,\epsilon)$ is
\[
\beta_1(G)+\sum_{v\in V_G}\epsilon(v),
\]
where $\beta_1(G)$ is the first Betti number of the geometric realization of $G$.
If $G$ is connected and $G\neq \updownarrow$, then
\[
\beta_1(G)=|iE_G|-|V_G|+1.
\]
%A genus-graded graph $(G,\epsilon)$ is \emph{stable}\mrnote{I am not sure we use stable anywhere} if $G$ is connected and, for every vertex $v$,
%\[
%2\epsilon(v)+|\nb(v)|-2>0.
%\]
We adopt the convention that the edge $\updownarrow$ has genus zero.

\begin{notation}\label{not:decorated-graphs}
Graphs in this paper are often equipped with additional structure, such as a vertex-ordering $\lambda$,
labelings $\ell$ (Definition~\ref{def: labelled graph}), and/or a genus grading $\epsilon$.
When we wish to display this data explicitly we write
\[
G=(G,\lambda,\ell,\epsilon),
\]
omitting any entries that are not present. When no confusion can arise, we suppress some or all of
these decorations and simply write $G$ for the underlying graph.
\end{notation}


\begin{definition}\label{def: iso of graphs}
An \emph{isomorphism of graphs} $\varphi:G\to G'$ is a triple of bijections
\[
\varphi_A:A_G\xrightarrow{\cong} A_{G'},\qquad
\varphi_D:D_G\xrightarrow{\cong} D_{G'},\qquad
\varphi_V:V_G\xrightarrow{\cong} V_{G'}
\]
such that the following diagram commutes:
\[
\begin{tikzcd}
A_G \arrow[d,"\varphi_A"'] \arrow[loop left,"\dagger_G"] &
D_G \arrow[l,hook'] \arrow[d,"\varphi_D"] \arrow[r,"t_G"] &
V_G \arrow[d,"\varphi_V"] \\
A_{G'} \arrow[loop left,"\dagger_{G'}"] &
D_{G'} \arrow[l,hook'] \arrow[r,"t_{G'}"'] &
V_{G'} .
\end{tikzcd}
\]
If $G$ and $G'$ are labelled graphs, a \emph{strict isomorphism} $\varphi:G\to G'$ is an isomorphism that preserves the vertex ordering and the labelings.
\end{definition}

\subsection{Operads and cyclic operads}
A symmetric \emph{sequence} in $\bE$ is a collection of objects $\calP=\{\calP(n)\}_{n\ge 1}$ in $\bE$, together with a right action of the symmetric group $\Sigma_n$ on each $\calP(n)$. In particular, for each $n\ge 1$ and each $\sigma\in\Sigma_n$, there is a morphism
\[
\sigma^*:\calP(n)\to \calP(n).
\]

\begin{definition}\label{def:operad}
An \emph{operad} in $\bE$ consists of a symmetric sequence $\calP=\{\calP(n)\}_{n\ge 1}$ together with:
\begin{itemize}
\item a distinguished element $\id\in\calP(1)$;
\item partial compositions
\[
\begin{tikzcd}
\calP(n)\times \calP(m)\arrow[r,"\circ_i"]& \calP(n+m-1),
\end{tikzcd}
\qquad 1\le i\le n,
\]
\end{itemize}
satisfying the following axioms, for all $x\in\calP(n)$, $y\in\calP(m)$, $z\in\calP(\ell)$:

\begin{enumerate}
\item For $1\le i\le n$ and $1\le j\le m$,
\[
(x\circ_i y)\circ_{i+j-1} z = x\circ_i (y\circ_j z).
\]

\item For $1\le i<k\le n$,
\[
(x\circ_i y)\circ_{k+m-1} z = (x\circ_k z)\circ_i y.
\]

\item For $x\in\calP(n)$ and $1\le i\le n$, $\id\circ_1 x = x,$ and $x\circ_i \id = x.$ 

\item  For $\sigma\in\Sigma_n$ and $\tau\in\Sigma_m$,
\[
(\sigma^*x)\circ_i (\tau^*y) =
(\sigma\circ_i \tau)^*(x\circ_{\sigma^{-1}(i)} y),
\]
where $\sigma\circ_i\tau\in\Sigma_{n+m-1}$ denotes the standard block permutation determined by
inserting $\tau$ into the $i$th input of $\sigma$.
\end{enumerate}
\end{definition}


A map $f:\calP\to\calQ$ of operads is a $\Sigma_n$-equivariant map of symmetric sequences that preserves the unit and partial compositions. We write $\Op(\bE)$ for the category of operads in $\bE$.

\begin{remark}
In the literature, operads often have operations of arity $0$. In this paper, however, we work exclusively with \emph{reduced} operads, and therefore restrict to arities $n\ge 1$.
\end{remark}

\medskip 

A cyclic operad is an operad in which the symmetric group actions on inputs extend compatibly to actions on the output together with the inputs. We write
\[
\Sigma_n^+ := \Aut(\{0,1,\ldots,n\})
\]
and let $z_n^+\in\Sigma_n^+$ denote the cyclic permutation $z_n^+(i)=i+1\pmod{n+1}$.

\begin{definition}\label{def: cyclic structure}
Let $\calP=\{\calP(n)\}_{n\ge 1}$ be an operad in $\bE$. A \emph{cyclic structure} on $\calP$ consists, for each $n\ge 1$, of a right action of
$\Sigma_n^+$ on $\calP(n)$, written
\[
\sigma^*:\calP(n)\to \calP(n),
\]
such that the following conditions hold:
\begin{enumerate}
\item the restriction along $\Sigma_n\subset \Sigma_n^+$ agrees with the given $\Sigma_n$-action on $\calP(n)$;

\item for $x\in\calP(n)$ and $y\in\calP(m)$,
\[
(z_{n+m-1}^+)^*(x\circ_i y)=
\begin{cases}
(z_n^+)^*(x)\circ_{i-1} y, & 2\le i\le n,\\[2pt] (z_m^+)^*y\circ_m (z_n^+)^*x, & i=1.
\end{cases}
\]
\end{enumerate}
A \emph{cyclic operad} is an operad equipped with a cyclic structure.
\end{definition}

A map $f:\calP\to\calQ$ of cyclic operads is a map of operads commuting with the $\Sigma_n^+$-actions. We write $\Cyc(\bE)$ for the category of cyclic operads in $\bE$. As discussed in Appendix~\ref{sec: homotopy theory of mod op and truncated mod op}, the categories $\Op(\bE)$ and $\Cyc(\bE)$ admit model category structures in which a map $f:\calP\to\calQ$ is a weak equivalence if each $f(n):\calP(n)\to\calQ(n)$ is a weak equivalence in $\bE$ for all $n\ge 1$.

\begin{remark}\label{rmk:reduce-circij}
Cyclic operads are often presented using two-slot partial compositions
\[
\circ_i^j:\calP(n)\times \calP(m)\longrightarrow \calP(n+m-1),
\]
where $i\in\{0,\ldots,n\}$ and $j\in\{0,\ldots,m\}$, rather than using the one-slot compositions $\circ_i$ of an ordinary operad together with an extended
symmetric group action. The two formalisms are equivalent.

Indeed, using the $\Sigma_m^+$-action on $\calP(m)$, one may first relabel the second operation so that its $j$th slot becomes the distinguished slot $0$.
One then composes using the operation $\circ_i^0$, and finally reindexes the output using the induced $\Sigma_{n+m-1}^+$-action. For example, if
$z_m^+\in\Sigma_m^+$ denotes the cycle $z_m^+(r)=r+1\pmod{m+1}$, then, up to the chosen convention for the final labelling, one may write
\[
x\circ_i^j y
:= \rho^*\Bigl(x\circ_i^0 \bigl(((z_m^+)^{-j})^*y\bigr)\Bigr),
\]
where $\rho\in\Sigma_{n+m-1}^+$ is the permutation determined by the chosen labelling convention for the boundary of the output. For an explicit computation
of $\rho$, see \cite{DCH}.
\end{remark}

\subsection{Modular operads}\label{sec: modular operads}
For a fixed monoidal category $\bE$, a \emph{genus-graded sequence} in $\bE$ is a collection of objects
\[
\calP=\{\calP(g,n+1)\}_{g\ge 0,\; n\ge 0}
\]
in $\bE$, such that each $\calP(g,n+1)$ carries a right action of the extended symmetric group $\Sigma_n^+ := \Aut(\{0,1,\ldots,n\})$.

\begin{definition}\label{def: modular operad}
A \emph{modular operad} in $\bE$ consists of a genus-graded sequence $\calP=\{\calP(g,n+1)\}$ together with:
\begin{enumerate}
 \item a distinguished unit object $\id\in\calP(0,2)$;
 \item a family of partial compositions
 \[
 \begin{tikzcd}
 \calP(g,n+1)\times \calP(h,m+1)\arrow[r,"\circ_{i}^{j}"]& \calP(g+h, n+m),
 \end{tikzcd}
 \]
 for $i\in\{0,1,\ldots,n\}$ and $j\in\{0,1,\ldots,m\}$;
 \item a family of contractions
 \[
 \begin{tikzcd}
 \calP(g,n+1)\arrow[r,"\xi_{i}^{j}"]& \calP(g+1,n-1),
 \end{tikzcd}
 \]
 for $0\leq i<j\leq n$ and $n\geq 2$.
\end{enumerate}
Moreover, we require the composition maps to satisfy associativity, equivariance, and unitality axioms as in \cite[Definition~3.3.1]{hv_cyclic} or \cite[Theorem~3.7]{gk_modular}. In addition, we require that the coherence relations governing contractions satisfy the following axioms: 
\begin{enumerate}
\setcounter{enumi}{3}
\item For $\sigma\in\Sigma_n^+$ and $a\in\calP(g,n+1)$,
\[
\bar\sigma^*(\xi_i^j a) = \xi_{\sigma(i)}^{\sigma(j)}(\sigma^*a),
\]
where $\xi_{\sigma(i)}^{\sigma(j)}$ means contraction along the unordered pair $\{\sigma(i),\sigma(j)\}$, and $\bar\sigma$ is the induced permutation of the remaining labels after the two contracted labels have been removed.
\item If $\{i,j\}\cap\{k,\ell\}=\varnothing$, then the two contractions commute:
\[
\xi_i^j\circ \xi_k^\ell
=
\xi_k^\ell\circ \xi_i^j,
\]
with the evident relabelling of indices after the first contraction.

\item For $a\in\calP(g,n+1)$ and $b\in\calP(h,m+1)$, with $n\geq 2$,
\[
\xi_1^2\bigl(a\circ_n^0 b\bigr) = (\xi_1^2 a)\circ_{n-2}^0 b.
\]

\item For $a\in\calP(g,n+1)$ and $b\in\calP(h,m+1)$, with $m\geq 2$,
\[
\xi_n^{n+1}\bigl(a\circ_n^0 b\bigr) = a\circ_n^0(\xi_1^2 b).
\]

\item For $a\in\calP(g,n+1)$ and $b\in\calP(h,m+1)$ with $n\ge 1$,
\[
\xi_{n-1}^{n}\bigl(a\circ_n^0 b\bigr)
=
\xi_{n+m-2}^{\,n+m-1}\Bigl(a\circ_{n-1}^0 (z_m^+)^*b\Bigr),
\]
where $z_{m}^+\in \Sigma_m^+$ is the cycle $(0\,1\,\dots\,m)$.
\end{enumerate}


\end{definition}


A map $f:\calP\to\calQ$ of modular operads is a map of the underlying genus-graded sequences that is equivariant for the $\Sigma_n^+$-actions and commutes with the unit, partial compositions, and contractions. We write $\Mod(\bE)$ for the category of modular operads in $\bE$. As we discuss in Theorem~\ref{model structure on Mod(E)}, $\Mod(\bE)$ admits a model category structure in which a map $f:\calP\to\calQ$ is a weak equivalence if each component
\[
f(g,n+1):\calP(g,n+1)\longrightarrow \calQ(g,n+1)
\]
is a weak equivalence in $\bE$ for all $g\ge 0$ and $n\ge 0$.

\begin{remark}\label{rmk:indexing-modular-operads}
Our indexing conventions for modular operads are chosen to match the mapping class group conventions of Section~\ref{subsec: mapping class groups}. Thus we write a genus-graded sequence as $\{\calP(g,n+1)\}_{g\geq 0,\;n\geq 0}$, so that the boundary labels are $\{0,1,\ldots,n\}$ and the symmetric group action is by $\Sigma_n^+=\Aut(\{0,1,\ldots,n\})$. This is a slight shift from the convention used in parts of the literature, where one often writes $\{\calP(g,n)\}_{g,n\geq 0}$.

We also exclude operations with no boundary components. Equivalently, in the above structure maps we only include compositions and contractions whose targets still have nonempty boundary; see Appendix~\ref{sec: homotopy theory of mod op and truncated mod op}.
\end{remark}


\begin{definition}\label{def:decorated graph}
Let $G=(G,\lambda,\ell,\epsilon)$ be a connected labelled, genus-graded graph with nonempty boundary, and let $\calP=\{\calP(g,n+1)\}$ be a genus-graded sequence in $\bE$. A \emph{$\calP$-decoration} of $G$ is an ordered tuple 
$\underline{p}=(p_{\lambda(1)},\ldots,p_{\lambda(k)})$, where $V_G=\{v_{\lambda(1)},\ldots,v_{\lambda(k)}\}$ and 
$p_{\lambda(i)}\in \calP\bigl(\epsilon(v_{\lambda(i)}),\,|\nb(v_{\lambda(i)})|\bigr)$ for each $i$.
Equivalently,
\[
\underline{p}\in \prod_{v_{\lambda(i)}\in V_G}\calP\bigl(\epsilon(v_{\lambda(i)}),\,|\nb(v_{\lambda(i)})|\bigr).
\]
\end{definition}

Let $G=(G,\lambda,\ell,\epsilon)$ and $G'=(G',\lambda',\ell',\epsilon')$ be labelled graphs with nonempty boundary. For
\[
i\in\{0,\ldots,|\partial(G)|-1\}
\qquad\text{and}\qquad
j\in\{0,\ldots,|\partial(G')|-1\},
\]
the \emph{grafting} $G\circ_i^j G'$ is the labelled graph obtained by gluing the boundary leg $\ell_G(i)\in\partial(G)$ to the boundary leg $\ell_{G'}(j)\in\partial(G')$, with the remaining structure inherited in the evident way. Similarly, for $0\le i<j\le |\partial(G)|-1$, the \emph{contraction} $\xi_i^j G$ is the labelled graph obtained by gluing together the boundary legs $\ell_G(i)$ and $\ell_G(j)$.

\begin{definition}\label{def: free modular operad}
Let $\calP=\{\calP(g,n+1)\}$ be a genus-graded sequence in $\Set$. The \emph{free modular operad} $\mathcal{F}\calP$ is defined by letting $\mathcal{F}\calP(g,n+1)$ be the set of isomorphism classes of $\calP$-decorated, connected, labelled, genus-graded graphs of total genus $g$ with $n+1$ boundary legs.

The group $\Sigma_n^+$ acts on $\mathcal{F}\calP(g,n+1)$ by permuting the boundary labels. The partial compositions and contractions in $\mathcal{F}\calP$ are given by grafting and contracting decorated graphs.
\end{definition}

To make this more explicit, for each $g\ge 0$ and $n\ge 0$ one may describe $\mathcal{F}\calP(g,n+1)$ as the quotient
\[
\mathcal{F}\calP(g,n+1)=
\Biggl(
\coprod_{[G]} \prod_{v\in V_G}\calP\bigl(\epsilon(v),|\nb(v)|\bigr)
\Biggr)\Big/\sim,
\]
where the coproduct ranges over isomorphism classes $[G]$ of connected labelled genus-graded graphs $G=(G,\lambda,\ell,\epsilon)$
with $|\partial(G)|=n+1$ and total genus
\[
g=\beta_1(G)+\sum_{v\in V_G}\epsilon(v).
\]

The equivalence relation $\sim$ identifies two decorated labelled genus-graded graphs $(G,\lambda,\ell,\epsilon,\underline{p})$ and $(G',\lambda',\ell',\epsilon',\underline{p}')$ whenever there exists an isomorphism of genus-graded graphs $\iota:G\xrightarrow{\cong} G'$ such that, for each vertex $v_{\lambda(i)}\in V_G$, the induced bijection $\iota:\nb(v_{\lambda(i)})\xrightarrow{\cong}\nb(\iota(v_{\lambda(i)}))$ determines a permutation
\[
\sigma_{v_{\lambda(i)}}:=
(\ell'_{\iota(v_{\lambda(i)})})^{-1}\circ \iota\circ \ell_{v_{\lambda(i)}}
\in \Sigma_{|\nb(v_{\lambda(i)})|-1}^+,
\]
and the decorations satisfy
\[
p'_{\lambda'(j)}=\bigl(\sigma_{v_{\lambda(i)}}\bigr)^*(p_{\lambda(i)})
\]
whenever $\iota(v_{\lambda(i)})=v'_{\lambda'(j)}$.  For further details and variants of this construction, see \cite[2.20]{gk_modular}, \cite[2.1]{hry_modular_nerve}, and \cite[13.2.2]{yau_modular}.





















 

\subsubsection{Truncations of modular operads and adjunctions}\label{subsec: intro to truncations}
In Appendix~\ref{subsec: truncations} we define the \emph{genus-$g$ truncation} of a modular operad
\[
\calP=\{\calP(h,n+1)\}_{h\geq 0,\;n\geq 0},
\]
denoted $\trun{g}\calP$, as the truncated modular operad obtained by retaining only the operations whose output has genus at most $g$. Thus compositions and contractions whose output has genus greater than $g$ are no longer defined. More concretely, we say 
\[
x\circ_i^j y
\qquad\text{and}\qquad
\xi_i^j w
\]
is undefined whenever
\begin{itemize}
\item $x\in\calP(h_1,m_1+1)$ and $y\in\calP(h_2,m_2+1)$ with
$h_1+h_2>g$, or
\item $w\in\calP(h,m+1)$ with $h\geq g$, so that
$\xi_i^j w\in\calP(h+1,m-1)$ has genus greater than $g$.
\end{itemize}
We write $\Mod_{\leq g}(\bE)$ for the category of genus-$g$ truncated modular
operads in $\bE$. The truncation construction defines a functor
\[
\trun{g}^*:\Mod(\bE)\longrightarrow \Mod_{\leq g}(\bE),
\]
which admits both a left adjoint and a right adjoint:
\[
\begin{tikzcd}
\Mod(\bE) \arrow[r, "\trun{g}^*"'] &
\Mod_{\leq g}(\bE)
\arrow[l, bend right=50, "(\trun{g})_!"']
\arrow[l, bend left=50, "(\trun{g})_*"]
.
\end{tikzcd}
\]

\begin{prop}
The category of cyclic operads is naturally equivalent to the category of genus-$0$ truncated modular operads. Moreover, the truncation functor
\[
\trun{0}^*:\Mod(\bE)\longrightarrow \Cyc(\bE)
\]
preserves weak equivalences and is the right adjoint in a Quillen adjunction.
\end{prop}

Similarly, there is a forgetful functor $u^*:\Cyc(\bE)\to \Op(\bE)$ which forgets the cyclic structure. This functor
admits a left adjoint, called the \emph{cyclic envelope}:
\[
\begin{tikzcd}
\Op(\bE)\arrow[r, shift left=1, "u_!"] &
\Cyc(\bE)\arrow[l, shift left=1, "u^*"]
.
\end{tikzcd}
\]
An explicit description of $u_!$ is somewhat involved; see \cite[Section~3.1]{DCH} or \cite[Section~9]{Ward}. The forgetful functor $u^*$ preserves weak equivalences, and the adjunction above is a Quillen adjunction.

\subsection{Mapping Class Groups}\label{subsec: mapping class groups}
A surface of type $(g,n+1)$ is a compact, oriented surface $S$ of genus $g$ with $n+1$ boundary components. 
% We equip each boundary component with a chosen \emph{collar}, i.e.\ a fixed neat embedding $[0,1)\times S^1\hookrightarrow S$ whose image is a neighbourhood of that boundary component.


\begin{definition}\label{def: labelled surface}
A \emph{labelled surface} of type $(g,n+1)$ is a pair $(S,\delta)$, where
\[
\begin{tikzcd}
\delta: \{0,\ldots,n\} \arrow[r, "\cong"] & \pi_0(\partial S)
\end{tikzcd}
\]
is a labelling of the boundary components, i.e. we write $\partial_i S:=\delta(i)$ for the $i$th boundary component. A \emph{parameterized labelled surface} $(S,\delta,\{x_i\})$ consists of a labelled surface together with a choice of point $x_i\in \partial_i S$ for each $0\le i\le n$. 
\end{definition}


\begin{definition}\label{def: mapping class group}
Let $S$ be a surface of type $(g,n+1)$. The associated \emph{pure mapping class group} is
    \[ \Gamma(S)\;:=\;\pi_0\,\Diff^{+}(S;\partial S), \]
where $\Diff^{+}(S;\partial S)$ denotes the group of orientation-preserving diffeomorphisms of $S$ that fix $\partial S$ pointwise. Elements of the mapping class group $\Gamma(S)$ are called \emph{mapping classes}. The group structure on $\Gamma(S)$ is induced by composition, and we use the usual order of composition of functions when writing the group multiplication.

Any two surfaces of type $(g,n+1)$ have isomorphic pure mapping class groups and we denote this group by $\Gamma_{g,n+1}$.
\end{definition}

\begin{remark}
    More generally, a \emph{mapping class} may refer to the isotopy class of a diffeomorphism between different manifolds that is not required to fix the boundary pointwise. Mapping classes that do fix the boundary pointwise are often called \emph{pure mapping classes}. Since all but one of the mapping classes considered in this paper (see Figure~\ref{fig:B-move}) fix the boundary pointwise, we will use the term \emph{mapping class} to mean \emph{pure mapping class}.
\end{remark}

\subsubsection{Presentations of mapping class groups}\label{subsubsec:presentations of mapping class groups} 

Mapping class groups of orientable surfaces are well understood, and several presentations are available \cite{GERVAIS2001703,humphries,Hatcher_Thurston,Wajnryb}. We highlight two cases that will play an important role throughout the paper. For a comprehensive reference, see \cite{farb-marg}.

\begin{example}
    The mapping class group of the disc is trivial, i.e. $\Gamma(D^2)=*$. This is known as Alexander’s lemma \cite[Lemma~2.1]{farb-marg}.
\end{example}

We now recall the mapping classes that are used as standard generators for higher genus surfaces.
Throughout this paper, we use the terms \emph{curve} and \emph{loop} to mean the isotopy class of an embedded closed curve in $S$. The geometric \emph{intersection number} $|a\cap b|$ of two curves $a$ and $b$ is the minimum number of intersection points among all representatives in their isotopy classes. Similarly, the \emph{self-intersection number} of a curve is the minimum self-intersection number among all representatives. A \emph{simple loop} is a loop of self-intersection number $0$. 

A fundamental family of mapping classes is given by \emph{Dehn twists} about simple loops. We denote the Dehn twist about a simple loop $a$ by $D_a$. 
% The mapping class group $\Gamma_{g,n+1}$ is generated by finitely many Dehn twists. In practice we work with a standard Humphries-type generating set, represented in Figure~\ref{fig: humphries generators}. 
See \cite[Section~4.4]{MR2850125} for details.



\begin{prop}\cite{Gervais96,Luo97}\label{prop: generators and relations of mcg}
    For a surface $S$ of type $(g,n+1)$ the mapping class group $\Gamma(S)$ is generated by all Dehn twists along simple closed curves together with the following relations:
    \begin{itemize}
        \item[(C)] If $|a\cap b|=0$,  then $D_aD_b=D_bD_a$.
        \item[(B)] If $|a\cap b|=1$ and $c=D_b(a)$, then $D_c = D_b D_a D_b^{-1}$.

        Relation (B) implies the \emph{braid relation} $D_a D_b D_a=D_b D_a D_b$, if $|a\cap b|=1$.
        \item[(D)] The \emph{doughnut relation}, $(D_{a}D_{b}D_{a})^4=D_\partial$ which takes place on a subsurface of $S$ of type $(1,1)$, $a$ and $b$ are the curves represented in Figure~\ref{fig:curves-for-doughnut-rel}, and $D_\partial$ denotes the twist around the boundary. 
        \item[(L)] The \emph{lantern relation}, $D_{\partial_0}D_{\partial_1}D_{\partial_2}D_{\partial_3}=D_{b_1}D_{b_2}D_{b_3}$, where the $\partial_i$ are loops in $S$ bounding a subsurface of type $(0,4)$ and the $b_j$ are as in Figure~\ref{fig:curves-for-lantern-rel}.
    \end{itemize}
\end{prop}

    % \[\Gamma(S) \cong \left<D_{a_{1}},D_{b_1},D_{a_2}, D_{b_2}\ldots,D_{b_{g-1}},D_{a_g},D_{m_1},D_{m_2},D_{c_0}, D_{c_1}\ldots,D_{c_{n}}\mid (\text{C}), (\text{B}), (\text{D}), \ \text{and} \ (\text{L})\right>\] where $D_{a_i}$, $D_{b_i}$, $D_{m_i}$, and $D_{c_i}$ denote the Dehn twists along the corresponding curves represented in Figure \ref{fig: humphries generators}. The relations are given as follows: 
% 

\begin{figure}[ht]
        \centering
        \hfill\hfill
        \begin{subfigure}[t]{0.45\textwidth}
            \centering
            \includegraphics[]{Figures/braid-rel.pdf}
            \caption{}
            \label{fig:curves-for-doughnut-rel}
            \end{subfigure}
                \hfill
            \begin{subfigure}[t]{0.45\textwidth}
                \centering
                \includegraphics[]{Figures/lantern-rel.pdf}
                \caption{}
                \label{fig:curves-for-lantern-rel}
            \end{subfigure}
            \hfill\hfill
            \caption{}
        \label{fig:curves for presentation of mcg}
\end{figure}


We will also make repeated use of the following:
\begin{lemma}\label{lemma: conjugation by mapping classes}
    Let $S,S'$ be parametrized labelled surfaces and let $f:S\to S'$ be a diffeomorphism taking the marked point on the $i$th boundary component of $S$ to the corresponding marked point of $S'$. Let $a$ be a simple loop in $S$. Then
    \[ D_{f(a)} = f\cdot D_a\cdot f^{-1} \]
    as elements of $\Gamma(S')$. Moreover, the mapping class $f\cdot D_a\cdot f^{-1}$ depends only on the isotopy class of $f$ relative to the boundary.
\end{lemma}

\begin{proof}
    The proof is standard; see \cite[Fact~3.7]{farb-marg}.
\end{proof}

\subsubsection{A diagrammatic representation of mapping classes}\label{subsubsec: diagrammatical representation of mapping classes} 
The goal of this subsection is to give a concrete combinatorial model for mapping classes in genus $0$. More precisely, we describe how isotopy classes of boundary-fixing diffeomorphisms can be encoded by embedded trees. This provides the diagrammatic framework needed later to give a presentation of our modular operad.


Let $(S,\delta,\{x_i\}_{i=0}^{n})$ and $(S',\delta',\{x'_i\}_{i=0}^{n})$ be parametrised labelled surfaces of type $(0,n+1)$, each equipped with a chosen point $x_i\in \partial_i S$ and $x'_i\in \partial_i S'$ on the $i$th boundary component. We write $\Diff^+\bigl(S,S'; \partial \bigr)$ for the space of
orientation-preserving diffeomorphisms $\phi:S\to S'$ that preserve the boundary labels and satisfy $\phi(x_i)=x'_i$ for all $0\le i\le n$. Its set of path components
\[
\pi_0\ \Diff^+\bigl(S,S'; \partial \bigr)
\]
is the set of isotopy classes of such diffeomorphisms. In genus $0$ we will encode these isotopy classes by combinatorial data coming from embedded trees.


\begin{definition}\label{def: marking of genus 0 surface}
Let $T$ be a tree with $|\partial(T)|=n+1$ equipped with a boundary labelling
\[
\begin{tikzcd}
\ell_T: \{0,1,\ldots,n\} \arrow[r, "\cong"] & \partial(T)
\end{tikzcd}
\]
Let $(S,\delta,\{x_i\})$ be a labelled surface of type $(0,n+1)$ with chosen points $x_i\in \partial_i S$ for $0\le i\le n$. A \emph{marking} of $(S,\delta,\{x_i\})$ by $T$ is an equivalence class of embeddings $m:T\hookrightarrow S$ such that
\[
m\bigl(\ell_T(i)\bigr)=x_i\in \partial_i S, \qquad \text{ for all }i\in\{0,\ldots,n\},
\] 
where $m\sim m'$ if $m'$ is obtained from $m$ by an isotopy relative to $\partial S$ or by precomposition with an element of $\Homeo(T)$

\medskip 

A \emph{marked genus-zero surface} is a pair $\bigl((S,\delta,\{x_i\}), [m]\bigr)$ consisting of a labelled genus-zero surface $(S,\delta,\{x_i\})$ and a marking $m:T\hookrightarrow S$ by a boundary-labelled tree $T$. 
\end{definition}

Any marked surface determines additional combinatorial structure on the underlying tree. Since $T$ is embedded in the oriented surface $S$, the incident half-edges at each vertex inherit a cyclic ordering from the orientation of $S$. A graph equipped with a cyclic ordering of the incident half-edges at each vertex is called a \emph{ribbon graph}, and the corresponding data is called a \emph{ribbon structure}.

In addition, if $U_T$ is a sufficiently small tubular neighbourhood of $T$ in $S$, then $\partial U_T$ is an oriented circle, and the legs of $T$ meet $\partial U_T$ in a cyclically ordered set. Via the boundary labelling $\ell_T$, this induces a cyclic order on the set of boundary labels $\{0,\ldots,n\}$.


\begin{figure}[h!t]
    \includegraphics[width=\textwidth]{Figures/Diagramatical_representation.pdf}
    \centering
    \caption{Diagrammatic representation of a diffeomorphism $S\rightarrow S'$.}
    \label{fig: diagramatic rep of diffeo}
\end{figure}


\medskip 

Let $S=(S,\delta,\{x_i\}_{i=0}^{n})$ and $S'=(S',\delta',\{x'_i\}_{i=0}^{n})$ be parametrized labelled surfaces of type $(0,n+1)$, and let $[m]$ and $[m']$ be markings on $S$ and $S'$, respectively, that induce the same cyclic order on the boundary labels. Choose representatives
\[
m:T\hookrightarrow S
\qquad\text{and}\qquad
m':T'\hookrightarrow S'
\]
of the markings $[m]$ and $[m']$. Using these representatives, we construct an orientation-preserving diffeomorphism $S\to S'$ up to isotopy relative to the boundary as follows. 


Choose small tubular neighbourhoods $U_T\subset S$ and $U_{T'}\subset S'$ of the embedded trees $m(T)$ and $m'(T')$. After smoothing corners, both $U_T$ and $U_{T'}$, as well as their complements
\[
S\setminus \operatorname{int}(U_T)
\qquad\text{and}\qquad
S'\setminus \operatorname{int}(U_{T'}),
\]
are discs. The markings and boundary data determine an orientation-preserving homeomorphism
\[
\partial U_T \longrightarrow \partial U_{T'}
\]
well-defined up to isotopy relative to the marked boundary points. Such a boundary map extends over each of the two discs uniquely up to isotopy relative to the boundary, since the mapping class group of the disc is trivial (Alexander’s lemma). We therefore obtain diffeomorphisms 
\[
\phi_{T,T'}^1:U_T\to U_{T'}
\qquad\text{and}\qquad
\phi_{T,T'}^2:S\setminus \operatorname{int}(U_T)\to S'\setminus \operatorname{int}(U_{T'}).
\]
Gluing these along $\partial U_T$ produces an orientation-preserving diffeomorphism
\[
\phi:S\to S',
\]
well-defined up to isotopy relative to $\partial S$.  The following lemma shows that the isotopy class of $\phi:S\rightarrow S'$ depends only on the isotopy classes of the markings on $S$ and $S'$, respectively.


\begin{lemma}\label{lemma: marking-independence} 
Let $S=(S,\delta,\{x_{i}\}_{i=0}^{n})$ and $S'=(S',\delta',\{x'_{i}\}_{i=0}^{n})$ be two parameterized labelled surfaces of type $(0,n+1)$, and let $[m]$ and $[m']$
be markings on $S$ and $S'$, respectively, inducing the same cyclic order on the boundary labels $\{0,\ldots,n\}$.

Let
\[
m:T\hookrightarrow S\qquad \text{and}\qquad \tilde m:\widetilde T\hookrightarrow S
\]
be two representatives of the marking $[m]$, and let
\[
m':T'\hookrightarrow S'\qquad \text{and}\qquad  \tilde m':\widetilde T'\hookrightarrow S'
\]
be two representatives of the marking $[m']$.
Then the diffeomorphisms obtained from the above construction using $(m,m')$ and $(\tilde m,\tilde m')$ determine the same isotopy class of the diffeomorphism $\phi: S\to S'.$ \qed
\end{lemma}



\subsection{The Complex of Markings}\label{sec: complex of markings}
In the study of presentations of mapping class groups, it is standard to consider a $2$-dimensional CW complex whose vertices encode decompositions of a surface and whose edges correspond to elementary moves between them. To give a presentation of our modular operad of mapping class groups, we use a variant of the complex of maximal markings of \cite[Section~5]{bk_marked_surfaces}.


A \emph{cut system} on a surface $S$ is a finite collection of pairwise disjoint simple loops $C=\{c_1,\ldots,c_k\}$, up to isotopy, such that the complement is a disjoint union of genus-zero punctured surfaces. Each puncture can be canonically replaced by a boundary component, yielding a disjoint union of genus-zero surfaces with boundary. By abuse of notation, we denote the resulting disjoint union of surfaces by
\[
\overline{S\setminus C}=\overline{S\setminus \bigcup_i c_i} \;=\; \coprod_j S_j.
\]
Each surface $S_j$ is canonically associated to a subsurface of $S$, and we refer to the surfaces $S_j$ as the \emph{cut-subsurfaces} determined by $C$. See, for example, \cite{Hatcher_Thurston,hls,moore1989classical}. 
%\lbnote{Slightly re-written: if a reader is very nitpicky, they would interpret the closure $\overline{S\setminus \bigcup_i c_i}$ in $S$. But then this closure is just the whole $S$. Original below}\mrnote{good point. }

We will enhance a cut system by equipping the cut-subsurfaces with compatible markings. These markings assemble into an embedding of a graph $G$ recording how the genus-zero pieces are connected: the vertices of $G$ correspond to cut-subsurfaces, while the internal edges correspond to the cut curves.

\begin{definition}\label{def: marked cut system}\cite[Definition~3.3]{bk_marked_surfaces}
Let $S=(S,\delta,\{x_i\}_{i=0}^n)$ be a labelled parameterized surface of type $(g,n+1)$ with $n\geq 0$. A \emph{marked cut system} on $S$ consists of a pair $(C,[m])$, where:
\begin{itemize}
    \item $C=\{c_1,\ldots,c_k\}$ is a cut system on $S$;

    \item $[m]$ is an equivalence class of embeddings $m:G\hookrightarrow S$, where $m\sim m'$ if $m'$ is obtained from $m$ by an isotopy relative to $\partial S$ or by precomposition with an element of $\Homeo(G)$.
\end{itemize}

These data are required to satisfy the following conditions:
\begin{enumerate}
    \item The embedding $m$ is compatible with the decomposition determined by
    $C$ in the following sense:
    \begin{itemize}
        \item the internal edges $iE_G$ of $G$ are in bijection with the curves in $C$, and the image of each internal edge meets the corresponding curve transversely in a single point and is disjoint from all other curves in $C$;
        
        \item the image of each vertex lies in a distinct cut-subsurface of $S$, and this determines a bijection of the vertices $V_G$ with the connected components of $\overline{S\setminus C}$;
        
        \item for each vertex $v\in V_G$, the half-edges incident to $v$ are in bijection with the boundary components of the corresponding cut-subsurface $S_v$.
    \end{itemize}

    \item The embedding $m$ sends each leg of $G$ to a marked boundary point of $S$. Equivalently, for every leg $e\in \partial(G)$ there exists $i\in\{0,\ldots,n\}$ such that the endpoint of $e$ is mapped to $x_i\in \delta(i)\subset \partial S$.

    \item For each vertex $v\in V_G$, the restriction of $m$ to the corolla subgraph $C_v\subset G$ induces a marking
    \[
    m_v:C_v\hookrightarrow S_v
    \]
    in the sense of Definition~\ref{def: marking of genus 0 surface}, where $S_v$ is the cut-subsurface corresponding to $v$ under (1).
\end{enumerate}

A marked cut system $(C,[m])$ is called a \emph{marked pants decomposition} if every cut-subsurface $S_j$ is a surface of type $(0,3)$.
\end{definition}


\begin{example}\label{ex:standard-markings}
For the pair of pants $P$, markings are represented by embeddings of the corolla $m:C_2\hookrightarrow P$. Figure~\ref{fig:standard-markings} depicts one such marking, which we call the \emph{standard marking} $m_{P}$ on $P$.

The surfaces $S_{\forkone}$ and $S_{\forktwo}$ are obtained by gluing two copies of $P$ in the two ways shown in Figure~\ref{fig:standard-markings}. The same figure shows the corresponding marked pants decompositions, which we call the \emph{standard markings} on $S_{\forkone}$ and $S_{\forktwo}$.

The surface $S_{\contraction}$ is obtained from a single copy of $P$ by gluing the boundary components $\partial_1$ and $\partial_2$. Its standard marking, also shown in Figure~\ref{fig:standard-markings}, is induced by gluing the corresponding half-edges of the embedded corolla in $P$.

\begin{figure}[h!]
    \centering
    \includegraphics[width=\linewidth]{Figures/standard-markings.pdf}
    \caption{Standard markings on $P$, $S_{\forkone}$, $S_{\forktwo}$, and
    $S_{\contraction}$.}
    \label{fig:standard-markings}
\end{figure}
\end{example}

The complex of markings is a CW complex whose vertices are marked cut systems. Its edges are generated by the following mapping classes, written using the diagrammatic representation introduced in
Section~\ref{subsubsec: diagrammatical representation of mapping classes}.

\begin{definition}\label{def:standard moves}
    Let $P$, $S_{\forkone}$, $S_{\forktwo}$, and $S_{\contraction}$ be the surfaces described in Example~\ref{ex:standard-markings}.
    \begin{itemize}
        \item The standard $T$-move is the mapping class
            \[
            \tau:P\longrightarrow P
            \]
            represented in  Figure~\ref{fig:T-move}. This is precisely the Dehn twist about the boundary component labelled $0$.
        
        \item The standard $B$-move on the pair of pants, is the mapping class 
            \[
            \beta:P\longrightarrow P
            \]
            represented in Figure~\ref{fig:B-move}. Note that this mapping class is actually an isotopy class of a diffeomorphism that does \emph{not} fix the boundary of $P$. In fact, it exchanges boundaries labelled $1$ and $2$.
            
            \begin{figure}[h!]
                \centering
                \begin{subfigure}[t]{0.45\textwidth}
                    \centering
                \includegraphics[]{Figures/sT-move.pdf}
                \caption{Standard $T$-move, $\tau$}
                \label{fig:T-move}
                \end{subfigure}
                \hfill
                \begin{subfigure}[t]{0.45\textwidth}
                    \centering
                    \includegraphics[]{Figures/B-move.pdf}
                    \caption{Standard $B$-move, $\beta$}
                    \label{fig:B-move}
                \end{subfigure}
                \caption{The standard $T$ and $B$-moves.}
                \label{fig:standard-moves-t-b}
            \end{figure}
        
        \item The standard $A$-move is the mapping class
            \[
            \alpha:S_{\forkone}\longrightarrow S_{\forktwo}
            \]
            represented in Figure~\ref{fig:A-move}.    
            \begin{figure}[h]
                \centering
                \includegraphics[]{Figures/A-move.pdf}
                \caption{Standard $A$-move, $\alpha$}
                \label{fig:A-move}
            \end{figure}
        
        \item The standard $S$-move is the mapping class
            \[
            \sigma:S_{\contraction}\longrightarrow S_{\contraction}
            \] 
            given by
            \[
            \sigma=D_aD_bD_a\in \Gamma(S_{\contraction}),
            \] where $a$ and $b$ are the curves represented in Figure~\ref{fig:S-move}. It is simple to check that $\sigma$ sends the marked cut system $(\{a\},[m_{\contraction}])$ to $(\{b\},[\sigma\circ m_{\contraction}])$, as in Figure~\ref{fig:S-move}. 
            \begin{figure}[h]
                \centering
                \includegraphics[]{Figures/S-move.pdf}
                \caption{Standard $S$-move, $\sigma$}
                \label{fig:S-move}
            \end{figure}
    \end{itemize}
\end{definition}


We now define the CW complex of markings used throughout the paper. It is a variant of the \emph{complex of maximal markings} of \cite[Section~5]{bk_marked_surfaces}. The main difference is that we restrict to markings over pants decompositions, so the cut-subsurfaces are all pairs of pants. We refer to Appendix~\ref{app: complex of markings} for the full list of $2$-cell relations, and record below those most relevant for this paper.

\begin{definition}\label{def: complex of markings}

Let $S=(S,\delta,\{x_i\}_{i=0}^{n})$ be a parametrized, labelled surface of type $(g,n+1)$, with $n\ge 0$. The \emph{complex of markings over pants decompositions on $S$}, denoted $\markS$, is the $2$-dimensional CW complex defined as follows.

\smallskip
\noindent\textbf{Vertices.}
The vertices of $\markS$ are marked pants decompositions $(C,[m])$ on $S$.

\smallskip
\noindent\textbf{$1$-cells.}
The $1$-cells are generated by the following moves.

\begin{enumerate}
\item \textbf{$T$-moves:}
Given marked pants decompositions $(C,[m])$ and $(C',[m'])$, we define a $1$-cell
\[
(C,[m])\overset{T_a}{\rightsquigarrow}(C',[m'])
\]
whenever $C=C'$ and the markings $m$ and $m'$ differ only on a cut-subsurface $S'\subset \overline{S\setminus C}$ having $a$ as one of its boundary curves, in the following way: under an identification $\phi:S'\xrightarrow{\cong} P$ satisfying $\phi\circ m|_{S'}=m_P$ and $\phi(a)=\partial_0$,
we have
    \[\phi\circ m'|_{S'}=\tau\circ m_P.\]

Thus a $T$-move keeps the underlying pants decomposition fixed and changes the
marking by a Dehn twist along one cut curve.



\item \textbf{$B$-moves:}
Given marked pants decompositions $(C,[m])$ and $(C',[m'])$, we define a $1$-cell
\[
(C,[m])\overset{B_{a,b,c}}{\rightsquigarrow}(C',[m'])
\]
whenever $C=C'$ and the markings $m$ and $m'$ differ only on a cut-subsurface $S'\subset \overline{S\setminus C}$ with boundary curves $a,b,c$, in the following way: under an identification
$\phi:S'\xrightarrow{\cong}P$ satisfying
\[
\phi\circ m|_{S'}=m_P,\qquad
\phi(a)=\partial_1,\qquad
\phi(b)=\partial_2,\qquad \text{and} \qquad
\phi(c)=\partial_0,
\]
we have
    \[ \phi\circ m'|_{S'}=\beta\circ m_P. \]
When the curves $a$ and $b$ uniquely determine the pair of pants $S'$, we omit $c$ from the notation and write $B_{a,b}$.

        
\item \textbf{$A$-moves:} 
Given marked pants decompositions $(C,[m])$ and $(C',[m'])$, we define a $1$-cell
\[
(C,[m])\overset{A_{a,b}}{\rightsquigarrow}(C',[m'])
\] 
whenever $C\setminus\{a\}=C'\setminus\{b\}$ and the markings $m$ and $m'$ differ only on a cut-subsurface 
\[
S'\subset
\overline{S\setminus (C\setminus\{a\})}
=
\overline{S\setminus (C'\setminus\{b\})}
\]
of type $(0,4)$, in the following way: under an identification $\phi:S'\xrightarrow{\cong} S_{\forkone}$ taking the local marked cut system $(\{a\},m|_{S'})$ to the standard marking on $S_{\forkone}$, the composite
\[
\alpha\circ\phi:S'\xrightarrow{\cong} S_{\forktwo}
\]
takes the local marked cut system $(\{b\},m'|_{S'})$ to the standard marking on $S_{\forktwo}$.

        
\item \textbf{$S$-moves:} 
Given marked pants decompositions $(C,[m])$ and $(C',[m'])$, we define a $1$-cell
\[
(C,[m])\overset{S_{c',c}}{\rightsquigarrow}(C',[m'])
\] 
whenever $C\setminus\{c\}=C'\setminus\{c'\}$ and the markings $m$ and $m'$ differ only on a cut-subsurface 
\[
S'\subset
\overline{S\setminus (C\setminus\{c\})}
=
\overline{S\setminus (C'\setminus\{c'\})}
\] 
of type $(1,1)$, in the following way: under an identification $S'\xrightarrow{\cong} S_{\contraction}$ taking the local marked cut system $(\{c\},m|_{S'})$ to the standard marking on $S_{\contraction}$, we have
\[
c'=\sigma(c)
\qquad\text{and}\qquad
m'|_{S'}=\sigma\circ m|_{S'}.
\]

\end{enumerate}
\smallskip
\smallskip
\smallskip
\noindent\textbf{$2$-cells.}
The $2$-cells are attached along loops expressing relations among the $1$-cells.
Most of these coincide with the relations in \cite{bk_marked_surfaces}; see also
Appendix~\ref{app: complex of markings}.

In describing these relations, we use the following conventions. For
$f\in\Gamma(S)$ and a marking $m:G\hookrightarrow S$, we write $f\circ m$ for
the marking represented by $G\xrightarrow{m}S\xrightarrow{f}S$. For a subsurface
$S'\subset S$, we write $m|_{S'}$ for the induced marking obtained by restricting
the image of $m$ to the corresponding subgraph. In particular, if $S'$ is a
cut-subsurface, then $m|_{S'}$ is the
corolla marking $m_v:C_v\hookrightarrow S'$ associated to the corresponding
vertex $v\in V_G$.

We include the following families of relations, which are described in detail in Appendix~\ref{app: complex of markings}:
\begin{enumerate}
\item commutativity of disjoint moves;
\item boundary twist relations;
\item pentagon relation;
\item two hexagon relations;
\item self-duality of associativity,
% \[
% A_{c,c'}A_{c',c}=\id;
% \]
\item  twist--braiding relation;
\item  gluing relations for Dehn twists.
\end{enumerate}

\smallskip
\noindent
In addition to the local relations above, we impose the following genus-one relations. For surfaces of type $(1,1)$, we impose
\begin{align}
    % S_{a,b} S_{b,a} B_{\partial_0,a,a} &= \id, \tag{G1.1}\label{eq:g1.1}\\
    SSB &= \id, \tag{G1.1}\label{eq:g1.1}\\
    T^{-1}ST^{-1}ST^{-1} &= S\tag{G1.2}\label{eq:g1.2}
\end{align}
shown in more detail in Figures~\ref{fig:G2-cell-G1-1}
and~\ref{fig:G2-cell-G1-2}.

\begin{figure}[h!]
    \centering
    \begin{subfigure}[t]{0.415\linewidth}
        \centering
        \includegraphics[width=\linewidth]{Figures/2-cell-g1-1.pdf}
        \caption{$2$-cell (G1.1)}
        \label{fig:G2-cell-G1-1}
    \end{subfigure}
    \hfill
    \begin{subfigure}[t]{0.545\linewidth}
        \centering
        \includegraphics[width=\linewidth]{Figures/2-cell-g1-2.pdf}
        \caption{$2$-cell (G1.2)}
        \label{fig:G2-cell-G1-2}
    \end{subfigure}
    \caption{The genus-one $2$-cells (G1.1) and (G1.2).}
    \label{fig:G2-cells-G1}
\end{figure}
For surfaces of type $(1,2)$, we also impose the relation shown in
Figure~\ref{fig:G2-cell-NS-complex-of-markings}. This relation differs slightly
from the corresponding relation in \cite{bk_marked_surfaces}.

\begin{figure}[h!]
    \centering
    \includegraphics[width=\linewidth]{Figures/2-cell-NS-complex-of-markings.pdf}
    \caption{$2$-cell (G2)}
    \label{fig:G2-cell-NS-complex-of-markings}
\end{figure}
\end{definition}



The relations above are chosen so that the resulting complex has the expected connectivity properties. The result we will use later is the following.

\begin{prop}\label{prop:our-complex-of-markings-is-simply-connected}
Let $S=(S,\delta,\{x_i\}_{i=0}^{n})$ be a parameterized, labelled surface of type $(g,n+1)$, with $n\geq 0$. Then the complex of markings over pants decompositions on $S$, $\markS,$ is connected and simply connected.
\end{prop}

A proof is given in Appendix~\ref{app: complex of markings}.


\section{A Modular Operad of Mapping Class Groupoids}\label{sec: modular operad in groupoids}
In this section, we define a modular operad in groupoids $\bS=\{\bS(g,n+1)\}_{g\ge 0,\, n\ge 0},$ which homotopically approximates a modular operad constructed from the mapping class groups in the sense that, \[\cs\bS(g,n+1)\simeq \cs\Gamma_{g,n+1}.\] In Section~\ref{sec: presentation of modular operad} we will compare each groupoid $\bS(g,n+1)$ with the corresponding complex of markings. This comparison will allow us to give a finite presentation of $\bS$ by generators and relations.


\subsection{The Modular Operad $\bS$}\label{sec: modular operad of mapping class groups}
Recall that we write $\Sigma_n^+ = \Aut({0,1,\ldots,n})$ for the extended symmetric group. If $C_{n+1}$ is a corolla with (n+1) boundary legs, then a labelling of its boundary by the set ${0,1,\ldots,n}$ is equivalent to a choice of element of $\Sigma_n^+$. Indeed, once a reference labelling is fixed, every other labelling is obtained by precomposition with a unique permutation in $\Sigma_n^+$.

In the case $n=2$, the alternating subgroup of $\Sigma_2^+ = \Aut(\{0,1,2\})$ is given by $
A_2^+ = \{\id,(012),(021)\}$. The quotient $\Sigma_2^+/A_2^+$ has two elements, represented by
\[
\id\cdot A_2^+
\qquad\text{and}\qquad
(12)\cdot A_2^+.
\]
%Geometrically, these elements correspond to the two planar orientations of the trivalent corolla with legs labelled by $\{0,1,2\}$. 
The permutations in $A_2^+$ preserve the cyclic ordering of the labels, so the three labellings in each coset determine the same cyclic ordering of the boundary legs. 

\begin{definition}\label{def:free modular operad generated by mu}
Let $\Upsilon=\{\Upsilon(g,n+1)\}$ be the free modular operad in $\Set$ generated by the set \[
\Upsilon(0,3)=\Sigma_2^+/A_2^+.
\] 
The set $\Upsilon(0,3)$ consists of the two planar trivalent corollas with boundary labelled by $\{0,1,2\}$, which we denote by
\[
P=\id\cdot A_2^+
\qquad\text{and}\qquad
(12)^*P=(12)\cdot A_2^+.  
\] Figure~\ref{fig:objects-of-PP2} depicts elements of the set $\Upsilon(0,3)$.
\end{definition}


 \begin{figure}[ht]
        \centering
        \includegraphics[width=0.4\textwidth]{Figures/generators.pdf}
        \caption{Generators of the modular operad $\Upsilon$.}
        \label{fig:objects-of-PP2}
    \end{figure}
    
\medskip 
Following Definition~\ref{def: free modular operad}, we see that an element of $\Upsilon(g,n+1)$ is an isomorphism class of connected decorated graphs $G=(G,\lambda, \ell_G,\epsilon,\underline{\upsilon})$ where:
\begin{enumerate}
    \item every vertex $v\in V_G$ is trivalent and has genus $\epsilon(v)=0$;
    \item The vertex decoration $\underline{\upsilon} = (\upsilon_{\lambda(1)}, \ldots, \upsilon_{\lambda(k)})$ consists of a sequence of elements of $\Upsilon(0,3)$. In other words, each vertex $v_{\lambda(i)}$ is decorated by one of the two planar trivalent corollas in $\Upsilon(0,3)$, equivalently by a bijection
    \[
    \ell_{v_{\lambda(i)}}:\{0,1,2\}\xrightarrow{\cong}\nb(v_{\lambda(i)})
    \]
    defined up to the action of $A_2^+$;
    \item the boundary of $G$ is equipped with a labelling
    \[
    \ell_G:\{0,\ldots,n\}\xrightarrow{\cong}\partial(G);
    \]
    \item the total genus of $G$ is $g$, equivalently $\beta_1(G)=g$.
\end{enumerate}

Condition~(2) is equivalent to equipping each vertex of $G$ with a cyclic ordering of its incident half-edges. In other words, $G$ is a ribbon graph, and the isomorphism classes above are isomorphism classes of ribbon graphs. Accordingly, we will often regard elements of $\Upsilon(g,n+1)$ as ribbon graphs with first Betti number $g$ and $n+1$ ordered legs. However, the interpretation of $\Upsilon(g,n+1)$ as the free modular operad generated by a symmetric sequence will be essential in later sections.


\medskip



Our next goal is to define, for each $g\geq 0$ and $n\geq 1$, a groupoid $\bS(g,n+1)$ whose objects are the decorated graphs in $\Upsilon(g,n+1)$. The construction is arranged so that, for every $G\in \Upsilon(g,n+1)$, there is a natural isomorphism
\[
\Aut_{\bS(g,n+1)}(G)\cong \Gamma(S_G),
\] where $S_G$ is the surface associated to $G$.  

We now describe this surface.  Let $G=(G,\lambda,\ell_G,\epsilon,\underline{\upsilon})\in \Upsilon(g,n+1)$
be a decorated graph. From $G$ we construct a labelled, parametrized surface
\[
S_G=(S_G,\delta,\{x_i\}_{i=0}^n)
\]
of type $(g,n+1)$, together with a marking over a pants decomposition in the
sense of Definition~\ref{def: marked cut system}, as follows.
\begin{enumerate}
    \item For each vertex $v_{\lambda(i)}\in V_G$, let $P_{v_{\lambda(i)}}=(P_{v_{\lambda(i)}},\delta_{v_{\lambda(i)}},\{x_0,x_1,x_2\})$ denote a copy of the standard pair of pants corresponding to the decoration $\upsilon_{\lambda(i)}\in \Upsilon(0,3).$ Thus the boundary components of $P_{v_{\lambda(i)}}$ are labelled by $\{0,1,2\}$ according to the labelling of the legs of the corolla decorated by $\upsilon_{\lambda(i)}$, i.e.  \[
    \begin{tikzcd}
    \delta := \ell_{P_{v_{\lambda(i)}}}: \{0,1,2\} \arrow[r] & \partial(P_{v_{\lambda(i)}}).
    \end{tikzcd}
    \] Moreover, $P_{v_{\lambda(i)}}$ comes equipped with the standard marking (Figure~\ref{fig:standard-markings})
\[
m_{v_{\lambda(i)}}:C_{v_{\lambda(i)}}\hookrightarrow P_{v_{\lambda(i)}},
\]
where the image of the $j$-th leg determines the marked point $x_j$ on the $j$th boundary component of $P_{v_{\lambda(i)}}$. 

    
\item If two vertices of $G$ are joined by an internal edge, we glue the corresponding boundary components of the associated pairs of pants. In this way one obtains a surface of type $(g,n+1)$, which we denote by $S_G$. 

    More precisely, if the graph $G$ is presented as the coequalizer
    \[
    \begin{tikzcd}
        \coprod\limits_{e\in iE_G}\updownarrow 
        \arrow[r, shift left=1.2] 
        \arrow[r, shift right=1.2]
        & 
        \coprod\limits_{v_{\lambda(i)}\in V_G} C_{v_{\lambda(i)}} 
        \arrow[r] 
        & 
        G,
    \end{tikzcd}
    \]
    as in Lemma~\ref{lem:coequalizer-graph}, then the associated surface $S_G$ is the coequalizer
    \[
    \begin{tikzcd}
        \coprod\limits_{e\in iE_G} S^1 
        \arrow[r, shift left=1.2] 
        \arrow[r, shift right=1.2]
        & 
        \coprod\limits_{v_{\lambda(i)}\in V_G} P_{v_{\lambda(i)}} 
        \arrow[r] 
        & 
        S_G.
    \end{tikzcd}
    \]
We note that to ensure that $S_G$ is a smooth surface one should perform the gluing along the collar neighbourhoods of the boundary components. We omit these technical details for brevity. 

\item The boundary components of $S_G$ are labelled via the boundary labelling of $G$, that is,
    \[
    \begin{tikzcd}
    \delta := \ell_G : \{0,1,\ldots,n\} \arrow[r] & \partial(S_G).
    \end{tikzcd}
    \]
    For each $i\in\{0,\ldots,n\}$, let $x_i\in \delta(i)$ denote the point where the marking $m_G$ meets the boundary component $\delta(i)$.

\item     By construction, the image of $\coprod_{e\in iE_G} S^1$ in $S_G$ defines a pants decomposition.  The local markings $m_{v_{\lambda(i)}}:C_{v_{\lambda(i)}}\hookrightarrow P_{v_{\lambda(i)}}$ glue together to give a natural embedding $m_G:G\hookrightarrow S_G.$ Together this defines a marked pants decomposition $(iE_G, [m_G])$ of the surface $S_G$. 

\end{enumerate}



\begin{definition}\label{def: surface modular operad}
We define a family of groupoids $\bS(g,n+1)$, for $2g+n\geq 1$, as follows. 

For $(g,n+1)=(0,2)$, let $\bS(0,2)$ be the groupoid with a single object $|$ and $\Hom_{\bS(0,2)}(|,|)=\mathbb Z.$ We regard this as the groupoid version of the mapping class group of a thin cylinder, whose morphisms correspond to Dehn twists along the boundary; e.g.~\cite[3.1.1]{Wahl_infinite_loop_space}. 


\medskip

For $g\ge 0$ and $n\ge 1$, the objects of $\bS(g,n+1)$ are the decorated graphs in $\Upsilon(g,n+1)$:
\[
\ob(\bS(g,n+1))=\Upsilon(g,n+1).
\]
Given any two graphs $G,G'\in \Upsilon(g,n+1)$, the morphisms are isotopy classes of diffeomorphisms
\[
\phi:S_G\to S_{G'}
\]
preserving the boundary labelling and sending the marked points on the boundary of $S_G$ to the corresponding marked points on the boundary of $S_{G'}$. Thus 
\[\Hom_{\bS(g,n+1)}(G,G'):=\pi_0\Diff(S_G,S_{G'};\partial).\] 
Composition in $\bS(g,n+1)$ is defined to be the composition of diffeomorphisms.

\end{definition} 

For $(g,n)=(0,2)$, the symmetric group $\Sigma_1^+$ acts trivially on $\bS(0,2)$. For all other pairs $(g,n)$, the extended symmetric group $\Sigma_n^+$ acts on $\bS(g,n+1)$ by permuting the labels on the boundary legs of the underlying decorated graphs. Thus $\bS=\{\bS(g,n+1)\}_{g\geq 0,\; n\geq 0}$ is a genus-graded sequence in groupoids. 

We now define the modular operad structure. For $0\leq i\leq n$ and
$0\leq j\leq m$, the composition functor
\[
\circ_i^j:\bS(g,n+1)\times \bS(h,m+1)\longrightarrow \bS(g+h,n+m)
\]
is defined on objects by grafting decorated graphs, so that
$(G,H)$ is sent to $G\circ_i^j H$, as in
Definition~\ref{def: free modular operad}. 

For $(g,n)=(0,2)$, there is a single object $|$ in $\bS(0,2)$ which acts as the identity for graph grafting. The generating morphism $1\in \mathbb Z=\Hom_{\bS(0,2)}(|,|)$ acts by Dehn twists: composing it at the $i$th input with a morphism $[\phi]:S_G\to S_H$ produces the mapping class $[D_{\partial_i}\circ \phi]:S_G\to S_H,$
where $\partial_i$ denotes the $i$th boundary component of $S_H$. 


For all other pairs $(g,n)$, the associated surface $S_{G\circ_i^j H}$ is obtained by gluing $S_G$ and $S_H$ along the boundary components labelled $\ell_G(i)$ and $\ell_H(j)$. The same gluing operation defines the functor on morphisms. Namely, if $\phi:S_G\to S_{G'}$ and $\psi:S_H\to S_{H'}$ are boundary-label-preserving diffeomorphisms, then they carry the chosen boundary components to the corresponding boundary components of $S_{G'}$ and $S_{H'}$. Hence they glue to a diffeomorphism $\phi\circ_i^j\psi:
S_{G\circ_i^j H}
\longrightarrow
S_{G'\circ_i^j H'}.$
This construction is compatible with isotopy classes, and therefore defines the desired functor.

% For $(g,n)=(0,2)$, there is a single object $|$ in $\bS(0,2)$ which behaves as the identity in graph grafting. The generating morphism $1\in\Z\in S(0,2)$ when composed with with a morphism $[\phi]:S_G\to S_H$ at an input $i$, gives as a result the mapping class $[D_{\partial_i}\cdot \phi]:S_G\to S_H$, where $\partial_i$ denotes the $i$th boundary of $S_H$. For all other pairs $(g,n)$, we see that the associated surface
% $S_{G\circ_i^j H}$ is obtained by gluing $S_G$ and $S_H$ along the boundary
% components labelled $\ell_G(i)$ and $\ell_H(j)$.

% The same gluing operation defines the functor on morphisms. Namely, if
% $\phi:S_G\to S_{G'}$ and $\psi:S_H\to S_{H'}$ are boundary-label-preserving
% diffeomorphisms, then they carry the chosen boundary components to the
% corresponding boundary components of $S_{G'}$ and $S_{H'}$. Hence they glue to a
% diffeomorphism
% $\phi\circ_i^j\psi:S_{G\circ_i^j H}\to S_{G'\circ_i^j H'}$.
% This construction is compatible with isotopy classes, and therefore defines the
% desired functor. 

Similarly, for distinct labels $i$ and $j$, the contraction functor
\[
\xi_i^j:\bS(g,n+1)\longrightarrow \bS(g+1,n-1)
\]
is defined on objects by grafting together the two boundary legs of $G$ labelled
$\ell_G(i)$ and $\ell_G(j)$. On surfaces, this is the operation of gluing the
two corresponding boundary components of $S_G$, thereby increasing the genus by
one and reducing the number of boundary components by two.

On morphisms, a boundary-label-preserving diffeomorphism
$\phi:S_G\to S_{G'}$ glues along the corresponding boundary components to give
a diffeomorphism $S_{\xi_i^jG}\to S_{\xi_i^jG'}$. This is again compatible with
isotopy classes, and hence defines the contraction functor.

\begin{definition}\label{def: modular operad of mapping class groups}
The \emph{modular operad of surfaces}, 
\[
\bS=\{\bS(g,n+1)\}_{2g+n\geq 1},
\]
is the modular operad in groupoids whose $(g,n+1)$-component is the groupoid $\bS(g,n+1)$ of genus $g$ surfaces with $n+1$ marked boundary components.  The modular operad structure is induced by gluing boundary components, or equivalently by the grafting and contraction operations on graphs.

\medskip

The \emph{cyclic operad of genus-zero surfaces}, denoted $\trun{0}\bS$, is
obtained by restricting $\bS$ to genus zero. Its component with $n+1$ boundary
components is
\[
\trun{0}\bS(n+1)=\bS(0,n+1).
\]
The cyclic operad structure is induced by the gluing operations in $\bS$ which glue distinct genus-zero surfaces along boundary components. On markings, these operations correspond to grafting the underlying labelled trees.
\end{definition} 

By construction, for every object $G\in \ob(\bS(g,n+1))$, the automorphism group of
$G$ is the mapping class group of the associated surface:
\[
\Aut_{\bS(g,n+1)}(G)=\Gamma(S_G)\cong \Gamma_{g,n+1}.
\]
Since $\bS(g,n+1)$ is a connected groupoid, the following lemma follows
immediately.

\begin{lemma}
For each $g\geq 0$ and $n\geq 0$ satisfying $2g+n\geq 1$, there is a weak equivalence
\[
\cs\bS(g,n+1)\simeq \cs\Gamma_{g,n+1}.
\]
\end{lemma}

%\begin{proof}
%By construction, $\bS(g,n+1)$ is a connected groupoid. Moreover, for any object $G\in \bS(g,n+1)$, its automorphism group is naturally isomorphic to the mapping class group:
%\[
%\Aut_{\bS(g,n+1)}(G)\cong \Gamma_{g,n+1}.
%\]
%Since a connected groupoid is equivalent to the one-object groupoid associated to the automorphism group of any of its objects, it follows that $\bS(g,n+1)$ is equivalent to the one-object groupoid with endomorphism group $\Gamma_{g,n+1}$. Therefore
%\[
%\cs\bS(g,n+1)\simeq \cs\Gamma_{g,n+1}.
%\]
%\end{proof}



\begin{remark}
The modular operad in Definition~\ref{def: modular operad of mapping class groups} is closely related to the surface operads of \cite[Construction 2.2]{Till_surface_operad} and \cite[Section 3.1]{Wahl_infinite_loop_space}. These constructions all build surfaces by gluing basic pieces along boundary components, but they differ in which pieces are taken as ``atomic''. In the constructions of Tillmann and Wahl, the atomic surfaces include the pair of pants, the one-holed torus, and the disc. In our setting, the pair of pants is the only atomic surface.
\end{remark}

\subsubsection{A presentation of the modular operad $\bS$}\label{sec: presentation of modular operad}
Fix an object $G\in\ob(\bS(g,n+1))$, regarded as a decorated graph, and let $S_G$ denote the associated surface equipped with its standard marked pants
decomposition
\[
(iE_G,[m_G]:G\hookrightarrow S_G).
\]
We compare the topological encoding of markings of $S_G$, given by the complex of markings, with the categorical encoding given by the coslice category $G\downarrow \bS(g,n+1)$.


Recall the complex of markings $\calM(S_G)$ from Definition~\ref{def: complex of markings}. We write $\mathbb{M}_G$ for the groupoid presented by this $2$-complex. Its objects are the $0$-cells of $\calM(S_G)$, its generating morphisms are the edges of
$\calM(S_G)$, and its relations are the boundaries of the $2$-cells. Equivalently, $\Hom_{\mathbb{M}_G}(v,w)$ consists of homotopy classes, relative endpoints, of edge paths from $v$ to $w$.


\begin{definition}\label{defn: coslice S under G}
Fix a graph $G\in \ob(\bS(g,n+1))$. The \emph{coslice category} $G\downarrow \bS(g,n+1)$ is the category whose objects are morphisms $\phi:G\to H$ in $\bS(g,n+1)$, as $H$ ranges over the objects of $\bS(g,n+1)$.

A morphism from $\phi_1:G\to H_1$ to $\phi_2:G\to H_2$ is a morphism $\psi:H_1\to H_2$ in $\bS(g,n+1)$ such that $\psi\phi_1=\phi_2$. In other words, a morphism from $\phi_1\rightarrow \phi_2$ is a morphism $\psi:H_1\to H_2$ which makes the diagram
\[
\begin{tikzcd}
& G \arrow[dl, "\phi_1"'] \arrow[dr, "\phi_2"] & \\
H_1 \arrow[rr, "\psi"'] && H_2
\end{tikzcd}
\]
commute.
\end{definition}

\begin{lemma}\label{lemma: the complexes of morphisms and markings}
Fix a graph $G\in \ob(\bS(g,n+1))$ and let $S_G$ be the associated labelled, parametrized surface. Then the coslice category $G\downarrow \bS(g,n+1)$ is isomorphic to the groupoid $\mathbb{M}_G$.
\end{lemma}

%This lemma will be crucial in our future proofs. In particular, the observation that morphisms that compose to $\id_G$ correspond to loops in $\calM(S_G)$.

\begin{proof}
We first construct a functor $I:\mathbb{M}_G\longrightarrow G\downarrow \bS(g,n+1)$ and then show that it is an isomorphism.

To define $I$ on objects,  let $(C,[m]:H\hookrightarrow S_G)$ be a vertex of the complex of markings $\calM(S_G)$. Since $C$ is a pants decomposition, the components of $\overline{S_G\setminus C}$ are pairs of pants. Thus the graph $H$ is trivalent, with internal edges corresponding to the curves of $C$ and vertices corresponding to the pair-of-pants
components. Since every pants decomposition of a surface of type $(g,n+1)$ has $2g+n-1 = |\chi(S_G)|$ pairs of pants, the graph $H$ has $2g+n-1$ trivalent vertices. An Euler characteristic argument then shows that $\beta_1(H)=g$.

The marking $[m]$ determines the additional structure required to regard $H$ as an object of $\bS(g,n+1)$. Indeed, the orientation of $S_G$ induces a cyclic ordering on the half-edges incident to each vertex, while the labelling of the boundary components of $S_G$ induces a labelling of the boundary legs of $H$. Thus the marked pants decomposition determines an object of $\bS(g,n+1)$, which we again denote by $H$.

%We now convert the marked cut system $(C,m)$ into an object of $\ob(\bS(g,n+1))$.  For each trivalent vertex $v\in V_{H}$ the corresponding subsurface $S'_v\subset S_G$ is a pair of pants. The embedding $m: H\hookrightarrow S_G$ determines a unique labelling permutation $\ell_v\in\Sigma_2^+$ (relative to the chosen ordering of the boundary components of the standard pair of pants), and we decorate $v$ by the corresponding element \[\mu_v = \mu(\ell_v(0),\ell_v(1),\ell_v(2)).\]

The object $H$ comes with its standard surface $S_H$ and standard marked pants decomposition $(C_H,[m_H])$. By construction, the marked pants decomposition $(C,[m])$ of $S_G$ and the standard marked pants decomposition of $S_H$ have the same associated graph. Therefore the local identifications of the
corresponding pairs of pants glue to a boundary-label-preserving diffeomorphism $\phi:S_G\longrightarrow S_H,$ which is well defined up to isotopy. We set
\[
I(C,[m]) = [\phi:G\to H],
\]
viewed as an object of the coslice category $G\downarrow \bS(g,n+1)$.

%Thus we obtain a planar, labelled, decorated graph \[H = (H, \lambda_{m}, \ell_{m}, {\mu}) \in \ob(\bS(g,n+1)).\] Recall that to such an object $H$ we associate a surface $S_H$ equipped with its standard marking over a pants decomposition $(C_H,m_H)$. For each vertex $v\in V_H$, let $S_v\subset S_G\setminus C$ and $S'_v\subset S_H\setminus C_H$ denote the corresponding subsurfaces, both of type $(0,3)$. The markings on $S_G$ and $S_H$ restrict to markings on $S_v$ and $S'_v$, which determine a unique mapping class \[\phi_v : S_v \longrightarrow S'_v\] as in Definition~\ref{def: diagramatic description of diffeo}. 

%Since the boundary parameterizations agree along the cut curves, the mapping classes $\phi_v$ are compatible along the glued boundary components and therefore assemble to a unique global mapping class    \[[\phi] : S_G \longrightarrow S_H.\] We define    \[I(C,m) := [\phi] \in \Hom_{\bS(g,n+1)}(G,H).\]

\medskip

We now define $I$ on morphisms. Let $\alpha:(C,[m])\longrightarrow (C',[m'])$ be a morphism in $\mathbb{M}_G$, and suppose
\[
I(C,[m])=[\phi:G\to H]
\qquad \text{and} \qquad
I(C',[m'])=[\phi':G\to H'].
\]
There is a unique morphism in the coslice category from $[\phi]$ to $[\phi']$, namely the mapping class
\[
[\phi'\phi^{-1}]:H\longrightarrow H'.
\]
We set $I(\alpha)=[\phi'\phi^{-1}]$. %Equivalently, $I(\alpha)$ is the unique morphism making the triangle
%\[
%\begin{tikzcd}
%& G \arrow[dl, "{[\phi]}"'] \arrow[dr, "{[\phi']}"] & \\
%H \arrow[rr, "{[\phi'\phi^{-1}]}"'] && H'
%\end{tikzcd}
%\]
%commute.
    

\medskip 
It remains to show that $I$ is an isomorphism. On objects, $I$ is bijective. Indeed, given an object $[\phi]:G\to H$ of $G\downarrow \bS(g,n+1)$, one can pull back the standard marked pants decomposition $(C_H,[m_H])$ of $S_H$ along $\phi$ to define a marked pants decomposition
\[
(\phi^{-1}(C_H),[\phi^{-1}m_H])
\] of $S_G$. This construction is inverse to the one above.

For morphisms, first observe that the coslice category $G\downarrow \bS(g,n+1)$ has exactly one morphism between any two objects. Indeed, if the objects are $[\phi]:G\to H$ and $[\phi']:G\to H'$, then a morphism between them is a morphism $[\psi]:H\to H'$ satisfying $[\psi]\circ[\phi]=[\phi']$. Since $\bS(g,n+1)$ is a groupoid, $[\phi]$ is invertible, so $[\psi]$ is necessarily $[\phi']\circ[\phi]^{-1}$. %Conversely, this morphism clearly satisfies the required relation.

On the other hand, $\mathbb{M}_G$ is the fundamental groupoid of $\calM(S_G)$ on its vertices. Since $\calM(S_G)$ is connected and simply connected by Proposition~\ref{prop:our-complex-of-markings-is-simply-connected}, there is exactly one homotopy class of paths between any two vertices. Thus $\mathbb{M}_G$ also has exactly one morphism between any two objects. Therefore $I$ is bijective on all Hom-sets. Hence $I$ is an isomorphism of categories. It follows that
\[
\mathbb{M}_G\cong G\downarrow \bS(g,n+1).
\]
\end{proof}

\medskip

Recall the mapping classes that described the standard $T$-, $B$-, $A$-, and $S$-moves in Definition~\ref{def:standard moves}. We write
\[
\tau \in \Hom_{\bS(0,2)}(|,|),\quad
\beta \in \Hom_{\bS(0,3)}\bigl(\triv,(12)^*\triv\bigr),\quad
\alpha \in \Hom_{\bS(0,4)}\bigl(\forkone,\forktwo\bigr),\quad
\sigma \in \Hom_{\bS(1,1)}\bigl(\contraction,\contraction\bigr)
\]
for the corresponding morphisms in $\bS$.

\begin{prop}\label{prop:S-generated-by-standard-moves}
Every morphism in $\bS$ can be obtained from $\tau$, $\beta$, $\alpha$, $\sigma$, and their inverses using categorical composition, operadic composition, symmetric group actions, and contractions.
\end{prop}

\begin{proof} 
Let $[\phi]:G\to H$ be a morphism in $\bS(g,n+1)$, represented by a boundary-label-preserving diffeomorphism
\[
\phi:S_G\longrightarrow S_H.
\]
By Lemma~\ref{lemma: the complexes of morphisms and markings}, this morphism corresponds to the marked pants decomposition of $S_G$ obtained by pulling back the standard marked pants decomposition of $S_H$ along $\phi$.

We first translate a single edge of the complex of markings into a morphism of $\bS$. Suppose that
\[
(C,[m])\rightsquigarrow (C',[m'])
\]
is an edge of $\calM(S_G)$ of type $T$, $B$, $A$, or $S$. Let
\[
[\phi]:G\to H
\qquad\text{and}\qquad
[\psi]:G\to H'
\]
be the objects of the coslice category $G\downarrow\bS(g,n+1)$ corresponding to $(C,[m])$ and $(C',[m'])$, respectively, under Lemma~\ref{lemma: the complexes of morphisms and markings}. The edge determines the unique morphism in the coslice category from $[\phi]$ to $[\psi]$; equivalently, there is a unique morphism $[f]:H\to H'$ such that
\[
[\psi]=[f]\circ[\phi].
\]
We claim that $[f]$ is obtained from identities and one of the standard generators $\tau,\beta,\alpha,\sigma$ by operadic composition, symmetric group actions, and contractions.


We verify the claim for a $B$-move; the other cases are analogous. Suppose that
\[
(C,[m]:H\hookrightarrow S_G)
\overset{B_{a,b,c}}{\rightsquigarrow}
(C',[m']:H'\hookrightarrow S_G)
\]
is a $B$-move. Then $C=C'$, so the two marked pants decompositions have the same underlying cut system. The move is supported on a single cut-subsurface
\[
S'\subset \overline{S_G\setminus C}
\]
with boundary components $a,b,c$, and the markings $m$ and $m'$ agree on the complement of $S'$.

Choose an identification $\xi:S'\xrightarrow{\cong}P$ with the standard pair of pants such that
\[
\xi\circ m|_{S'}=m_P,\qquad
\xi(a)=\partial_1,\qquad
\xi(b)=\partial_2,\qquad
\xi(c)=\partial_0.
\]
By the definition of the $B$-move, under this identification the new marking is given by
\[
\xi\circ m'|_{S'}=\beta\circ \xi\circ m|_{S'}.
\]

Let $v\in V(H)$ and $v'\in V(H')$ be the vertices corresponding to the component $S'$. Since the move is supported on $S'$, the graphs $H$ and $H'$ agree away from these vertices, with all boundary labels preserved. We therefore define a mapping class
\[
f:S_H\longrightarrow S_{H'}
\]
by taking it to be the identity on every pair-of-pants component except the one corresponding to $v$, and by taking it on that component to be the standard $B$-move $\beta$ under the above identifications. If the local identification with the standard pair of pants differs by a cyclic relabelling, then the local map is the corresponding cyclic translate of $\beta$.

Thus $[f]:H\to H'$ is obtained from $\beta$ by operadic composition with identities, together with the relevant symmetric group action. Moreover, the marking corresponding to $[\psi]$ is exactly the marking obtained from the one corresponding to $[\phi]$ by applying $[f]$ after $[\phi]$. Hence
\[
[\psi]=[f]\circ[\phi].
\]


We now complete the proof. Let $[\phi]\in \bS(g,n+1)(G,H)$, and let $(C_\phi,[m_\phi])$ be the corresponding marked pants decomposition of $S_G$. Since $\calM(S_G)$ is connected, there is an edge path
\[
(C_0,[m_0])\rightsquigarrow (C_1,[m_1])
\rightsquigarrow \cdots \rightsquigarrow (C_k,[m_k])
\]
from the standard marking $(C_0,[m_0])=(iE_G,[m_G])$ to $(C_k,[m_k])=(C_\phi,[m_\phi])$.

Under Lemma~\ref{lemma: the complexes of morphisms and markings}, the standard marking corresponds to the identity morphism $\id_G:G\to G$. By the single-edge argument above, each edge in the path corresponds to postcomposition with a morphism
\[
[f_i]:H_{i-1}\to H_i
\]
obtained from identities and one of $\tau,\beta,\alpha,\sigma$ by operadic composition, symmetric group actions, and contractions. Hence
\[
[\phi]=[f_k]\circ\cdots\circ [f_1]\circ \id_G.
\]
Thus $[\phi]$ lies in the subcategory generated by $\tau,\beta,\alpha,\sigma$ under categorical composition, operadic composition, symmetric group actions, and contractions. Since $[\phi]$ was arbitrary, the claim follows.
\end{proof}

%By definition of the complex $\calM(S_{G})$, a loop based at the standard cut system $(iE_G, m_G: G\hookrightarrow S_G)$ in $\calM_{S_G}$ is a sequence of edges \[(iE_G,m_G)\overset{e_1}{\rightarrow} (C_1,m_1) \overset{e_2}{\rightarrow}\dots \overset{e_{k-1}}{\rightarrow}(C_{k-1},m_{k-1})\overset{e_k}{\rightarrow} (iE_G,m_G)\] and, by Lemma~\ref{lemma: the complexes of morphisms and markings}, it corresponds to a sequence of diffeomorphisms \[\phi_{e_{k-1}}\phi_{e_{k-2}}\dots\phi_{e_1}=\id_{G} \in \Hom_{\bS(g,n+1)}(G,G).\]  Theorem 5.1 of \cite{bk_marked_surfaces} says that the complex $\calM(S_G)$ is connected and simply connected. In particular the fundamental group $\pi_1(\calM(S_G); (iE_G, m_G))$ is trivial and the boundaries of the $2$-cells based at the point $(iE_G, m_G)$ in $\calM(S_{G})$ translate to sequences of diffeomorphisms $\phi_{e_{k-1}}\phi_{e_{k-2}}\dots\phi_{e_1}=\id_{G} \in \Hom_{\bS(g,n+1)}(G,G)$ in which each of the $\phi_{e_i}$ are one of the elementary diffeomorphisms $\tau, \beta, \alpha$ and $\sigma$.  Translating the diffeomorphisms $\phi_{e_{k-1}}\phi_{e_{k-2}}\dots\phi_{e_1}=\id_{G} \in \Hom_{\bS(g,n+1)}(G,G)$ into a series of modular operadic and categorical compositions, provides a series of equations which are satisfied by the mapping classes $\tau, \beta, \alpha$ and $\sigma$ in $\bS(g,n+1)$. The following theorem states that the equations imposed by the $2$-cells of $\calM(S_{G})$ are sufficient to determine a map of modular operads out of $\bS$. 


The preceding proposition identifies the generators of the morphisms in $\bS$. The relations among these generators are encoded by the $2$-cells of the complex of markings.  Fix $G\in \ob(\bS(g,n+1))$. A loop in $\calM(S_G)$ based at the standard marking $(iE_G,[m_G])$ is represented by an edge path
\[
(iE_G,[m_G])
\overset{e_1}{\rightsquigarrow}
(C_1,[m_1])
\overset{e_2}{\rightsquigarrow}
\cdots
\overset{e_k}{\rightsquigarrow}
(iE_G,[m_G]).
\]
Under Lemma~\ref{lemma: the complexes of morphisms and markings}, this loop corresponds to a sequence of morphisms in $\bS(g,n+1)$ whose composite is the identity morphism of $G$:
\[
[f_k]\circ \cdots \circ [f_1]=\id_G
\qquad\text{in }\Aut_{\bS(g,n+1)}(G).
\]
Each edge $e_i$ is one of the elementary moves $T$, $B$, $A$, or $S$, so each $[f_i]$ is obtained from $\tau,\beta,\alpha,\sigma$ by operadic composition, symmetric group actions, and contractions.

The boundaries of the $2$-cells in $\calM(S_G)$ therefore give equations among
the generators $\tau,\beta,\alpha,\sigma$. The goal of the next few sections is
to show that these are the only equations satisfied by these generators. This
will give a presentation of $\bS$, and hence a classification of maps out of the
modular operad $\bS$.

Before proving this in full generality, we first study the genus-zero part of $\bS$ and relate it to operads whose presentations are already known.

\subsection{Genus-zero surfaces and cyclic ribbon braids} \label{sec: operad of ribbon braids}
As a first step toward a presentation of the modular operad of surfaces, we study its genus-zero truncation. We use the known presentation of the operad of parenthesized ribbon braids to obtain a presentation of $\trun{0}\bS$, and then show that this identification is compatible with the cyclic structures. %The key input is the classical identification of pure ribbon braid groups with mapping class groups of genus-zero surfaces with boundary.

\subsubsection{Braids and ribbon braids}

For $n\geq 1$, let $\Br_n$ denote the braid group on $n$ strands. It may be defined as the fundamental group of the configuration space of $n$ unordered distinct points in the complex plane. It admits the standard presentation
\[
\Br_n=
\left\langle
\beta_1,\ldots,\beta_{n-1}
\middle|
\beta_i\beta_j=\beta_j\beta_i \ \text{for } |i-j|\ge 2,
\beta_i\beta_{i+1}\beta_i=\beta_{i+1}\beta_i\beta_{i+1}
\right\rangle .
\]
For each $n\geq 1$, there is a short exact sequence
\[
\begin{tikzcd}
1 \arrow[r] & \PB_n \arrow[r] & \Br_n \arrow[r, "\pi"] & \Sigma_n \arrow[r] & 1,
\end{tikzcd}
\]
where $\pi$ sends the braid generator $\beta_i$ to the transposition $(i\ i+1)$. The kernel of $\pi$ is the \emph{pure braid group}, denoted $\PB_n$. It is generated by the elements
\[
x_{ij}=\beta_{j-1}\cdots \beta_{i+1}\beta_i^2\beta_{i+1}^{-1}\cdots \beta_{j-1}^{-1},
\qquad 1\le i<j\le n.
\]

The \emph{ribbon braid group} on $n$ strands, denoted $\RB_n$, may be defined as the fundamental group of the configuration space of $n$ unordered distinct points in the plane, each equipped with a framing (equivalently, each equipped a point of $S^1$). It is generated by the braid generators $\beta_1,\ldots,\beta_{n-1}$ together with twists $\twist_1,\ldots,\twist_n$, subject to the braid relations above and the additional relations
\[
\beta_i\twist_j=\twist_j\beta_i \quad \text{for } j\notin\{i,i+1\},\qquad
\beta_i\twist_{i+1}=\twist_i\beta_i,\qquad
\twist_i\twist_j=\twist_j\twist_i \quad \text{for } i\neq j.
\]
The \emph{pure ribbon braid group}, denoted $\PRB_n$, is the kernel of the natural projection
\[
\RB_n\longrightarrow \Sigma_n.
\]

The ribbon braid groups are closely related to mapping class groups of genus-zero surfaces. Indeed, there is a natural homomorphism
\[
\PRB_n\longrightarrow \Gamma_{0,n+1}
\]
sending the pure braid generator $x_{ij}$ to the Dehn twist about the simple closed curve separating the boundary components labelled $i$ and $j$ from the remaining boundary components, and sending the twist $\twist_k$ to the full Dehn twist about the $k$th boundary component. The following proposition is standard, c.f., for example, \cite[Section 1.5]{Wahl_Thesis}. 

\begin{prop}\label{prop: RBn is Gamma0n+1}
There is an isomorphism of groups
\[
\PRB_n\cong \Gamma_{0,n+1}.
\]
\end{prop}


\subsubsection{The operad of parenthesized ribbon braids} 
The objects of the operad of parenthesized ribbon braids are parenthesized binary products. We encode these products by planar rooted binary trees.% whose input legs are labelled.
\begin{definition}\label{def: magma operad}
The \emph{magma operad} $\Omega=\{\Omega(n)\}_{n\geq 1}$ is the operad in $\Set$ whose elements are isomorphism classes of planar rooted binary trees with $n$ input legs, equipped with a labelling
\[
\ell_T:\{0,1,\dots,n\}\xrightarrow{\cong}\partial(T),
\]
where $\ell_T(0)$ is the root.

Equivalently, $\Omega$ is the free operad generated by the two rooted binary corollas
\[
\mu=\mu(0;1,2)
\qquad\text{and}\qquad
(12)^*\mu=\mu(0;2,1).
\]
These are the two corollas obtained by decorating the rooted binary corolla by the two elements of $\Sigma_2$, or equivalently by choosing the order of the two input legs relative to the fixed root. Operadic composition is given by grafting the root of one tree onto an input leg of another. For further details, see \cite[Section~6.1]{FresseBook1} and \cite[Definition~6.10]{Boavida-Horel-Robertson}.
\end{definition}

\begin{remark}
Elements of $\Omega(n)$ may be identified with fully parenthesized words in the symbols $\{1,\dots,n\}$. Under this identification, the two generators $\mu$ and $(12)^*\mu$ correspond to the binary words $(12)$ and $(21)$, respectively, and every parenthesized word is obtained from these generators by operadic composition. For example,
\[
(((12)3)4), \qquad ((32)(14)) \qquad \text{and} \qquad (1((23)4))
\]
are elements of $\Omega(4)$.
\end{remark}

\begin{definition}
For each $n\geq 1$, let $\PaRB(n)$ be the groupoid whose objects are the elements of $\Omega(n)$. %Thus an object is a parenthesized word, or equivalently a labelled planar rooted binary tree with $n$ input legs.

Given two objects $T_1,T_2\in\Omega(n)$, their boundary labellings determine a unique bijection
\[
\rho_{T_1,T_2}:\partial(T_1)\xrightarrow{\cong}\partial(T_2)
\]
such that $\rho_{T_1,T_2}\circ \ell_{T_1}=\ell_{T_2}$. A morphism $r\in \Hom_{\PaRB(n)}(T_1,T_2)$ is a ribbon braid $r\in \RB_n$ whose underlying permutation is the permutation of the input legs induced by $\rho_{T_1,T_2}$.
\end{definition}

Composition in $\PaRB(n)$ is given by multiplication in the ribbon braid group $\RB_n$. In particular, for any object $T\in\Omega(n)$, the endomorphism group is naturally identified with the pure ribbon braid group:
\[
\Aut_{\PaRB(n)}(T)\cong \PRB_n.
\] The symmetric group $\Sigma_n$ acts on $\PaRB(n)$ by permuting the labels of the input legs. The resulting symmetric sequence
\[
\PaRB=\{\PaRB(n)\}_{n\geq 1}
\] forms an operad in groupoids. On objects, operadic composition is the grafting operation in $\Omega$. On morphisms, it is given by inserting a ribbon braid in $\RB_m$ into the $i$th strand of a ribbon braid in $\RB_n$. See \cite[Definition~6.11]{Boavida-Horel-Robertson} for details.

\medskip
An important feature of $\PaRB$ is that its morphisms are generated by three elementary morphisms:
\[
\tau\in\Hom_{\PaRB(1)}(\mid,\mid),\qquad
\beta\in\Hom_{\PaRB(2)}(\mu,(12)^*\mu),
\qquad
\alpha\in\Hom_{\PaRB(3)}(\mu\circ_1\mu,\mu\circ_2\mu).
\]
Here $\tau$ is the strand twist, $\beta$ is the braiding, and $\alpha$ is the re-parenthesization isomorphism from $\mu\circ_1\mu$ to $\mu\circ_2\mu$. More precisely, every morphism in $\PaRB(n)$ can be written as an iterated categorical and operadic composite of these morphisms and their inverses; see \cite[Section~7]{Boavida-Horel-Robertson}, following the general framework of \cite[Chapter~6]{FresseBook1}.  This presentation above allows one to classify maps out of $\PaRB$. 

The following is a restatement of Lemmas~7.1 and~7.4 of \cite{Boavida-Horel-Robertson}, based on \cite[Theorem~6.2.4]{FresseBook1}. We write $\cdot$ for categorical composition of morphisms in the groupoids $\PaRB(n)$, and $\circ_i$ for operadic composition.
\begin{thm}\label{thm: maps_out_of_P_0}
Let $\calQ$ be an operad in groupoids. To give an operad map
\[
f:\PaRB\longrightarrow \calQ
\]
is equivalent to giving a quadruple $(\mathbf m,\mathbf t,\mathbf b,\mathbf a)$ consisting of an object $f(\mu)=\mathbf m\in \ob(\calQ(2))$
and morphisms
\[
f(\tau)=\mathbf t\in \Hom_{\calQ(1)}(\mathbf 1,\mathbf 1),\qquad
f(\beta)=\mathbf b\in \Hom_{\calQ(2)}(\mathbf m,(12)^*\mathbf m) \qquad \text{and} \qquad f(\alpha)= \mathbf a\in \Hom_{\calQ(3)}(\mathbf m\circ_1 \mathbf m,\mathbf m\circ_2 \mathbf m),
\]
satisfying the relations below.

\begin{equation}\label{twist_rel_morph}\tag{T}
\mathbf t\circ_1 1_{\mathbf m}
=
\mathbf b\cdot (12)^*\mathbf b\cdot
\bigl(1_{\mathbf m}\circ (\mathbf t,\mathbf t)\bigr)
\qquad \text{in } \calQ(3).
\end{equation}

\begin{equation}\label{hex_rel_1_morph}\tag{H1}
\bigl((213)^*(1_{\mathbf m}\circ_{1}\mathbf b)\bigr)
\cdot \bigl((213)^*\mathbf a\bigr)
\cdot (1_{\mathbf m}\circ_{1}\mathbf b)
=
\bigl((231)^*\mathbf a\bigr)
\cdot (\mathbf b\circ_{2}1_{\mathbf m})
\cdot \mathbf a
\qquad \text{in } \calQ(4).
\end{equation}

\begin{equation}\label{hex_rel_2_morph}\tag{H2}
\bigl((132)^*(1_{\mathbf m}\circ_{1}\mathbf b)\bigr)
\cdot \bigl((132)^*\mathbf a^{-1}\bigr)
\cdot (1_{\mathbf m}\circ_{1}\mathbf b)
=
\bigl((312)^*\mathbf a^{-1}\bigr)
\cdot (\mathbf b\circ_{1}1_{\mathbf m})
\cdot \mathbf a^{-1}
\qquad \text{in } \calQ(4).
\end{equation}

\begin{equation}\label{pent_rel_1_morph}\tag{P}
(\mathbf a\circ_1 1_{\mathbf m})
\cdot (\mathbf a\circ_3 1_{\mathbf m})
=
(1_{\mathbf m}\circ_1 \mathbf a)
\cdot (\mathbf a\circ_2 1_{\mathbf m})
\cdot (1_{\mathbf m}\circ_2 \mathbf a)
\qquad \text{in } \calQ(5).
\end{equation}
\end{thm}


\medskip 

Recall from Proposition~\ref{prop:S-generated-by-standard-moves} that the
genus-zero generating morphisms in $\bS$ are
\[ 
\tau \in \Hom_{\bS(0,2)}(|,|),\quad \beta \in \Hom_{\bS(0,3)}\bigl(\triv,(12)\triv\bigr)\quad \text{and} \quad \alpha \in \Hom_{\bS(0,4)}\bigl(\forkone,\forktwo\bigr),
\]
where these are the mapping classes associated to the standard $T$-, $B$-, and $A$-moves in the complex of markings, respectively. To distinguish these mapping classes from the generating morphisms of $\PaRB$, throughout this subsection we denote the corresponding morphisms in $u^*\trun{0}\bS$ by
\[
\mathbf{t} \in \Hom_{\bS(0,2)}(|,|),\quad \mathbf{b} \in \Hom_{\bS(0,3)}\bigl(\triv,(12)^*\triv\bigr)\quad \text{and} \quad \mathbf{a} \in \Hom_{\bS(0,4)}\bigl(\forkone,\forktwo\bigr).
\]

The genus-zero truncation is naturally a cyclic operad, whose arity-$(n+1)$ component is $\trun{0}\bS(n+1)=\bS(0,n+1)$. Its underlying ordinary operad $u^*\trun{0}\bS$ is obtained by choosing the boundary component labelled $0$ as the root, so the same morphisms appear with one fewer input, i.e. $u^*\trun{0}\bS(n)=\trun{0}\bS(n+1)=\bS(0,n+1)$. The main difference in these groupoids are the actions by the symmetric groups. Notably, they all have the same objects. In particular, the two generating objects in arity $2$ are still $P$ and $(12)^*P$, corresponding to the two rooted planar corollas obtained from the two classes in $\Sigma_2^+/A_2^+$ after choosing the boundary labelled $0$ as the root. 

\begin{lemma}\label{lemma:operad-map-PaRB-to-S0}
There is an isomorphism of operads
\[
I:\PaRB\longrightarrow u^*\trun{0}\bS
\]
defined on generators by
\[
I(\mu)=P\in\ob\bigl(u^*\trun{0}\bS(2)\bigr),\quad I(\tau)=\mathbf{t}\in\Hom_{u^*\trun{0}\bS(1)}(|,|),
\]
\[
I(\beta)=\mathbf{b}\in\Hom_{u^*\trun{0}\bS(2)}\bigl(\triv,(12)^*\triv\bigr)\quad \text{and} \quad
I(\alpha)=\mathbf{a}\in\Hom_{u^*\trun{0}\bS(3)}\bigl(\forkone,\forktwo\bigr).
\]
\end{lemma}


\begin{proof}
By Theorem~\ref{thm: maps_out_of_P_0}, it is enough to check that the values assigned to the generators of the operad $\PaRB$ satisfy the defining relations \eqref{twist_rel_morph}, \eqref{hex_rel_1_morph}, \eqref{hex_rel_2_morph}, and \eqref{pent_rel_1_morph}.  

We first check the twist--braiding relation \eqref{twist_rel_morph}. Under the assignment in the statement of the lemma, this relation becomes
\begin{equation}\label{eq:I-applied-to-twist}
\mathbf t\circ_1 \id_P = \mathbf b\cdot (12)^*\mathbf b\cdot \bigl(\id_P\circ(\mathbf t,\mathbf t)\bigr)
\end{equation}
in $\Aut_{u^*\trun{0}\bS(2)}(P)$. We interpret this equation using the equivalence $\mathbb{M}_{P}\cong P\downarrow\bS(0,3)$ from Lemma~\ref{lemma: the complexes of morphisms and markings}. Under this equivalence, the left-hand side is the full Dehn twist about the boundary
component labelled $0$, which we denote by $T_0$. The right-hand side identifies with $B_{1,2}B_{2,1}T_1T_2$: here
$(12)^*B_{1,2}=B_{2,1}$, and the two inserted twists in $\id_P\circ(\mathbf t,\mathbf t)$ become the boundary twists $T_1$ and
$T_2$.

Thus \eqref{eq:I-applied-to-twist} is equivalent to $T_0 = B_{1,2}B_{2,1}T_1T_2.$  Since Dehn twists about boundary components commute, this is equivalent to
\[
B_{1,2}B_{2,1}=T_0T_1^{-1}T_2^{-1},
\]
which is precisely the twist--braiding $2$-cell in the marking complex $\mathbb{M}_{P}$. Hence the image of \eqref{twist_rel_morph} under $I$ holds.

We next check the pentagon relation \eqref{pent_rel_1_morph}. We rewrite it as the statement that
\[
(\alpha\circ_1 \id_{\mu})^{-1}\cdot (\alpha\circ_3 \id_{\mu})^{-1}\cdot (\id_{\mu} \circ_2 \alpha)\cdot (\alpha\circ_2 \id_{\mu})\cdot (\id_{\mu}\circ_1\alpha)=1
\] in $\Aut_{\PaRB(4)}((\mu\circ_1\mu)\circ_1\mu)$.

Let $\forkone=(P\circ_1P)\circ_1P$. After applying $I$, and then using the equivalence $\mathbb{M}_{S_{\forkone}}\cong \forkone\downarrow \bS(0,5)$ from Lemma~\ref{lemma: the complexes of morphisms and markings}, the left-hand side is identified with the loop of $A$-moves
\[
(A_{b,e})^{-1}(A_{a,d})^{-1}A_{c,e}A_{b,d}A_{a,c}
\]
in the complex of markings. The curves defining these $A$-moves are shown in Figure~\ref{fig:pentagon}. This loop is the boundary of the pentagon $2$-cell, and hence represents the identity morphism in $\mathbb{M}_{S_{\forkone}}$. Therefore its image is the identity in $\Aut_{u^*\trun{0}\bS(4)}(\forkone)$, so $I$ respects the pentagon relation \eqref{pent_rel_1_morph}.

The hexagon relations \eqref{hex_rel_1_morph} and \eqref{hex_rel_2_morph} are verified in the same way. Each may be rewritten as a closed composite in $\PaRB(3)(\mu\circ_1\mu,\mu\circ_1\mu)$ and $\PaRB(3)(\mu\circ_2\mu,\mu\circ_2\mu)$, respectively. Applying $I$, and then using the corresponding equivalences with the fundamental groupoids of the marking complexes, these automorphisms are identified with the two hexagon loops shown in Figure~\ref{fig:hexagons}. Each loop bounds a hexagon $2$-cell in the complex of markings. Hence both automorphisms are trivial, and $I$ respects the hexagon relations \eqref{hex_rel_1_morph} and \eqref{hex_rel_2_morph}. It follows from Theorem~\ref{thm: maps_out_of_P_0}
that the assignment on generators extends to a map of operads
\[
I:\PaRB\longrightarrow u^*\trun{0}\bS.
\]


It remains to show that this map is an isomorphism. On objects, $I$ is the identity. For every object $T\in\PaRB(n)$, the induced map on automorphism groups is the standard identification $\PRB_n\cong \Gamma_{0,n+1}$ from Proposition~\ref{prop: RBn is Gamma0n+1}.
Thus each functor $I(n):\PaRB(n)\longrightarrow u^*\trun{0}\bS(n)$ is fully faithful. Since it is also bijective on objects, it is an isomorphism of groupoids for each $n$. Therefore $I$ is an isomorphism of operads.
\end{proof}


\medskip 


The operad $\PaRB$ admits a cyclic structure in the sense of Definition~\ref{def: cyclic structure}; see \cite{campos2019configuration,MW,robertson2025grothendieck}. We write $\PaRB^{\mathrm{cyc}}$ for the resulting cyclic operad.

It is enough to describe the action of the transposition $(01)$, since this action, together with the ordinary symmetric group actions, determines the full extended symmetric group action. On objects, $(01)$ changes the choice of root: the leg labelled $1$ becomes the new root. In particular, on the generating object one has $(01)^*\mu=(12)^*\mu.$  Since $\PaRB$ admits a finite presentation,  it remains to specify the action of $(01)$ on the generating morphisms. This action is given by
\begin{equation}\label{cyclic action on PaRB generators}\tag{$\star$}
(01)^*(\tau)=\tau,\qquad (01)^*(\beta)=\beta^{-1}\cdot(\id_{\mu}\circ_2\tau^{-1}) \qquad \text{and} \qquad (01)^*(\alpha)=(231)^*\alpha^{-1}.
\end{equation}


We now show that the operad isomorphism constructed in Lemma~\ref{lemma:operad-map-PaRB-to-S0} respects this root-changing action, and therefore upgrades to an isomorphism of cyclic operads.

\begin{prop}\label{prop:PaRB-cyclic-is-S0}
The isomorphism of operads $I:\PaRB\to u^*\trun{0}\bS$ is compatible with the cyclic structure on $\PaRB$. Hence it induces an isomorphism of cyclic operads
\[
\PaRB^{\mathrm{cyc}}\cong \trun{0}\bS.
\]
\end{prop}

\begin{proof}
Since the $(01)$-action determines the cyclic structure on $\PaRB$, it suffices to verify $I((01)^*x)=(01)^*I(x)$ for $x=\mu,\tau,\beta,\alpha$.


For the generating object, recall that $I(\mu)=P$, where $P$ is the rooted trivalent corolla with boundary labels $\{0,1,2\}$. The objects of $\trun{0}\bS$ are generated by the two rooted planar corollas represented by $P=\id\cdot A_2^+$ and $(12)^*P=(12)\cdot A_2^+$, where $A_2^+=\langle(012)\rangle$. Since $(01)$ and $(12)$ represent the same nontrivial coset modulo $A_2^+$, we have $(01)^*P=(12)^*P$. Therefore $I((01)^*\mu)=(01)^*I(\mu)$.

For the twist generator, $I(\tau)=\mathbf{t}$. Recall that, by definition, the transposition
$(01)$ acts trivially on $\bS(0,2)$. Hence \[ I((01)^*\tau)=I(\tau)=\mathbf{t}=(01)^*I(\tau). \]

For the braiding generator, we compare the cyclic translate of $\mathbf b:P\to(12)^*P$ using the marking complex for the correct source object. Applying $(01)^*$ gives a morphism $(01)^*\mathbf{b}:(12)^*P\longrightarrow P.$
Through the equivalence
\[
\mathbb{M}_{(12)^*P}\cong (12)^*P\downarrow \trun{0}\bS(3),
\]
the morphism $(01)^*\mathbf{b}$ is identified with the opposite $B$-move, together with the boundary-twist correction needed to rewrite the $(01)$-relabelled marking in standard form. With our conventions this gives
\[
(01)^*B_{1,2}=T_0^{-1}T_1\cdot B_{2,1}.
\]
Using the twist--braiding relation in the marking complex,
\[
B_{1,2}B_{2,1}=T_0T_1^{-1}T_2^{-1},
\]
we obtain
\[
(01)^*B_{1,2} = T_0^{-1}T_1\cdot B_{1,2}^{-1}T_0T_1^{-1}T_2^{-1} = B_{1,2}^{-1}T_2^{-1}.
\]
The twist $T_2^{-1}$ is the twist inserted on the second input. In operadic notation, this says
\[
(01)^*\mathbf{b} = \mathbf{b}^{-1}\cdot(\id_P\circ_2\mathbf{t}^{-1}).
\]
This is exactly the image under $I$ of the relation $(01)^*\beta=\beta^{-1}\cdot(\id_{\mu}\circ_2\tau^{-1}).$  Therefore $I((01)^*\beta)=(01)^*I(\beta)$.


It remains to check the associativity generator. Applying $(01)^*$ to the morphism $\mathbf{a}:P\circ_1P\longrightarrow P\circ_2P$ in $\trun{0}\bS(4)$   gives a morphism
\[
(01)^*\mathbf{a}: (01)^*(P\circ_1P)\longrightarrow (01)^*(P\circ_2P).
\]
A direct calculation on labelled rooted corollas gives
\[
(01)^*(P\circ_1P)=(231)^*(P\circ_2P)
\qquad \text{and} \qquad (01)^*(P\circ_2P)=(231)^*(P\circ_1P).
\]
Thus both $(01)^*\mathbf{a}$ and $(231)^*\mathbf{a}^{-1}$ are morphisms
\[
(231)^*(P\circ_2P)\longrightarrow (231)^*(P\circ_1P).
\]

To compare them, we compose with $(231)^*\mathbf{a}$ and consider the resulting
automorphism:
\[
(01)^*\mathbf{a}\cdot (231)^*\mathbf{a}: (231)^*(P\circ_2P)\rightarrow (231)^*(P\circ_2P). 
\]
Passing through the equivalence $\mathbb M_{(231)^*(P\circ_2P)} \cong (231)^*(P\circ_2P)\downarrow \trun{0}\bS(4),$  this loop is identified with the loop in the marking complex obtained by performing the $(01)$-relabelled $A$-move and then the corresponding
$(231)$-relabelled $A$-move back. This is precisely the self-duality loop for the $A$-move. By the self-duality relation
\[
A_{c,c'}A_{c',c}=\id,
\]
the loop is trivial. Hence $(01)^*\mathbf{a}=(231)^*\mathbf{a}^{-1}.$  Therefore $I((01)^*\alpha)=(01)^*I(\alpha)$.

It follows that $I$ is compatible with the cyclic structures and thus induces an isomorphism $\PaRB^{\mathrm{cyc}}\cong \trun{0}\bS.$ 
\end{proof}


\begin{remark}
In the proof of Lemma~\ref{lemma:operad-map-PaRB-to-S0}, we freely used diagrammatic depictions of genus-zero mapping classes and their corresponding paths in the complexes of markings. This is justified by Lemma~\ref{lemma: the complexes of morphisms and markings}, together with the discussion in Section~\ref{subsubsec: diagrammatical representation of mapping classes}, which shows that these diagrammatic representatives are well defined. For example, Figure~\ref{fig: cyclic} depicts the action of the $(01)$-transposition on the $B$- and $A$-moves.

\begin{figure}[h!t]
    \centering
    \includegraphics[width=\textwidth]{Figures/Cyclic_B.pdf}
    
    \vspace{.5cm}

    \includegraphics[width=\textwidth]{Figures/Cyclic_A.pdf}
    \caption{Action of the $(01)$-transposition on the $B$- and $A$-moves.}
    \label{fig: cyclic}
\end{figure}
\end{remark}







\subsection{A presentation for maps out of $\bS$}
\label{sec: maps out of bS}

We use $\cdot$ for categorical composition in the groupoids $\bS(g,n+1)$, write $\circ_j^i$ for cyclic operadic composition, and write $\xi_i^j$ for
contractions. When only the underlying operad structure is involved, we use the convention $\circ_j^0=\circ_j$.

The goal of this subsection is to prove the main presentation theorem for maps out of the modular operad of surfaces $\bS$.  As in the genus-zero case, the proof uses the comparison between the groupoids $\bS(g,n+1)$ and the $2$-dimensional CW complexes of reduced marked cut systems on surfaces of type $(g,n+1)$. This comparison is established in Lemma~\ref{lemma: the complexes of morphisms and markings}. It allows us to apply Proposition~\ref{prop:our-complex-of-markings-is-simply-connected}, which states that $\calM(S)$ is connected and simply connected, and hence to obtain a presentation for the morphisms of $\bS(g,n+1)$.


















%\subsubsection{A presentation of the modular operad $\bS$} Proof of the main theorem. 




\begin{thm}\label{thm: presentation for maps out of bS}
Let \(\calQ\) be a modular operad in groupoids. A map of modular operads \(f:\bS\longrightarrow \calQ\) is uniquely determined by a tuple
\[
(\mathbf{m},\mathbf{t},\mathbf{b},\mathbf{a},\mathbf{s})
\]
consisting of an object $f(P)=\mathbf{m}\in \ob(\calQ(0,3))$ 
and morphisms
\[
f(\tau)=\mathbf{t}\in \Hom_{\calQ(0,2)}(|,|),\quad
f(\beta)=\mathbf{b}\in
\Hom_{\calQ(0,3)}(\mathbf{m},(12)^*\mathbf{m}),
\]
\[
f(\alpha)=\mathbf{a}\in
\Hom_{\calQ(0,4)}
(\mathbf{m}\circ_1\mathbf{m},\mathbf{m}\circ_2\mathbf{m}),
\quad
f(\sigma)=\mathbf{s}\in
\Hom_{\calQ(1,1)}(\xi_{12}\mathbf{m},\xi_{12}\mathbf{m}),
\]
satisfying the following relations. %All equations below are equations between morphisms in the indicated groupoid.

\begin{equation}\label{sym_rel_morph}\tag{S}
(012)^*(\mathbf{m})=\mathbf{m}
\qquad \text{in } \ob(\calQ(0,3)).
\end{equation}
%where $z_3=(012)$ is the cyclic permutation of the three boundary labels.

\begin{equation}\label{twist_rel_morphS}\tag{T}
(\mathbf{t}\circ_1 \id_{\mathbf{m}})
=
\mathbf{b}\cdot (12)^*\mathbf{b}\cdot
(\id_{\mathbf{m}}\circ(\mathbf{t},\mathbf{t}))
\qquad \text{in } \calQ(0,3).
\end{equation}

\begin{equation}\label{hex_rel_1_morphS}\tag{H1}
((213)^*(\id_{\mathbf{m}}\circ_{1}^{0}\mathbf{b}))
\cdot ((213)^*\mathbf{a})
\cdot (\id_{\mathbf{m}}\circ_{1}^{0}\mathbf{b})
=
((231)^*\mathbf{a})
\cdot (\mathbf{b}\circ_{2}^{0}\id_{\mathbf{m}})
\cdot \mathbf{a}
\qquad \text{in } \calQ(0,4).
\end{equation}

\begin{equation}\label{hex_rel_2_morphS}\tag{H2}
((132)^*(\id_{\mathbf{m}}\circ_{1}^{0}\mathbf{b}))
\cdot ((132)^*\mathbf{a}^{-1})
\cdot (\id_{\mathbf{m}}\circ_{1}^{0}\mathbf{b})
=
((312)^*\mathbf{a}^{-1})
\cdot (\mathbf{b}\circ_{1}^{0}\id_{\mathbf{m}})
\cdot \mathbf{a}^{-1}
\qquad \text{in } \calQ(0,4).
\end{equation}

\begin{equation}\label{pent_rel_1_morphS}\tag{P}
(\mathbf{a}\circ_1 \id_{\mathbf{m}})
\cdot(\mathbf{a}\circ_3 \id_{\mathbf{m}})
=
(\id_{\mathbf{m}}\circ_1\mathbf{a})
\cdot (\mathbf{a}\circ_2 \id_{\mathbf{m}})
\cdot (\id_{\mathbf{m}}\circ_2\mathbf{a})
\qquad \text{in } \calQ(0,5).
\end{equation}

\begin{equation}\label{g1_rel_morph_1}\tag{G1a}
\mathbf{s}^2
=
\xi_{0,2}\bigl(\mathbf{b}^{-1}\bigr)
\qquad \text{in } \calQ(1,1).
\end{equation}

\begin{equation}\label{g1_rel_morph_2}\tag{G1b}
\mathbf{d}^{-1}
\cdot \mathbf{s}
\cdot \mathbf{d}^{-1}
\cdot \mathbf{s}
\cdot \mathbf{d}^{-1}
=
\mathbf{s}
\qquad \text{in } \calQ(1,1),
\end{equation}
where $\mathbf{d}=\xi_{12}(\id_{\mathbf{m}}\circ_1\mathbf{t}).$

\begin{multline}\label{g2_rel_morph_1}\tag{G2}
\xi_{12}\bigl((012)^*\mathbf{b}\circ_2
(\mathbf{t}^{-1},\mathbf{t}^{-1})\circ \mathbf{b}\bigr)
\\
=
\xi_{12}(\mathbf{a})
\cdot (\id_{\mathbf{m}}\circ_1\mathbf{s})
\cdot \xi_{12}(\mathbf{a}^{-1})
\cdot
\xi_{12}\bigl(\id_{\mathbf{m}}\circ_2
(\mathbf{t}^{-1},\mathbf{t})\circ\id_{\mathbf{m}}\bigr)
\\
\cdot \xi_{12}(\mathbf{a})
\cdot (\id_{\mathbf{m}}\circ_1\mathbf{s})
\cdot \xi_{12}(\mathbf{a}^{-1})
\qquad \text{in } \calQ(1,2).
\end{multline}
\end{thm}





Relations \ref{g1_rel_morph_1}, \ref{g1_rel_morph_2} and \ref{g2_rel_morph_1} are added to include the higher genus structure. For instance, 
equation~\ref{g2_rel_morph_1} can be pictured diagrammatically as in Figure~\ref{fig:2-cell from NS}.

\begin{figure}
    \centering
    \includegraphics[width=\linewidth]{Figures/2-cell-from-NS2.pdf}
    \caption{Diagram of relation (G2)}
    \label{fig:2-cell from NS}
\end{figure}





\begin{proof}
A map of modular operads $f:\bS\to \calQ$ necessarily determines the images of the generating morphisms:
\[ f(\mu)=\mathbf{m},\quad f(\tau)=\mathbf{t},\quad f(\beta)=\mathbf{b},\quad f(\alpha)=\mathbf{a},\quad f(\sigma)=\mathbf{s}.
\]
Moreover, such a map consists of functors $f(g,n+1):\bS(g,n+1)\longrightarrow \calQ(g,n+1)$ compatible with the symmetric group actions, operadic compositions, and contractions. Under the equivalence of Lemma~\ref{lemma: the complexes of morphisms and markings}, the boundaries of the relevant $2$-cells in the marking complexes $\calM(S_G)$ are identified with the relations \eqref{sym_rel_morph}, \eqref{twist_rel_morphS}, \eqref{pent_rel_1_morphS}, \eqref{hex_rel_1_morphS}, \eqref{hex_rel_2_morphS}, \eqref{g1_rel_morph_1}, \eqref{g1_rel_morph_2}, and \eqref{g2_rel_morph_1}. Therefore any map $f:\bS\to\calQ$ gives a tuple $(\mathbf{m},\mathbf{t},\mathbf{b},\mathbf{a},\mathbf{s})$ satisfying the stated relations.

\medskip 

We now prove the converse. Suppose we are given $\mathbf{m},\mathbf{t},\mathbf{b},\mathbf{a},\mathbf{s}$ in $\calQ$ satisfying the stated relations.  Since $\ob(\bS)$ is the free modular operad in sets generated by the genus-zero corolla $P$ and $(12)^*P$, the assignment $\ob(f)(P)=\mathbf{m}\in\ob(\calQ(0,3))$ extends uniquely to a map of modular operads in sets
\[
\ob(f):\ob(\bS)\longrightarrow\ob(\calQ).
\]
In particular, the object assignment is compatible with the extended symmetric group actions, operadic compositions, and contractions.  Note that since the object $P\in \bS(0,3)$ is represented by the coset $\id\cdot A_2^+$, where $A_2^+=\langle(012)\rangle$. Thus $(012)^*P=P$, and the condition
\[
(012)^*\mathbf{m}=\mathbf{m}
\]
is precisely the condition needed for the assignment $P\mapsto \mathbf{m}$ to respect the stabiliser of the generating object.

\medskip

Using the cyclic operad isomorphism $\PaRB^{\mathrm{cyc}}\cong \trun{0}\bS$ from Proposition~\ref{prop:PaRB-cyclic-is-S0}, together with Theorem~\ref{thm: maps_out_of_P_0}, the assignments
\[
\mu\mapsto P,\qquad \tau\mapsto \mathbf{t},\qquad \beta\mapsto \mathbf{b},\qquad \text{and} \qquad \alpha\mapsto \mathbf{a}
\]
satisfying \eqref{twist_rel_morphS}, \eqref{pent_rel_1_morphS}, \eqref{hex_rel_1_morphS}, and \eqref{hex_rel_2_morphS} define the genus-zero
part of $f$, namely a map of cyclic operads $\trun{0}f:\trun{0}\bS\longrightarrow \trun{0}\calQ.$

\medskip 

It remains to extend the genus-zero map to all of $\bS$. By Proposition~\ref{prop:S-generated-by-standard-moves}, every morphism in $\bS$ is obtained from the generating morphisms $\alpha,\beta,\tau,\sigma$ by modular operad operations and categorical composition. Thus, for a morphism
\[
\phi\in\Hom_{\bS(g,n+1)}(G,G'),
\]
we choose an expression for $\phi$ in terms of these generators and define its image in $\calQ(g,n+1)$ by replacing $\alpha,\beta,\tau,$ and $\sigma$ by $\mathbf{a},\mathbf{b},\mathbf{t}$ and $\mathbf{s}$. 

We must check that this description of $f(\phi)$ is independent of the chosen expression. So suppose that $\phi$ has two expressions
\[
\phi=\phi^1_k\cdots \phi^1_1=\phi^2_m\cdots \phi^2_1,
\]
in $\bS(g,n+1)$, where each $\phi^r_i$ is one of the generating morphisms of $\bS$. Then
\[
(\phi^2_m\cdots \phi^2_1)^{-1}(\phi^1_k\cdots \phi^1_1)=\id_G \ \in\Aut_{\bS(g,n+1)}(G).
\]

Under the equivalence of Lemma~\ref{lemma: the complexes of morphisms and markings}, this relation is represented by a loop in the marking complex $\calM(S_G)$. Since $\calM(S_G)$ is connected and simply connected \cite[Theorem~5.1]{bk_marked_surfaces}, this loop is generated by the boundaries of the $2$-cells of $\calM(S_G)$.  The boundary relations of these $2$-cells are exactly the relations imposed above, together with the formal relations already forced by the modular operad axioms. 

More precisely, the twist--braiding, pentagon, and two hexagon cells are encoded by \eqref{twist_rel_morphS}, \eqref{pent_rel_1_morphS}, \eqref{hex_rel_1_morphS}, and \eqref{hex_rel_2_morphS}. The self-duality of the associativity move follows from Mac Lane coherence, and therefore from the pentagon and hexagon relations already imposed.  The genus-one cells in types $(1,1)$ and $(1,2)$ are encoded by \eqref{g1_rel_morph_1}, \eqref{g1_rel_morph_2}, and \eqref{g2_rel_morph_1}. The remaining cells express formal compatibility relations for disjointly supported moves and for the modular operad operations. These hold in $\calQ$ because $\calQ$ is a modular operad. 

It follows that
\[
(f(\phi^2_m)\cdots f(\phi^2_1))^{-1} (f(\phi^1_k)\cdots f(\phi^1_1)) = \id_{f(G)}.
\]
Thus the two expressions for $\phi$ have the same image in $\calQ$, so the assignment on morphisms is well defined. By construction it is compatible with categorical composition, symmetric group actions, operadic compositions, and contractions. Hence it defines a map of modular operads $f:\bS\longrightarrow \calQ.$ The theorem follows. 

\end{proof}



\section{The operadic two-level principle}
\label{sec: two level principle}

Grothendieck's \emph{two-level principle} predicts that Galois actions on the Teichmüller tower should be determined by their effect in the first two levels,
namely on
\[
\widehat{\Gamma}_{0,4},\quad
\widehat{\Gamma}_{0,5},\quad
\widehat{\Gamma}_{1,1},\quad
\text{and}\quad
\widehat{\Gamma}_{1,2}.
\]
In this section we formulate and prove an operadic version of this principle. Our main goal is to show that there exists an action of $\galQ$ on the profinite completion of the surface modular operad $\widehat{\bS}$ (Theorem~\ref{main theorem NS action}). Moreover, this action is completely determined by the genus one truncation of $\widehat{\bS}$ (Corollary~\ref{cor: truncation and endomorphisms of the surface mod operad}). 

We begin with a brief discussion of profinite completion for modular operads in groupoids. 

\subsection{Profinite completion of modular operads in groupoids}\label{sec: profinite completion of modular operads in groupoids}
Proposition~\ref{prop: profinite completion preserves products of groupoids} allows us to define profinite completion entrywise for modular operads in groupoids satisfying the relevant finiteness hypotheses. 

\begin{definition}\label{def: profinite completion of modular operads in groupoids}
Let $\calP=\{\calP(g,n+1)\}$ be a modular operad in groupoids such that each $\calP(g,n+1)$ has finitely many objects. Its profinite completion is the modular operad in profinite groupoids
\[
\widehat{\calP}=\{\widehat{\calP}(g,n+1)\}.
\]
The structure maps are obtained by applying profinite completion to the structure maps of $\calP$, using the product comparison to identify products of completions with completions of products.
\end{definition}

More precisely, the partial composition
\[
\widehat{\calP}(g,n+1)\times \widehat{\calP}(h,m+1)
\longrightarrow
\widehat{\calP}(g+h,n+m)
\]
is the composite
\[
\begin{tikzcd}
\widehat{\calP}(g,n+1)\times \widehat{\calP}(h,m+1)
\arrow[r, "\cong"]
&
\reallywidehat{\calP(g,n+1)\times \calP(h,m+1)}
\arrow[r, "\widehat{(\circ_i^j)}"]
&
\widehat{\calP}(g+h,n+m).
\end{tikzcd}
\]
Here the first map is the inverse of the product isomorphism from Proposition~\ref{prop: profinite completion preserves products of groupoids}, and the second map is obtained by applying profinite completion to the composition map of $\calP$. The contraction maps
\[
\widehat{(\xi_{i}^{j})}:\widehat{\calP}(g,n+1)\longrightarrow \widehat{\calP}(g+1,n-1)
\]
are defined by applying profinite completion to the contraction maps of $\calP$.


\subsubsection{The genus-one truncation}
The presentation theorem (Theorem~\ref{thm: presentation for maps out of bS}) shows that maps out of $\mathbf S$ are determined by the images of the generators and relations appearing in genus at most one. We express this using the truncation formalism introduced in Section~\ref{subsec: intro to truncations} and developed further in Appendix~\ref{subsec: truncations}: $\trun{1}\calP$ denotes the genus-one truncation of a modular operad $\calP$, and the truncation functor admits a left adjoint $(\trun{1})_!$.

\begin{cor}\label{cor: maps out of bS}
The counit of the truncation adjunction gives an isomorphism of modular operads in groupoids
\[
(\trun{1})_!(\trun{1}\bS)\cong \bS.
\]
\end{cor}

\begin{proof}
By Theorem~\ref{thm: presentation for maps out of bS}, a map out of $\bS$ is uniquely determined by the images of the generators appearing in $\trun{1}\bS$, and every choice of such images satisfying the stated relations extends uniquely to a map out of $\bS$. Thus, for every modular operad $\calQ$, restriction induces a natural bijection
\[
\Map_{\Mod(\Grpd)}(\bS,\calQ) \cong \Map_{\Mod_{\leq 1}(\Grpd)}(\trun{1}\bS,\trun{1}\calQ).
\]
By the adjunction $(\trun{1})_!\dashv \trun{1}$, the right-hand side is naturally isomorphic to $\Map_{\Mod(\Grpd)}((\trun{1})_!(\trun{1}\bS),\calQ).$ Under these identifications, the natural bijection is induced by precomposition with the counit
\[
(\trun{1})_!(\trun{1}\bS)\longrightarrow \bS.
\]
Therefore, by Yoneda, the counit is an isomorphism.
\end{proof}


We will use this after profinite completion.  Let $\eta_{\calP}:\calP\to |\widehat{\calP}|$ denote the unit of the profinite-completion adjunction, where $|-|:\widehat{\Grpd}\to\Grpd$ is the forgetful functor. Thus, for any modular operad $\mathcal R$ in profinite groupoids, composition with $\eta_{\calP}$ gives a natural bijection
\[
\Map_{\Mod(\widehat{\Grpd})}(\widehat{\calP},\mathcal R) \cong \Map_{\Mod(\Grpd)}(\calP,|\mathcal R|).
\]
Equivalently, every map $f:\calP\to|\mathcal R|$ extends uniquely to a continuous map $\widetilde f:\widehat{\calP}\to\mathcal R$ with $|\widetilde f|\circ\eta_{\calP}=f$.

In particular, taking $\calP=\bS$ and $\mathcal R=\widehat{\bS}$, the presentation theorem describes maps $\bS\to|\widehat{\bS}|$, and the adjunction upgrades each such map uniquely to a continuous endomorphism $\widehat{\bS}\to\widehat{\bS}$. 


\begin{cor}\label{cor: truncation and endomorphisms of the surface mod operad}
The truncation functor $\trun{1}$ induces an isomorphism of monoids
\[
\End_0(\widehat{\bS})
\xrightarrow{\cong}
\End_0(\trun{1}\widehat{\bS}).
\]
\end{cor}

\begin{proof}
By the profinite-completion adjunction, endomorphisms of $\widehat{\bS}$ are equivalent to maps $\bS\to|\widehat{\bS}|$. Since the unit $\eta_{\bS}:\bS\to|\widehat{\bS}|$ does not change the object set, this bijection restricts to object-fixing maps:
\[
\End_0(\widehat{\bS})\cong \Map_0(\bS,|\widehat{\bS}|).
\]
Corollary~\ref{cor: maps out of bS} identifies this set with $\Map_0(\trun{1}\bS,|\trun{1}\widehat{\bS}|)$. Since profinite completion is applied entrywise, it commutes with genus truncation, so $\trun{1}\widehat{\bS}$ is the profinite completion of $\trun{1}\bS$. A second application of the profinite-completion adjunction gives
\[
\Map_0(\trun{1}\bS,|\trun{1}\widehat{\bS}|) \cong \End_0(\trun{1}\widehat{\bS}).
\]
Combining these bijections gives the desired isomorphism of monoids.
\end{proof}
\subsection{The action of $\galQ$ on $\widehat{\bS}$}\label{sec: NS}
To study Galois actions on higher-genus mapping class groups, we work with explicit profinite groups containing $\galQ$ and admitting well-understood presentations. The most prominent example is the Grothendieck--Teichmüller group $\GT$, introduced by Drinfeld \cite[Section~4]{Drin}. 
In their study of Galois actions on higher-genus mapping class groups, Nakamura and Schneps introduced in \cite{Nakamura-Schneps} a subgroup $\NS$ of the Grothendieck--Teichmüller group, satisfying
\[
\galQ\hookrightarrow \NS\hookrightarrow \GT.
\]
Rather than defining the Galois action directly on the full profinite modular operad $\widehat{\bS}$, we construct it via these explicit profinite groups.

Let $\widehat{\F}_2$ denote the profinite completion of the free group on two generators $x$ and $y$. We use the following standard convention: if $f=f(x,y)\in \widehat{\F}_2$ and $a,b$ are elements of a profinite group $\mathsf G$, then $f(a,b)$ denotes the image of $f$ under the unique
continuous homomorphism $\widehat{\F}_2\longrightarrow \mathsf G$ sending $x\mapsto a$ and $y\mapsto b$.

\begin{definition}\label{defn:GT}
The \emph{Grothendieck--Teichmüller monoid} $\underline{\GT}$ consists of pairs $(\lambda,f)\in \widehat{\mathbb Z}\times \widehat{\F}_2$ such that the assignment
\[
x\longmapsto f(x,y)^{-1}x^\lambda f(x,y),
\qquad
y\longmapsto y^\lambda
\]
defines a continuous endomorphism $\varphi_{(\lambda,f)}\colon \widehat{\F}_2\longrightarrow \widehat{\F}_2$
and the following equations hold:
\begin{equation}\label{I}\tag{I}
f(x,y)f(y,x)=1,
\end{equation}
\begin{equation}\label{II}\tag{II}
f(z,x)z^\nu f(y,z)y^\nu f(x,y)x^\nu=1,
\qquad
z=(xy)^{-1},
\qquad
\nu=(\lambda-1)/2.
\end{equation}
and
\begin{equation}\label{III}\tag{III}
f(x_{12},x_{23}x_{24})f(x_{13}x_{23},x_{34})
=
f(x_{23},x_{34})f(x_{12}x_{13},x_{24}x_{34})f(x_{12},x_{23}).
\end{equation}
The first two equations hold in \(\widehat F_2\), and the last equation holds
in the profinite completion of the pure braid group \(\widehat{\PB}_4\). The multiplication is defined by composition of the associated endomorphisms.
Thus, for \(F=(\lambda,f)\) and \(G=(\mu,g)\), we set
\[
F\cdot G
=
(\lambda,f)\cdot(\mu,g)
:=
\left(
\lambda\mu,\,
g(x,y)\,
f\bigl(g(x,y)^{-1}x^\mu g(x,y),\,y^\mu\bigr)
\right),
\]
so that
\[
\varphi_{F\cdot G}=\varphi_G\circ\varphi_F.
\]
The Grothendieck--Teichmüller group \(\GT\) is the group of invertible elements of the monoid \(\underline{\GT}\).

%The multiplication is given by\[(\lambda,f)\cdot(\mu,g) = \bigl(\lambda\mu, f(gx^\mu g^{-1},y^\mu)\cdot g\bigr).\] The Grothendieck--Teichmüller group $\GT$ is the group of invertible elements of $\underline{\GT}$.
\end{definition}


\begin{remark}
The reader familiar with $\GT$ might note that our multiplication convention on $\GT$ differs slightly from that of \cite{Drin}. We follow the composition convention used in \cite{hls}, where the product is defined so that the associated endomorphisms of $\widehat{\F}_2$ compose by $\varphi_{F\cdot G}=\varphi_G\circ\varphi_F.$
\end{remark}

Drinfeld described the action of $\GT$ on profinite braid groups in his original definition of the Grothendieck--Teichmüller group~\cite{Drin}. An operadic interpretation of this action was developed in~\cite{Horel_profinite_groupoids}, building on work such as~\cite{BN,FresseBook1}. This was extended to an action on profinite ribbon braid operads by the second author and collaborators in~\cite{Boavida-Horel-Robertson}. The compatibility with the cyclic structure on $\PaRB$ was recently established by the second author and Singh~\cite{robertson2025grothendieck}.

By Proposition~\ref{prop:PaRB-cyclic-is-S0}, there is an isomorphism of cyclic operads
\[
\trun{0}\widehat{\bS}\cong \widehat{\PaRB}^{\mathrm{cyc}}.
\]
We use this identification to regard the genus-zero $\GT$-action on $\widehat{\PaRB}^{\mathrm{cyc}}$ as an action on $\trun{0}\widehat{\bS}$. Explicitly, an element $F=(\lambda,f)\in\underline{\GT}$ determines an object-fixing endomorphism $\iota_F\in\End_0(\trun{0}\widehat{\bS})$ by its action on the standard generators $\tau,\beta,\alpha$:
\[
\iota_F(\tau)=\tau^\lambda,\qquad
\iota_F(\beta)=\beta\cdot\bigl(\beta\cdot(12)^*\beta\bigr)^\nu,\qquad
\iota_F(\alpha)=\alpha\cdot f(D_{c'},D_c),
\]
where $\nu=(\lambda-1)/2$, and where $c$ and $c'$ are the curves shown in Figure~\ref{fig:A-move}.

\begin{prop}\label{prop:genus-zero-GT}
There is an isomorphism of monoids
\[
\End_0(\trun{0}\widehat{\bS})
\cong
\underline{\GT}.
\]
\end{prop}

\begin{proof}
This follows from the identification $\trun{0}\widehat{\bS}\cong \widehat{\PaRB}^{\mathrm{cyc}}$ above and the isomorphism $\underline{\GT}\cong \End_0(\widehat{\PaRB}^{\mathrm{cyc}})$ proved in~\cite[Theorem~5.4]{robertson2025grothendieck}.
\end{proof}

\begin{comment}
\begin{cor}
There exists an isomorphism of monoids \[\End_0(\PaRB^{cyc})\cong\End_0(\trun{0}\hat\bS)\cong \underline{\GT}.\]
\end{cor}

\begin{proof}
    This follows from Lemma~\ref{lemma: iso between S0 and cyclic PaRB}, Theorem~\ref{thm: maps_out_of_P_0}, Theorem~\ref{thm: PaRB and GT}, and Theorem~\ref{thm: presentation for maps out of bS} together with the fact that any endomorphism of $\trun{0}\bS$ fixing the objects will immediately satisfy relation \ref{sym_rel_morph}.
\end{proof}

\end{comment}

In \cite{Nakamura-Schneps}, Nakamura and Schneps introduced a subgroup $\NS\subseteq \GT$ by imposing two additional relations, discovered through comparisons of Galois representations in $\widehat{\Gamma}_{0,5}$ and $\widehat{\Gamma}_{1,2}$. These relations encode the genus-one compatibility needed to extend the genus-zero Grothendieck--Teichmüller action to the full tower of profinite mapping class groups.


\begin{definition}\label{def: NS}
The monoid $\underline{\NS}\subseteq \underline{\GT}$ consists of those elements $(\lambda,f)\in\underline{\GT}$ satisfying the following two additional relations:
\begin{equation}\label{eq: III'}\tag{NS I}
f(\beta_1\beta_3, \beta_2^2) (\beta_1\beta_2\beta_3)^{4\rho_{2}} = h((\beta_1\beta_2)^{3}, (\beta_2\beta_3)^3) f(\beta_1^2,\beta_2^2) f(\beta_3^2,(\beta_1\beta_2)^{3})
\quad \text{in } \widehat{\Gamma}_{0,5},
\end{equation}
and
\begin{equation}\label{eq: IV}\tag{NS II}
f(\beta_1,\beta_2^{4}) = \beta_{2}^{8\rho_{2}} f(\beta_1^2,\beta_2^2) \beta^{4}_{1} (\beta_1\beta_2)^{-6\rho_2} \quad \text{in } \widehat{\Br}_3/\mathbb{Z}\cong \widehat{\Gamma}_{1,1}.
\end{equation}

Here $h=h(x,y)$ is the unique element of $\widehat{\mathbb F}_2$ such that
\[
f(x,y)=h(y,x)^{-1}h(x,y)
\]
\cite{lochak1997cohomological}, and $\rho_2$ denotes the Kummer $1$-cocycle with respect to the roots of $2$. Equivalently,
\[
h(x,y)\equiv (xy)^{\rho_2}
\]
in the abelianisation $\widehat{\mathbb F}_2^{ab}\cong \widehat{\mathbb Z}\times\widehat{\mathbb Z}$ (see \cite[Definition~5.1]{Nakamura-Schneps}). The group $\NS$ is the group of invertible elements of of the monoid $\underline{\NS}$.
\end{definition}

In \eqref{eq: III'}, the elements $\beta_1,\beta_2,\beta_3$ are the standard generators of the braid group $\Br_4$, viewed as elements of the mapping class group $\Gamma_{0,5}$; see \cite[Proposition~5.4]{Nakamura-Schneps}. Nakamura and Schneps prove that $\NS$ is a subgroup of $\GT$ and that the absolute Galois group maps into $\NS$ \cite[Theorem~1.2]{Nakamura-Schneps}. More explicitly, for $g\in\galQ$, the corresponding element is $(\lambda_g,f_g)$, where $\lambda_g$ is the cyclotomic character and $f_g(x,y)$ lies in the commutator subgroup of $\pi_1^{\mathrm{et}}(\calM_{0,4})\cong\widehat{\mathbb F}_2$.

Nakamura and Schneps construct actions of $\NS$ on profinite mapping class groups using surfaces equipped with quilt, or seam, decompositions. For the purposes of the formulae below, a marking $m:G\hookrightarrow S$ provides the same data: it determines the standard curves and Dehn twist generators on which the Nakamura--Schneps action is defined. In \cite[Theorem~11.2]{Nakamura-Schneps}, they give explicit formulae for the action of $\NS$ on these Dehn twists. We translate those formulae into the language of marked surfaces and use them to construct an action on the profinite modular operad $\widehat{\bS}$. The resulting action extends the genus-zero $\underline{\GT}$-action identified in 
Proposition~\ref{prop:genus-zero-GT}.


\begin{theorem}\label{main theorem NS action}
There is an inclusion of monoids $\underline{\NS}\hookrightarrow \End_0(\widehat{\bS})$ whose restriction to genus zero agrees with the isomorphism \[\underline{\GT}\cong \End_0(\trun{0}\widehat{\bS})\] from Proposition~\ref{prop:genus-zero-GT}, after restricting along $\underline{\NS}\hookrightarrow\underline{\GT}$.
\end{theorem}

\begin{proof}
Let $F=(\lambda,f)\in\underline{\NS}$. Recall that an object-fixing map of modular operads $\iota_F:\bS\longrightarrow |\widehat{\bS}|$ extends uniquely to a continuous object-fixing endomorphism $\hat{\iota}_F:\widehat{\bS}\longrightarrow \widehat{\bS}.$ Thus it is enough, using Theorem~\ref{thm: presentation for maps out of bS}, to choose morphisms
\[
\hat{\iota}_F(\tau)=\mathbf{t}_F \in \Hom_{\widehat{\bS}(0,2)}(|,|),
\qquad
\hat{\iota}_F(\beta)\mathbf{b}_F \in \Hom_{\widehat{\bS}(0,3)}(P,(12)^*P),
\]
and
\[
\hat{\iota}_F(\alpha)=\mathbf{a}_F \in
\Hom_{\widehat{\bS}(0,4)}(P\circ_1 P,P\circ_2 P),
\qquad
\hat{\iota}_F(\sigma)=\mathbf{s}_F \in
\Hom_{\widehat{\bS}(1,1)}(\xi_{12}P,\xi_{12}P)
\]
and check that these choices satisfy the defining relations. Since $\iota_F$ is object-fixing, the object-level symmetry relation \eqref{sym_rel_morph} is automatic. Thus it remains to check the morphism relations \eqref{twist_rel_morphS}, \eqref{pent_rel_1_morphS}, \eqref{hex_rel_1_morphS}, \eqref{hex_rel_2_morphS}, \eqref{g1_rel_morph_1}, \eqref{g1_rel_morph_2}, and \eqref{g2_rel_morph_1}.

We chose the values of our generating morphisms using the action described in \cite{Nakamura-Schneps}. Explicitly, letting $\nu=(\lambda-1)/2$, we set
\begin{align}
    \mathbf{t}_F& = \tau^\lambda
    &
    \mathbf{b}_F &=\beta \cdot (\beta \cdot (12)^*\beta)^{\nu} \nonumber\\
    \mathbf{a}_F &= \alpha \cdot f(D_{d},D_{c})
    &
    \mathbf{s}_F &= \sigma^{\lambda}\cdot D_{b}^{8\rho_2(F)}\cdot f(D_{b}^2,D_a^2) \cdot D_a^{-8\rho_2(F)} .
    \label{NS Action}
\end{align}
Here $\rho_2(F)$ denotes the value, uniquely determined by $F=(\lambda,f)$, of the Nakamura--Schneps extension of the Kummer $1$-cocycle with respect to the roots of $2$. The symbol $D_i$ denotes the Dehn twist about the curve $i$. The curves $a,b,c,d$ appearing in \eqref{NS Action} are shown in Figures~\ref{fig:S-move} and \ref{fig:action-on-A-move}.
    
    \begin{figure}[ht]
        \centering
        \includegraphics[]{Figures/action-on-alpha.pdf}
        \caption{}
        \label{fig:action-on-A-move}
    \end{figure}

We now check the defining relations. The genus-zero relations can be dealt with first. By Proposition~8.1 of \cite{Boavida-Horel-Robertson} and Proposition~\ref{prop:genus-zero-GT}, the object-fixing endomorphisms of $\trun{0}\widehat{\bS}$ are identified with $\underline{\GT}$, and the element $F=(\lambda,f)\in\underline{\NS}\subseteq\underline{\GT}$ acts on the generators by the formulae defining $\mathbf{t}_F$, $\mathbf{b}_F$, and $\mathbf{a}_F$. Therefore these three morphisms define the genus-zero part of $\iota_F$, \[\trun{0}(\iota_F):\trun{0}\bS\longrightarrow |\trun{0}\widehat{\bS}|,\] and the genus-zero relations \eqref{twist_rel_morphS}, \eqref{pent_rel_1_morphS}, \eqref{hex_rel_1_morphS}, and \eqref{hex_rel_2_morphS} are satisfied.

It remains to check the genuinely genus-one relations \eqref{g1_rel_morph_1}, \eqref{g1_rel_morph_2}, and \eqref{g2_rel_morph_1}. In the following calculations we use the mapping class group relations from Proposition~\ref{prop: generators and relations of mcg}, together with the conjugation formulae in Lemma~\ref{lemma: conjugation by mapping classes}.

We first check the two relations in type $(1,1)$, namely \eqref{g1_rel_morph_1} and \eqref{g1_rel_morph_2}. These are relations in
\[
\Aut_{\widehat{\bS}(1,1)}(\xi_{12}P) \cong \widehat{\Gamma}_{1,1}.
\] We will use the curves shown in Figure~\ref{fig: g1-curves}. Since the calculations below are notation-heavy, we write $a$, $b$, and $c$ instead of $D_{a}, D_{b}$ and $D_{c}$ for the Dehn twists about the curves labelled $a$, $b$, and $c$.

\begin{figure}[h!]
    \centering
    \includegraphics[width=0.25\linewidth]{Figures/G1-curves.pdf}
    \caption{Curves used in the verification of the genus-one relations in $\widehat{\Gamma}_{1,1}$.}
    \label{fig: g1-curves}
\end{figure}

\emph{Checking relation \eqref{g1_rel_morph_1}:}  we first note that $\xi_{02}\beta^{-1}=(aba)^2$ and $\xi_{02}(\beta\cdot (12)^*\beta)=\partial_0^{-1}$. Then 
\begin{align*}
    \hat{\iota}_{F}(G1a)&=(aba)^{\lambda} b^{8\rho_2}f(b^2,a^2)a^{-8\rho_2} \cdot (aba)^{\lambda} b^{8\rho_2}f(b^2,a^2)a^{-8\rho_2} \cdot  (aba)^{-2} \partial_0^{-\nu}.
\end{align*}
Using that $(aba)^{-1}a(aba)=b$ and $(aba)^{-1}b(aba)=a$, we can move the middle $(aba)^\lambda$ term to the left to get:
\begin{align*}
    \hat{\iota}_{F}(G1a)&=(aba)^{2\lambda} a^{8\rho_2}f(a^2,b^2)b^{-8\rho_2} b^{8\rho_2}f(b^2,a^2)a^{-8\rho_2} (aba)^{-2} \partial_0^{-\nu}
\end{align*}
Recalling that $f(a^2,b^2)=f(b^2,a^2)^{-1}$, this gives us
\begin{align*}
    \hat{\iota}_{F}(G1a)&=(aba)^{4\nu} (aba)^2 (aba)^{-2} \partial_0^{-\nu}
\end{align*}
Since $(aba)^4=\partial_0$, relation \eqref{g1_rel_morph_1} follows.

\medskip 

\emph{Checking relation \eqref{g1_rel_morph_2}:} Recall that $d=\xi_{12}(\id_P\circ_1\tau)\in \Aut_{\bS(1,1)}(\xi_{12}P)$. Under the identification $\Aut_{\bS(1,1)}(\xi_{12}P)\cong \Gamma_{1,1}$, the morphism $d$ is the Dehn twist about the curve labelled $a$ in Figure~\ref{fig: g1-curves}. Hence 
    \[a=\xi_{1,2}(\id_P\circ_1 \tau)\]
and therefore its image under $\iota_F$ is represented in $\widehat{\Gamma}_{1,1}$ by $a^\lambda$. We compute the image of the corresponding loop, replacing each occurrence of $d$ by $a^\lambda$:
\begin{align*}
    \iota_F(G1b)
    &=(aba)^{\lambda} b^{8\rho_2}f(b^2,a^2)a^{-8\rho_2}\cdot a^{\lambda}\cdot  a^{8\rho_2}f(a^2,b^2)b^{-8\rho_2}(aba)^{-\lambda}\cdot a^{\lambda}\cdot  a^{8\rho_2}f(a^2,b^2)b^{-8\rho_2}(aba)^{-\lambda}\cdot a^{\lambda}.
\end{align*}


Using that $a^{-1}ba=c$ we get \[\iota_{F}(G1)=(aba)^{\lambda} b^{8\rho_2}f(b^2,a^2)a^{-8\rho_2}\cdot a^{\lambda}\cdot  a^{8\rho_2}f(a^2,b^2)b^{-8\rho_2}(aba)^{-\lambda}\cdot a^{\lambda}\cdot  a^{8\rho_2}f(a^2,b^2)b^{-8\rho_2}\mathbf{(aba)^{-1}a}(aca)^{-2\nu}\cdot a^{2\nu}\]

Since $(aba)^{-1}a= (ba)^{-1}$ and 
    \begin{align*}
        (ba)a(ba)^{-1}&= c & 
        (ba)b(ba)^{-1}&= a,
            \end{align*}
we get
\[\iota_{F}(G1b)=(aba)^{\lambda} b^{8\rho_2}f(b^2,a^2)a^{-8\rho_2}\cdot a^{\lambda}\cdot  a^{8\rho_2}f(a^2,b^2)b^{-8\rho_2}(aba)^{-\lambda}\cdot \mathbf{a(ba)^{-1}} \cdot c^{2\nu}\cdot  c^{8\rho_2}f(c^2,a^2)a^{-8\rho_2}(aca)^{-2\nu}\cdot a^{2\nu}\]

Since $a(ba)^{-1}=b$ and 
    \begin{align*}
        bab^{-1}&=c & 
        bbb^{-1}&=b,
    \end{align*}
we get
\[\iota_{F}(G1b)=(aba)^{\lambda} b^{8\rho_2}f(b^2,a^2)a^{-8\rho_2}\cdot a^{\lambda}\cdot  a^{8\rho_2}f(a^2,b^2)b^{-8\rho_2} \mathbf{(aba)^{-1}b^{-1}} \cdot (cbc)^{-2\nu}\cdot c^{2\nu}\cdot  c^{8\rho_2}f(c^2,a^2)a^{-8\rho_2}(aca)^{-2\nu}\cdot a^{2\nu}\]
Now we use that 
    \begin{align*}
        (b(aba))a((aba)^{-1}b^{-1})&=b &  (b(aba))b((aba)^{-1}b^{-1})&=c 
    \end{align*}
giving us
\[\iota_{F}(G1b)=(aba)^{\lambda} b^{8\rho_2}f(b^2,a^2)a^{-8\rho_2}\cdot \mathbf{a(aba)^{-1}b^{-1}} \cdot b^{2\nu}\cdot  b^{8\rho_2}f(b^2,c^2)c^{-8\rho_2}  (cbc)^{-2\nu}\cdot c^{2\nu}\cdot  c^{8\rho_2}f(c^2,a^2)a^{-8\rho_2}(aca)^{-2\nu}\cdot a^{2\nu}\]

Since $a(aba)^{-1}b^{-1}=(bab)^{-1}$ and 
    \begin{align*}
    (bab)a(bab)^{-1} &=b &  (bab)b(bab)^{-1}&=a
    \end{align*}
we get
\[\iota_{F}(G1b)=\mathbf{(aba)(aba)^{-1}} \cdot(bab)^{2\nu} a^{8\rho_2}f(a^2,b^2)b^{-8\rho_2}\cdot  b^{2\nu}\cdot  b^{8\rho_2}f(b^2,c^2)c^{-8\rho_2}  (cbc)^{-2\nu}\cdot c^{2\nu}\cdot  c^{8\rho_2}f(c^2,a^2)a^{-8\rho_2}(aca)^{-2\nu}\cdot a^{2\nu}\]

We now note that $(aba)^2$, $(aca)^2$, and $(bcb)^2$ are all equal and generate the centre of $\widehat{\Gamma}_{1,1}$, so we have
\[\iota_{F}(G1b)=(aba)^{-2\nu} a^{8\rho_2}f(a^2,b^2)\cdot  b^{2\nu}\cdot  f(b^2,c^2)\cdot c^{2\nu}\cdot  f(c^2,a^2)\cdot a^{2\nu}a^{-8\rho_2}\]
Finally, using that $c=a^{-1}ba$, we can check that $b^2a^2c^2=(aba)^2$ so using \cite[Relation (1.5.1)]{Nakamura-Schneps} (recall they write multiplication in the opposite order), we have
\[\iota_{F}(G1b)=(aba)^{-2\nu} a^{8\rho_2}\cdot(aba)^{2\nu}\cdot a^{-8\rho_2}=1.\]
Hence, relation \eqref{g1_rel_morph_2} holds.

\medskip

\emph{Checking relation \eqref{g2_rel_morph_1}:} We will study relation \eqref{g2_rel_morph_1} through the composition of morphisms depicted in Figure~\ref{fig:2-cell from NS} starting at the top of the diagram. In relation \eqref{g2_rel_morph_1}, the mapping class $\alpha$ appears conjugating certain mapping classes. However, by Lemma~\ref{lemma: conjugation by mapping classes}, these conjugations can be expressed entirely in terms of mapping classes on the surface $S_{\graphgtwo}$.\footnote{Here the graph $\graphgtwo$ is the composition $\xi_{2}^{3}(P\circ_1P)$ in $\Upsilon$.} Hence all our computations can be made on the profinite completion of the mapping class group of $S_{\graphgtwo}$.

For such computations, we will use as a reference the curves in Figure~\ref{fig:curves for g2} :

\begin{figure}[h!]
\centering
\includegraphics[width=0.85\linewidth]{Figures/curves-for-g2-2.pdf}
% \includegraphics[width=0.85\linewidth]{Figures/curves-for-g2.pdf}
\caption{}
\label{fig:curves for g2}
\end{figure} 

We start by noting the following equality on mapping class groups
    \[\xi_{12}[(z_3^*\beta)\circ_2(\tau^{-1},\tau^{-1})\circ \beta)]=\partial_0^{\frac{1}{2}}\partial_1^{\frac{1}{2}}\]
where $\partial_0^{\frac{1}{2}}\partial_1^{\frac{1}{2}}$ denotes the simultaneous half-twists along boundaries $\partial_0$ and $\partial_1$. Together with the definition of $\iota_{F}$ and relation \eqref{twist_rel_morphS}, this yields
    \[\hat{\iota}_{F}[\xi_{12}((z_3^*\mathbf{b})\circ_2(\mathbf{t}^{-1},\mathbf{t}^{-1})\circ \mathbf{b})]=\partial_0^{\nu}\partial_1^{\nu}(\partial_0^{\frac{1}{2}}\partial_1^{\frac{1}{2}}).\]
Since $\lambda=2\nu+1$, we abuse notation and denote this element by $\partial_0^{\frac{\lambda}{2}}\partial_1^{\frac{\lambda}{2}}$. Hence, checking relation \eqref{g2_rel_morph_1} amounts to showing that the following expression is trivial in the profinite mapping class group of $S_{\graphgtwo}$:
\begin{multline*}
    \iota_{F}(G2)=\partial_0^{-\frac{\lambda}{2}}\partial_1^{-\frac{\lambda}{2}}f(a_1,e_3)\cdot (a_3 a_2 a_3)^{\lambda} a_2^{8\rho_2} f(a_2^2,a_3^2) a_3^{-8\rho_2} \cdot f(e_3,a_1) \cdot a_3^{-\lambda} a_1^{\lambda} \cdot f(a_1,e_3) \cdot (a_3 a_2 a_3)^{\lambda} a_2^{8\rho_2} f(a_2^2,a_3^2) a_3^{-8\rho_2}\cdot f(e_3,a_1)
\end{multline*}

As in the proof of relation \eqref{g1_rel_morph_2}, we repeatedly use the fact that $\lambda=2\nu+1$ to separate one copy of each profinite power and move it to the left-hand side of the expression. Starting from the rightmost such term:
\begin{multline*}
    \iota_{F}(G2)=\partial_0^{-\frac{\lambda}{2}}\partial_1^{-\frac{\lambda}{2}}f(a_1,e_3)\cdot (a_3 a_2 a_3)^{\lambda} a_2^{8\rho_2} f(a_2^2,a_3^2) a_3^{-8\rho_2} \cdot f(e_3,a_1) \cdot a_3^{-\lambda} a_1^{\lambda} \cdot f(a_1,e_3) \cdot \\\mathbf{(a_3 a_2 a_3)}\cdot (a_3 a_2 a_3)^{2\nu} a_2^{8\rho_2} f(a_2^2,a_3^2) a_3^{-8\rho_2}\cdot f(e_3,a_1)
\end{multline*}

Now using that 
    \begin{align*}
        (a_3 a_2 a_3)^{-1}e_3(a_3 a_2 a_3)&= e_3 & (a_3 a_2 a_3)^{-1}a_1(a_3 a_2 a_3)&=e_2\\
        (a_3 a_2 a_3)^{-1}a_3(a_3 a_2 a_3)&= a_2
    \end{align*}
we get
\begin{multline*}
    \iota_{F}(G2)=\partial_0^{-\frac{\lambda}{2}}\partial_1^{-\frac{\lambda}{2}}f(a_1,e_3)\cdot (a_3 a_2 a_3)^{\lambda} a_2^{8\rho_2} f(a_2^2,a_3^2) a_3^{-8\rho_2} \cdot f(e_3,a_1) \cdot \mathbf{[a_3^{-1}a_1(a_3 a_2 a_3)]}\cdot a_2^{-2\nu} e_2^{2\nu} \cdot  \\ f(e_2,e_3) \cdot(a_3 a_2 a_3)^{2\nu} a_2^{8\rho_2} f(a_2^2,a_3^2) a_3^{-8\rho_2}\cdot f(e_3,a_1)
\end{multline*}

Since $a_1$ and $a_3$ are disjoint, they commute. Hence the expression in bold is equal to $a_1a_2a_3$. We now use the following relations in the mapping class group
    \begin{align*}
        (a_1a_2a_3)^{-1} a_1 (a_1a_2a_3)&=e_2 & (a_1a_2a_3)^{-1}e_3(a_1a_2a_3)&= e_1\\
        (a_1a_2a_3)^{-1} a_3 (a_1a_2a_3)&= a_2 & (a_1a_2a_3)^{-1} a_2 (a_1a_2a_3)&= a_1
    \end{align*}
then pulling the bold terms to the left, we get
\begin{multline*}
    \iota_{F}(G2)=\partial_0^{-\frac{\lambda}{2}}\partial_1^{-\frac{\lambda}{2}}f(a_1,e_3)\cdot \mathbf{[(a_3a_2a_3)(a_1 a_2 a_3)]}\cdot(a_2 a_1 a_2)^{2\nu} a_1^{8\rho_2} f(a_1^2,a_2^2) a_2^{-8\rho_2} \cdot f(e_1,e_2) \cdot  a_2^{-2\nu} e_2^{2\nu} \cdot  \\ f(e_2,e_3) \cdot(a_3 a_2 a_3)^{2\nu} a_2^{8\rho_2} f(a_2^2,a_3^2) a_3^{-8\rho_2}\cdot f(e_3,a_1)
\end{multline*}

Now note that
    \begin{align*}
        [(a_3a_2a_3)(a_1 a_2 a_3)]^{-1}e_3[(a_3a_2a_3)(a_1 a_2 a_3)]&=e_1 \\
        [(a_3a_2a_3)(a_1 a_2 a_3)]^{-1}a_1[(a_3a_2a_3)(a_1 a_2 a_3)]&=a_3
    \end{align*}
And we get
    \begin{multline*}
        \iota_{F}(G2)=\partial_0^{-\frac{\lambda}{2}}\partial_1^{-\frac{\lambda}{2}}\mathbf{[(a_3a_2a_3)(a_1 a_2 a_3)]}\cdot f(a_3,e_1)\cdot (a_2 a_1 a_2)^{2\nu} a_1^{8\rho_2} f(a_1^2,a_2^2) a_2^{-8\rho_2} \cdot f(e_1,e_2) \cdot  a_2^{-2\nu} e_2^{2\nu} \cdot  \\ f(e_2,e_3) \cdot(a_3 a_2 a_3)^{2\nu} a_2^{8\rho_2} f(a_2^2,a_3^2) a_3^{-8\rho_2}\cdot f(e_3,a_1)
    \end{multline*}
Since
    \begin{align*}
        (a_3a_2a_3)(a_1 a_2 a_3)&= \partial_0^{\frac{1}{2}}\partial_2^{\frac{1}{2}},
    \end{align*}
our expression for $\iota_{F}(G2)$ simplifies to 
    \begin{multline*}
        \iota_{F}(G2)=\partial_0^{-\nu}\partial_1^{-\nu}\cdot f(a_3,e_1)\cdot (a_2 a_1 a_2)^{2\nu} a_1^{8\rho_2} f(a_1^2,a_2^2) a_2^{-8\rho_2} \cdot f(e_1,e_2) \cdot  a_2^{-2\nu} e_2^{2\nu} \cdot  \\ f(e_2,e_3) \cdot(a_3 a_2 a_3)^{2\nu} a_2^{8\rho_2} f(a_2^2,a_3^2) a_3^{-8\rho_2}\cdot f(e_3,a_1)
    \end{multline*}
which is shown to be trivial in \cite[Proof of Claim 8.3, relation (6AS), end of page 542]{Nakamura-Schneps} (recall that in that paper they compose mapping classes in the opposite order).

\medskip 

\emph{Homomorphism:} We have now constructed a map $\NS\hookrightarrow\Aut_0(\hat{S})$ and we finish by showing it preserves the composition in $\NS$. Recall that, for \(F=(\lambda,f)\) and \(G=(\mu,g)\), our multiplication
convention is defined by
\[
F\cdot G
=
(\lambda\mu,\,
g(x,y)\,f(g(x,y)^{-1}x^\mu g(x,y),y^\mu)),
\]
so that \(\varphi_{F\cdot G}=\varphi_G\circ\varphi_F\).  It remains to verify that the composition $\iota_G\circ \iota_F$ agrees with $\iota_{F\cdot G}$ on the generators $\tau$, $\beta$, $\alpha$, and $\sigma$. The verification for $\tau$ and $\beta$ is immediate. The arguments for $\alpha$ and $\sigma$ are similar, so we only discuss the latter, which is slightly more involved.

We now show that
    \[\iota_G\circ \iota_F(\sigma)=\iota_{F\cdot G}(\sigma).\]
The left-hand side is equal to 
    \begin{align}\label{eq:homomorphism eq for sigma}
        \iota_G(\sigma^{\lambda}\cdot {b}^{8\rho_2(F)}\cdot f({b}^2,a^2) \cdot a^{-8\rho_2(F)})=\iota_G(\sigma)^{\lambda}\cdot {\iota_G(b)}^{8\rho_2(F)}\cdot f({\iota_G(b)}^2,\iota_G(a)^2) \cdot \iota_G(a)^{-8\rho_2(F)}
    \end{align}
The term $\iota_G(\sigma)$ is explicitly described in \eqref{NS Action} and, as discussed in the proof of \eqref{g2_rel_morph_1}, we have $\iota_G(a)=a^\mu$. Moreover, since $b=\sigma^{-1}a\sigma$, we get 
    \begin{align*}
        \iota_G(b)&=\iota_G(\sigma)^{-1}  \cdot \iota_G(a)\cdot \iota_G(\sigma)\\
        &=[\sigma^{\mu}\cdot b^{8\rho_2(G)}\cdot g(b^2,a^2) \cdot a^{-8\rho_2(G)}]^{-1} \cdot a^\mu \cdot [\sigma^{\mu}\cdot b^{8\rho_2(G)}\cdot g(b^2,a^2) \cdot a^{-8\rho_2(G)}]\\
        &=a^{8\rho_2(G)}\cdot  g(b^2,a^2)^{-1} \cdot b^{-8\rho_2(G)}\cdot \sigma^{-\mu} \cdot a^\mu \cdot \sigma^{\mu}\cdot b^{8\rho_2(G)}\cdot g(b^2,a^2) \cdot a^{-8\rho_2(G)}.
    \end{align*}
Since $\mu$ is of the form $2\nu'+1$, and that $\sigma^2$ commutes with $a$ and $b$, we get
    \begin{align*}
        \iota_G(b)
        &=a^{8\rho_2(G)}\cdot  g(b^2,a^2)^{-1} \cdot b^{-8\rho_2(G)}\cdot  b^\mu \cdot b^{8\rho_2(G)}\cdot g(b^2,a^2) \cdot a^{-8\rho_2(G)}\\
        &=a^{8\rho_2(G)}  g(b^2,a^2)^{-1}  b^\mu g(b^2,a^2)  a^{-8\rho_2(G)}.
    \end{align*}
Substituting these expressions into \eqref{eq:homomorphism eq for sigma}, we obtain
    \begin{multline}\label{eq: homomorphism on sigma-2}
        \iota_G\circ \iota_F(\sigma)=[\sigma^{\mu}\cdot b^{8\rho_2(G)}\cdot g(b^2,a^2) \cdot a^{-8\rho_2(G)}]^\lambda \cdot a^{8\rho_2(G)}  g(b^2,a^2)^{-1}  b^{8\rho_2(F)\mu} g(b^2,a^2)  a^{-8\rho_2(G)}\cdot\\
         \cdot f(a^{8\rho_2(G)}  g(b^2,a^2)^{-1}  b^{2\mu} g(b^2,a^2)  a^{-8\rho_2(G)},a^{2\mu})\cdot a^{-8\rho_2(F)\mu}
    \end{multline}
We begin by simplifying the first factor on the right-hand side using that:
    \begin{align*}
        [\sigma^{\mu}\cdot b^{8\rho_2(G)}\cdot g(b^2,a^2) \cdot a^{-8\rho_2(G)}]^2&=\sigma^{\mu}\cdot b^{8\rho_2(G)}\cdot g(b^2,a^2) \cdot a^{-8\rho_2(G)}\cdot\sigma^{\mu}\cdot b^{8\rho_2(G)}\cdot g(b^2,a^2) \cdot a^{-8\rho_2(G)}\\
        &=\sigma^{2\mu}\cdot a^{8\rho_2(G)}\cdot g(a^2,b^2) \cdot b^{-8\rho_2(G)}\cdot b^{8\rho_2(G)}\cdot g(b^2,a^2) \cdot a^{-8\rho_2(G)}\\
        &=\sigma^{2\mu}
    \end{align*}
Since $\lambda=2\nu+1$, equation~\eqref{eq: homomorphism on sigma-2} becomes
    \begin{multline*}
        \iota_G\circ \iota_F(\sigma)=[\sigma^{\lambda\mu}\cdot b^{8\rho_2(G)}\cdot g(b^2,a^2) \cdot a^{-8\rho_2(G)}] \cdot a^{8\rho_2(G)}  g(b^2,a^2)^{-1}  b^{8\rho_2(F)\mu} g(b^2,a^2)  a^{-8\rho_2(G)}\cdot\\
         \cdot f(a^{8\rho_2(G)}  g(b^2,a^2)^{-1}  b^{2\mu} g(b^2,a^2)  a^{-8\rho_2(G)},a^{2\mu})\cdot a^{-8\rho_2(F)\mu}.
    \end{multline*}
After cancellation, this simplifies to
    \begin{align*}
        \iota_G\circ \iota_F(\sigma)&=\sigma^{\lambda\mu}\cdot b^{8(\rho_2(G)+\rho_2(F)\mu)} \cdot g(b^2,a^2)\cdot f(g(b^2,a^2)^{-1}  b^{2\mu} g(b^2,a^2),a^{2\mu})\cdot a^{-8(\rho_2(G)+\rho_2(F)\mu)}.
    \end{align*}
Using that $\rho_2(F\cdot G)=\rho_2(G)+\rho_2(F)\mu$, as discussed in \cite[Corollary 5.2]{Nakamura-Schneps}, we see that this is precisely the formula defining $\iota_{F\cdot G}(\sigma)$, as required.  Finally, the homomorphism is injective because its restriction to the
genus-zero truncation agrees with the inclusion
\(\underline{\NS}\subseteq \underline{\GT}\) under the isomorphism
\[
\underline{\GT}\cong
\End_0(\operatorname{tr}_{\leq 0}\widehat{\mathbf S})
\]
of Proposition~\ref{prop:genus-zero-GT}.
\end{proof}

\begin{remark}
The curve $e_2$ in Figure~\ref{fig:curves for g2} differs slightly from the curve labelled $e_2$ in \cite{Nakamura-Schneps}. The relation used in their calculation, however, is
\[
(a_2a_3)^{-1}a_1(a_2a_3)=e_2,
\]
which is the relation satisfied by the curve $e_2$ drawn here. Thus our figure should be read as using the curve determined by this conjugation relation.

We also note that Nakamura--Schneps write $\mu$ for the quantity we denote by $\nu=(\lambda-1)/2$, since $\mu$ is already reserved in this paper for the generating object of $\PaRB$.
\end{remark}

Theorem~\ref{main theorem NS action}, together with the inclusion $\galQ\hookrightarrow \NS$ constructed by Nakamura--Schneps, immediately gives the corresponding Galois action on the profinite modular operad of
surfaces.

\begin{cor}
The absolute Galois group $\galQ$ acts on $\widehat{\bS}$.
\end{cor}

\begin{proof}
This follows from Theorem~\ref{main theorem NS action} and the inclusion
$\galQ\hookrightarrow \NS$ of Nakamura--Schneps.
\end{proof}
\section{Modular $\infty$-operads}\label{sec:modular infinity operads}

To take profinite completions of modular operads in spaces, we need a model for homotopy-coherent modular operads. The relevant indexing category is the modular dendroidal category $\bM$, whose objects are connected genus-graded labelled graphs with non-empty boundary. Strict modular operads can be regarded as certain presheaves on $\bM$: their value on a graph is determined by their values on the corollas at its vertices. More precisely, the modular dendroidal nerve of a strict modular operad satisfies
\[
X_G \cong \prod_{v\in V_G} X_{\corolla_v}.
\]
See, \cite[Theorem~3.6]{hry_modular_nerve}. A modular $\infty$-operad is obtained by replacing this strict product condition by a Segal condition.

We first recall the modular dendroidal category. Let $G=(G,\lambda,\ell,\epsilon)$ be a genus-graded labelled graph. For each vertex $v\in V_G$, let $H_v$ be a genus-graded labelled graph whose boundary is identified with the half-edges incident to $v$. We require $H_v$ to have total genus $\epsilon(v)$, equivalently
\[
\epsilon(v)=\beta_1(H_v)+\sum_{w\in V_{H_v}}\epsilon_{H_v}(w),
\] where $\epsilon_{H_v}$ denotes the genus function on $H_v$.  Then we can define a new graph $G\{H_v\}_{v\in V_G}$ by the coequaliser
\[
\begin{tikzcd}
\displaystyle\coprod_{e\in iE_G}\updownarrow
\arrow[r, shift left=1, "u"]
\arrow[r, shift right=1, "d"']
&
\displaystyle\coprod_{v\in V_G} H_v
\arrow[r, "\pi"]
&
G\{H_v\}_{v\in V_G}.
\end{tikzcd}
\] Here, the maps $u$ and $d$ record the two incidences of each internal edge of $G$. The coequaliser glues the corresponding boundary legs of the graphs $H_v$. This is the \emph{graph substitution} construction of \cite[Construction~1.18]{hry1}.


\begin{definition}\label{def: modular dendroidal category}
The \emph{modular dendroidal category} $\bM$ is the category whose objects are connected genus-graded labelled graphs with non-empty boundary. Its morphisms, called \emph{graphical maps}, are generated by isomorphisms and graph substitutions. More precisely, a graphical map records the data of replacing vertices of a graph by graphs with matching boundary and genus data. See \cite[Definition~1.31]{hry1} for the full definition.
\end{definition}

\begin{figure}[h!]
\[
\begin{tikzpicture}[x=0.75pt,y=0.75pt,yscale=-1,xscale=1]
%uncomment if require: \path (0,823); %set diagram left start at 0, and has height of 823

%Straight Lines [id:da6488749677933408] 
\draw    (114.84,123.97) -- (143.88,188.12) ;
%Straight Lines [id:da18732783760567262] 
\draw    (73.89,176.06) -- (107.15,133.2) ;
%Straight Lines [id:da7889407840448126] 
\draw    (51.03,124.07) -- (102.31,123.97) ;
%Straight Lines [id:da4431489597110816] 
\draw    (143.88,188.12) -- (115.59,245.03) ;
%Shape: Circle [id:dp31356981671318795] 
\draw  [color={rgb, 255:red, 4; green, 146; blue, 194 }  ,draw opacity=1 ][fill={rgb, 255:red, 255; green, 255; blue, 255 }  ,fill opacity=1 ][dash pattern={on 4.5pt off 4.5pt}] (102.31,123.97) .. controls (102.31,117.05) and (107.92,111.44) .. (114.84,111.44) .. controls (121.76,111.44) and (127.37,117.05) .. (127.37,123.97) .. controls (127.37,130.89) and (121.76,136.5) .. (114.84,136.5) .. controls (107.92,136.5) and (102.31,130.89) .. (102.31,123.97) -- cycle ;
%Curve Lines [id:da771171194316091] 
\draw    (122.04,114.52) .. controls (166.89,52.51) and (199.98,153.31) .. (126.66,129.14) ;
%Straight Lines [id:da8453143109032858] 
\draw    (170.22,268.12) -- (115.59,245.03) ;
%Shape: Ellipse [id:dp2086114067940088] 
\draw  [color={rgb, 255:red, 237; green, 130; blue, 14 }  ,draw opacity=1 ][fill={rgb, 255:red, 255; green, 255; blue, 255 }  ,fill opacity=1 ][dash pattern={on 4.5pt off 4.5pt}] (156,191.3) .. controls (154.25,197.99) and (147.41,202) .. (140.71,200.25) .. controls (134.01,198.5) and (130.01,191.65) .. (131.76,184.95) .. controls (133.51,178.26) and (140.36,174.25) .. (147.05,176) .. controls (153.75,177.75) and (157.76,184.6) .. (156,191.3) -- cycle ;
%Shape: Boxed Bezier Curve [id:dp583141332441497] 
\draw    (121.93,256.52) .. controls (159.11,323.41) and (53.78,310.69) .. (106.73,254.51) ;
%Shape: Circle [id:dp21641337625529522] 
\draw  [color={rgb, 255:red, 140; green, 219; blue, 54 }  ,draw opacity=1 ][fill={rgb, 255:red, 255; green, 255; blue, 255 }  ,fill opacity=1 ][dash pattern={on 4.5pt off 4.5pt}] (127.71,248.21) .. controls (125.96,254.9) and (119.12,258.91) .. (112.42,257.16) .. controls (105.72,255.41) and (101.72,248.56) .. (103.47,241.86) .. controls (105.22,235.17) and (112.07,231.16) .. (118.76,232.91) .. controls (125.46,234.66) and (129.47,241.51) .. (127.71,248.21) -- cycle ;
%Shape: Circle [id:dp3919998773078771] 
\draw  [color={rgb, 255:red, 4; green, 146; blue, 194 }  ,draw opacity=1 ][dash pattern={on 4.5pt off 4.5pt}] (42.04,70.45) .. controls (42.04,50.91) and (57.88,35.07) .. (77.42,35.07) .. controls (96.96,35.07) and (112.8,50.91) .. (112.8,70.45) .. controls (112.8,89.99) and (96.96,105.83) .. (77.42,105.83) .. controls (57.88,105.83) and (42.04,89.99) .. (42.04,70.45) -- cycle ;
%Shape: Circle [id:dp5469801557415326] 
\draw  [color={rgb, 255:red, 237; green, 130; blue, 14 }  ,draw opacity=1 ][dash pattern={on 4.5pt off 4.5pt}] (172.85,175.51) .. controls (172.85,160.37) and (185.12,148.1) .. (200.26,148.1) .. controls (215.4,148.1) and (227.68,160.37) .. (227.68,175.51) .. controls (227.68,190.65) and (215.4,202.93) .. (200.26,202.93) .. controls (185.12,202.93) and (172.85,190.65) .. (172.85,175.51) -- cycle ;
%Straight Lines [id:da10145939265222326] 
\draw    (453.36,104.93) -- (480.57,165.38) ;
%Straight Lines [id:da6840236886882134] 
\draw    (349.18,152.84) -- (379.77,113.42) ;
%Straight Lines [id:da3915242304438935] 
\draw    (328.16,105.03) -- (375.32,104.93) ;
%Straight Lines [id:da13158542747729218] 
\draw    (398.37,104.93) -- (441.84,104.93) ;
%Straight Lines [id:da9875493906550447] 
\draw    (421.52,51.15) -- (386.85,104.93) ;
%Straight Lines [id:da3512817237369321] 
\draw    (421.52,51.15) -- (453.36,104.93) ;
%Shape: Circle [id:dp11097495085655085] 
\draw  [fill={rgb, 255:red, 255; green, 255; blue, 255 }  ,fill opacity=1 ] (375.32,104.93) .. controls (375.32,98.57) and (380.48,93.41) .. (386.85,93.41) .. controls (393.21,93.41) and (398.37,98.57) .. (398.37,104.93) .. controls (398.37,111.3) and (393.21,116.45) .. (386.85,116.45) .. controls (380.48,116.45) and (375.32,111.3) .. (375.32,104.93) -- cycle ;
%Shape: Circle [id:dp6559922512513051] 
\draw  [fill={rgb, 255:red, 255; green, 255; blue, 255 }  ,fill opacity=1 ] (410,51.15) .. controls (410,44.79) and (415.15,39.63) .. (421.52,39.63) .. controls (427.88,39.63) and (433.04,44.79) .. (433.04,51.15) .. controls (433.04,57.52) and (427.88,62.68) .. (421.52,62.68) .. controls (415.15,62.68) and (410,57.52) .. (410,51.15) -- cycle ;
%Shape: Ellipse [id:dp28458230797732953] 
\draw  [fill={rgb, 255:red, 255; green, 255; blue, 255 }  ,fill opacity=1 ] (441.84,104.93) .. controls (441.84,98.57) and (447,93.41) .. (453.36,93.41) .. controls (459.73,93.41) and (464.89,98.57) .. (464.89,104.93) .. controls (464.89,111.3) and (459.73,116.45) .. (453.36,116.45) .. controls (447,116.45) and (441.84,111.3) .. (441.84,104.93) -- cycle ;
%Straight Lines [id:da7405793901835506] 
\draw    (91.99,70.14) -- (103.35,95.25) ;
%Straight Lines [id:da2704008409978078] 
\draw    (47.65,90.52) -- (60.67,73.75) ;
%Straight Lines [id:da23190885369214975] 
\draw    (42.04,70.45) -- (58.78,70.14) ;
%Straight Lines [id:da3111650015820755] 
\draw    (91.99,70.14) -- (112,78.43) ;
%Straight Lines [id:da6112752101701214] 
\draw    (91.99,70.14) -- (109.19,53.12) ;
%Straight Lines [id:da591579984214321] 
\draw    (68.59,70.14) -- (87.08,70.14) ;
%Straight Lines [id:da42403731742552686] 
\draw    (78.44,47.25) -- (63.68,70.14) ;
%Straight Lines [id:da8604597036244791] 
\draw    (78.44,47.25) -- (91.99,70.14) ;
%Shape: Ellipse [id:dp48890303362623866] 
\draw  [fill={rgb, 255:red, 255; green, 255; blue, 255 }  ,fill opacity=1 ] (58.78,70.14) .. controls (58.78,67.43) and (60.97,65.23) .. (63.68,65.23) .. controls (66.39,65.23) and (68.59,67.43) .. (68.59,70.14) .. controls (68.59,72.85) and (66.39,75.04) .. (63.68,75.04) .. controls (60.97,75.04) and (58.78,72.85) .. (58.78,70.14) -- cycle ;
%Shape: Ellipse [id:dp4835970024442289] 
\draw  [fill={rgb, 255:red, 255; green, 255; blue, 255 }  ,fill opacity=1 ] (73.53,47.25) .. controls (73.53,44.55) and (75.73,42.35) .. (78.44,42.35) .. controls (81.14,42.35) and (83.34,44.55) .. (83.34,47.25) .. controls (83.34,49.96) and (81.14,52.16) .. (78.44,52.16) .. controls (75.73,52.16) and (73.53,49.96) .. (73.53,47.25) -- cycle ;
%Shape: Ellipse [id:dp09294465671430874] 
\draw  [fill={rgb, 255:red, 255; green, 255; blue, 255 }  ,fill opacity=1 ] (87.08,70.14) .. controls (87.08,67.43) and (89.28,65.23) .. (91.99,65.23) .. controls (94.69,65.23) and (96.89,67.43) .. (96.89,70.14) .. controls (96.89,72.85) and (94.69,75.04) .. (91.99,75.04) .. controls (89.28,75.04) and (87.08,72.85) .. (87.08,70.14) -- cycle ;

%Straight Lines [id:da6207155341595856] 
\draw    (187.17,152.19) -- (198.74,177.74) ;
%Straight Lines [id:da49058651435496536] 
\draw    (198.74,177.74) -- (187.47,200.41) ;
%Curve Lines [id:da03965419119301694] 
\draw    (201.73,173.18) .. controls (219.59,148.48) and (232.77,188.63) .. (203.56,179) ;
%Shape: Circle [id:dp18101776392820634] 
\draw  [fill={rgb, 255:red, 255; green, 255; blue, 255 }  ,fill opacity=1 ] (193.75,177.74) .. controls (193.75,174.98) and (195.98,172.75) .. (198.74,172.75) .. controls (201.49,172.75) and (203.73,174.98) .. (203.73,177.74) .. controls (203.73,180.5) and (201.49,182.73) .. (198.74,182.73) .. controls (195.98,182.73) and (193.75,180.5) .. (193.75,177.74) -- cycle ;
%Straight Lines [id:da26185445967315535] 
\draw    (480.57,165.38) -- (441.22,250.03) ;
%Curve Lines [id:da7928497690636905] 
\draw    (487.6,154.65) .. controls (529.62,96.56) and (560.61,190.99) .. (491.93,168.35) ;
%Shape: Ellipse [id:dp5472298127523607] 
\draw  [fill={rgb, 255:red, 255; green, 255; blue, 255 }  ,fill opacity=1 ] (468.83,165.38) .. controls (468.83,158.89) and (474.08,153.64) .. (480.57,153.64) .. controls (487.05,153.64) and (492.31,158.89) .. (492.31,165.38) .. controls (492.31,171.86) and (487.05,177.12) .. (480.57,177.12) .. controls (474.08,177.12) and (468.83,171.86) .. (468.83,165.38) -- cycle ;
%Straight Lines [id:da9275037610543019] 
\draw    (491.23,271.46) -- (440.85,250.17) ;
%Straight Lines [id:da17246183653719316] 
\draw    (441.22,250.03) -- (390.84,228.74) ;
%Shape: Ellipse [id:dp6439493847369149] 
\draw  [fill={rgb, 255:red, 255; green, 255; blue, 255 }  ,fill opacity=1 ] (379.26,228.74) .. controls (379.26,222.34) and (384.45,217.16) .. (390.84,217.16) .. controls (397.24,217.16) and (402.43,222.34) .. (402.43,228.74) .. controls (402.43,235.14) and (397.24,240.33) .. (390.84,240.33) .. controls (384.45,240.33) and (379.26,235.14) .. (379.26,228.74) -- cycle ;
%Shape: Ellipse [id:dp15675917651922633] 
\draw  [fill={rgb, 255:red, 255; green, 255; blue, 255 }  ,fill opacity=1 ] (429.27,250.17) .. controls (429.27,243.78) and (434.45,238.59) .. (440.85,238.59) .. controls (447.25,238.59) and (452.44,243.78) .. (452.44,250.17) .. controls (452.44,256.57) and (447.25,261.76) .. (440.85,261.76) .. controls (434.45,261.76) and (429.27,256.57) .. (429.27,250.17) -- cycle ;
%Shape: Boxed Bezier Curve [id:dp7358717336524252] 
\draw    (447.15,260.8) .. controls (481.33,322.31) and (384.47,310.61) .. (433.17,258.95) ;
%Curve Lines [id:da34900819714117004] 
\draw    (459.98,95.54) .. controls (501.23,38.51) and (531.65,131.21) .. (464.23,108.98) ;
%Shape: Ellipse [id:dp8284371495354373] 
\draw  [color={rgb, 255:red, 4; green, 146; blue, 194 }  ,draw opacity=1 ][dash pattern={on 4.5pt off 4.5pt}] (362.3,89.16) .. controls (362.3,57.92) and (387.62,32.59) .. (418.87,32.59) .. controls (450.11,32.59) and (475.44,57.92) .. (475.44,89.16) .. controls (475.44,120.41) and (450.11,145.73) .. (418.87,145.73) .. controls (387.62,145.73) and (362.3,120.41) .. (362.3,89.16) -- cycle ;
%Shape: Ellipse [id:dp6470131354645342] 
\draw  [color={rgb, 255:red, 237; green, 130; blue, 14 }  ,draw opacity=1 ][dash pattern={on 4.5pt off 4.5pt}] (455.36,165.38) .. controls (455.36,151.45) and (466.64,140.17) .. (480.57,140.17) .. controls (494.49,140.17) and (505.78,151.45) .. (505.78,165.38) .. controls (505.78,179.3) and (494.49,190.59) .. (480.57,190.59) .. controls (466.64,190.59) and (455.36,179.3) .. (455.36,165.38) -- cycle ;
%Shape: Circle [id:dp2882731415599845] 
\draw  [color={rgb, 255:red, 140; green, 219; blue, 54 }  ,draw opacity=1 ][dash pattern={on 4.5pt off 4.5pt}] (368.91,232.75) .. controls (368.91,204.39) and (391.9,181.39) .. (420.27,181.39) .. controls (448.64,181.39) and (471.63,204.39) .. (471.63,232.75) .. controls (471.63,261.12) and (448.64,284.12) .. (420.27,284.12) .. controls (391.9,284.12) and (368.91,261.12) .. (368.91,232.75) -- cycle ;
%Straight Lines [id:da6792600183091927] 
\draw    (79.07,214.04) -- (69.3,235.44) ;
%Straight Lines [id:da7033421088998842] 
\draw    (79.72,240.37) -- (69.13,235.51) ;
%Straight Lines [id:da9885858772107592] 
\draw    (69.3,235.44) -- (46.16,225.67) ;
%Shape: Ellipse [id:dp2330553093227623] 
\draw  [fill={rgb, 255:red, 255; green, 255; blue, 255 }  ,fill opacity=1 ] (40.84,225.67) .. controls (40.84,222.73) and (43.22,220.35) .. (46.16,220.35) .. controls (49.1,220.35) and (51.48,222.73) .. (51.48,225.67) .. controls (51.48,228.61) and (49.1,230.99) .. (46.16,230.99) .. controls (43.22,230.99) and (40.84,228.61) .. (40.84,225.67) -- cycle ;
%Shape: Ellipse [id:dp8591699508470692] 
\draw  [color={rgb, 255:red, 140; green, 219; blue, 54 }  ,draw opacity=1 ][dash pattern={on 4.5pt off 4.5pt}] (36.09,227.51) .. controls (36.09,214.48) and (46.65,203.92) .. (59.67,203.92) .. controls (72.7,203.92) and (83.26,214.48) .. (83.26,227.51) .. controls (83.26,240.54) and (72.7,251.1) .. (59.67,251.1) .. controls (46.65,251.1) and (36.09,240.54) .. (36.09,227.51) -- cycle ;
%Straight Lines [id:da4787390203115717] 
\draw    (69.13,235.51) -- (55.35,250.77) ;
%Straight Lines [id:da5055444695966994] 
\draw    (69.3,235.44) -- (74.85,245.89) ;
%Shape: Ellipse [id:dp14002901248011645] 
\draw  [fill={rgb, 255:red, 255; green, 255; blue, 255 }  ,fill opacity=1 ] (63.81,235.51) .. controls (63.81,232.57) and (66.19,230.19) .. (69.13,230.19) .. controls (72.07,230.19) and (74.45,232.57) .. (74.45,235.51) .. controls (74.45,238.45) and (72.07,240.83) .. (69.13,240.83) .. controls (66.19,240.83) and (63.81,238.45) .. (63.81,235.51) -- cycle ;

%Straight Lines [id:da32349549764022567] 
\draw    (259.5,179.74) -- (318.5,179.74) ;
\draw [shift={(320.5,179.74)}, rotate = 180] [color={rgb, 255:red, 0; green, 0; blue, 0 }  ][line width=0.75]    (10.93,-3.29) .. controls (6.95,-1.4) and (3.31,-0.3) .. (0,0) .. controls (3.31,0.3) and (6.95,1.4) .. (10.93,3.29)   ;

% Text Node
\draw (109.45,118.88) node [anchor=north west][inner sep=0.75pt]   [align=left] {$\displaystyle v$};
% Text Node
\draw (137.57,182.75) node [anchor=north west][inner sep=0.75pt]   [align=left] {$\displaystyle w$};
% Text Node
\draw (110.1,239.69) node [anchor=north west][inner sep=0.75pt]   [align=left] {$\displaystyle u$};
% Text Node
\draw (104.99,24.86) node [anchor=north west][inner sep=0.75pt]   [align=left] {$\displaystyle H_{v}$};
% Text Node
\draw (219.29,134.89) node [anchor=north west][inner sep=0.75pt]   [align=left] {$\displaystyle H_{w}$};
% Text Node
\draw (22,184.83) node [anchor=north west][inner sep=0.75pt]   [align=left] {$\displaystyle H_{u}$};

\end{tikzpicture}
\]
\caption{An example of a graphical map}\label{graph map}
\end{figure} 


The category of \emph{modular dendroidal objects} in $\bE$, denoted $\dMod(\bE)$, is the category of contravariant functors $X:\bM^{\mathrm{op}}\longrightarrow \bE$ and natural transformations. We write $X_G$ for the value of $X\in\dMod(\bE)$ on a graph $G\in\bM$, and we write
\[
\bM[G]:=\Hom_{\bM}(-,G)
\]
for the representable functor at $G$.

\begin{definition}\label{defn: Segal core}\cite[Section~3.1]{hry1}
Let $G$ be a graph with at least one vertex. The \emph{Segal core} of $G$ is the coequaliser in $\dMod(\Set)$
\[ \begin{tikzcd}
\displaystyle\coprod_{e\in iE_G}\bM[\updownarrow] \arrow[r, shift left=1] \arrow[r, shift right=1] & \displaystyle\coprod_{v\in V_G}\bM[\corolla_v] \arrow[r] & \Sc[G].
\end{tikzcd} \] The two arrows record the two incidences of each internal edge with the corollas at its endpoints. The Segal core is a subobject of the representable $\bM[G]$, with inclusion $\Sc[G]\longrightarrow \bM[G]$ induced by the inclusions of the local vertex data $\bM[\corolla_v]\to \bM[G]$.
\end{definition}

The Segal core detects whether the value of a modular dendroidal object on a graph is determined, up to homotopy, by its values on the corollas at the vertices. When $\bE$ is a Cartesian monoidal model category, the category $\dMod(\bE)$ admits the projective model structure, in which a map $X\to Y$ is a weak equivalence, respectively a fibration, if and only if $X_G\longrightarrow Y_G$ is a weak equivalence, respectively a fibration, in $\bE$ for every $G\in\bM$; see \cite[Theorem~11.6.1]{hirsch}.

A modular dendroidal object $X$ is called \emph{Segal} if, for every $G\in\bM$, the map induced by the Segal core inclusion
\[
\mathbb{R}\Map(\bM[G],X) \longrightarrow \mathbb{R}\Map(\Sc[G],X)
\]
is a weak equivalence of spaces.\footnote{Here $\bM[G]$ and $\Sc[G]$ are regarded as discrete objects of $\dMod(\bE)$ via the functor $\Set\to\bE$.} A modular dendroidal object $X$ is called \emph{reduced} if $X_{\updownarrow}=*$.
\begin{definition}\label{def: infty modular}
A \emph{modular $\infty$-operad} in $\bE$ is a modular dendroidal object
$X:\bM^{\mathrm{op}}\to\bE$ satisfying:
\begin{itemize}
    \item $X_{\updownarrow}=*$;
    \item for every graph $G\in\bM$, the Segal map
    \[
    \mathbb{R}\Map(\bM[G],X)
    \longrightarrow
    \mathbb{R}\Map(\Sc[G],X)
    \]
    is a weak equivalence of spaces.
\end{itemize}
\end{definition}

The condition $X_{\updownarrow}=*$ means that we are considering one-coloured modular $\infty$-operads. It also simplifies the Segal condition. Indeed, for a reduced object $X$, the edge-compatibility conditions in the Segal core are trivial, since each internal edge contributes a copy of $X_{\updownarrow}=*$. Thus, after replacing $X$ by an entrywise fibrant object if necessary, the Segal map may be identified with the comparison
\[
X_G \longrightarrow \prod_{v\in V_G} (X_{\corolla_v})_f .
\]
Here $(X_{\corolla_v})_f$ denotes a fibrant replacement in $\bE$, so that the product computes the corresponding homotopy product. Equivalently, a reduced
modular dendroidal object is a modular $\infty$-operad precisely when, for every graph $G$, the value $X_G$ is weakly equivalent to the homotopy product
of its values on the corollas at the vertices of $G$. 

In the cases used in this paper we do not need to worry much about fibrant replacements. If $\bE=\Grpd$, every object is fibrant, so the homotopy product is computed by the ordinary product. If $\bE=\sSet$ and $X$ is entrywise fibrant, for example if $X$ is fibrant in the localised Reedy model structure, the ordinary product again computes the homotopy product. In this situation the Segal condition can be written simply as requiring
\[
X_G \longrightarrow \prod_{v\in V_G} X_{\corolla_v}
\]
to be a weak equivalence.

\subsection{The modular dendroidal nerve}
Let $\Mod_{\infty}(\bE)$ denote the $\infty$-category of modular $\infty$-operads, as in Definition~\ref{def: infinity cat of infinity modular operads}.
A strict modular operad assigns operations to genus-graded corollas. Viewing a general genus-graded graph as a pattern for composing such operations, one obtains a modular dendroidal object by decorating each vertex with an operation of the corresponding genus and arity. This construction defines the modular dendroidal \emph{nerve} functor
\[
\nerve:\Mod(\bE)\longrightarrow \Mod_{\infty}(\bE).
\]

\begin{definition}\label{def: modular nerve}
The \emph{nerve} of a modular operad $\calP$ is the contravariant functor $\nerve\calP:\bM^{\mathrm{op}}\to\bE$ defined on a graph $G=(G,\lambda,\ell,\epsilon)$ by
\[
(\nerve\calP)_G := \prod_{v\in V_G} \calP\bigl(\epsilon(v),|\nb(v)|\bigr).
\]
In particular, if $C_v$ is the corolla at a vertex $v$, then $(\nerve\calP)_{C_v}=\calP\bigl(\epsilon(v),|\nb(v)|\bigr).$
\end{definition}

For a strict modular operad $\calP$, the Segal map is an isomorphism rather than just a weak equivalence. Indeed, for every graph $G$, the definition of the nerve identifies \[ (\nerve\calP)_G \cong \prod_{v\in V_G}(\nerve\calP)_{C_v}. \] Thus the nerve of a strict modular operad satisfies the Segal condition strictly. In particular, when ordinary finite products compute finite homotopy products in $\bE$, the object $\nerve\calP$ is a modular $\infty$-operad. This applies in the cases used in this paper, for example when $\bE=\Grpd$, and for entrywise fibrant objects in $\sSet$.


The nerve is fully faithful and admits a left adjoint $\tau:\dMod(\bE)\longrightarrow \Mod(\bE)$ such that $\tau\nerve\cong \id_{\Mod(\bE)}$ (\cite[Theorem~3.6]{hry_modular_nerve}). In Appendix~\ref{sec: nerve is homotopically fully faithful}, we prove a homotopical analogue of this statement. Namely, if $\bE=\sSet$ or $\Grpd$, then the modular dendroidal nerve is homotopically fully faithful. Thus, for any $\calP,\calQ\in\Mod(\bE)$, the induced map
\[ \mathbb{R}\Map_{\Mod(\bE)}(\calP,\calQ) \longrightarrow \mathbb{R}\Map_{\Mod_{\infty}(\bE)}(\nerve\calP,\nerve\calQ) \] is a weak equivalence of spaces.
\section{From profinite groupoids to profinite spaces}\label{sec: homotopy action} 
In previous sections, we showed that the modular operad in profinite groupoids $\widehat{\bS}$ admits an action of the monoid ${\NS}$ (Theorem~\ref{main theorem NS action}). The goal of this section is to show that $\widehat{\cs\nerve\bS}$ forms a modular $\infty$-operad in profinite spaces admitting an induced action of ${\NS}$. This modular $\infty$-operad may be regarded both as the profinite completion of the modular operad $\cs\bS$ of moduli spaces of surfaces and as a model for a Teichmüller tower built from étale homotopy types.

In this setting, the strict action of ${\NS}$ on $\widehat{\bS}$ induces an action in the homotopy category of modular $\infty$-operads. We therefore begin by explaining the passage from strict endomorphisms to homotopy endomorphisms.


\subsection{Profinite completion of modular $\infty$-operads}\label{sec: profinite modular infty operads}
We first explain how profinite completion is applied to modular dendroidal objects. Let $\bE$ be either $\sSet$ or $\Grpd$, and let $\mathscr C\subseteq \bE$ denote the full subcategory of $\pi$-finite spaces or finite groupoids, respectively.

We use the profinite completion of diagram categories developed in \cite[Theorem~2.4]{BM20}. For a suitably filtered small category $\mathscr K$, the pro-$\mathscr C$ completion of the category of coskeletal diagrams in $\operatorname{Fun}(\mathscr K^{\mathrm{op}},\bE)$ can be identified with the category of $\Pro(\mathscr C)$-valued diagrams $\operatorname{Fun}(\mathscr K^{\mathrm{op}},\Pro(\mathscr C))$. In this sense, profinite completion of coskeletal diagrams is computed entrywise.

We apply this result with $\mathscr K=\bM$. The graphical category $\bM$ is a generalised Reedy category \cite[Theorem~4.10 and Remark~4.17]{hry1}, and diagrams $X\in\operatorname{Fun}(\bM^{\mathrm{op}},\bE)$ are coskeletal in the relevant sense; see \cite{riehl2017inductive} and \cite[Section~6]{bm_reedy}. Hence the construction applies to modular dendroidal objects.

\begin{definition}\label{def: profinite modular operad}
Let $X:\bM^{\mathrm{op}}\to\bE$ be a modular dendroidal object. The \emph{profinite completion} of $X$ is the modular dendroidal object
\[
\widehat X:\bM^{\mathrm{op}}\longrightarrow \Pro(\mathscr C)
\]
obtained by applying pro-$\mathscr C$ completion entrywise.
\end{definition}

Since profinite completion is a left adjoint, it does not in general preserve finite limits. Thus, even if a modular dendroidal object $X$ satisfies the Segal condition, its entrywise profinite completion
\[
\widehat X:\bM^{\op}\longrightarrow \Pro(\mathscr C)
\]
need not satisfy the Segal condition. In general, profinite completion produces a modular dendroidal object in profinite spaces or profinite groupoids, rather
than a profinite modular $\infty$-operad.

There is, however, a useful exception in the groupoid case. As discussed in Section~\ref{sec: profinite completion of modular operads in groupoids}, and in particular in Proposition~\ref{prop: profinite completion preserves products of groupoids}, profinite completion is compatible with finite products under suitable finiteness hypotheses. The following lemma applies this product comparison to the Segal condition for modular $\infty$-operads in groupoids. %It gives the form of profinite completion that we will use to pass from the strict action of $\NS$ on the profinite modular operad $\widehat{\bS}$ to a homotopy action.

Recall that we write $C_{(g,n+1)}$ for the corolla of genus $g$ with $n+1$ legs. Thus, if $v$ is a vertex of a graph $G$, the corolla at $v$ is $C_{(\epsilon(v),|\nb(v)|)}$, with the labelling and genus grading inherited from $G$.

\begin{lemma}\label{lemma: profinite-completion-preserves-segal-groupoids}
Let $X:\bM^{\op}\to\Grpd$ be a modular $\infty$-operad such that each corolla value
$X_{C_{(g,n+1)}}$ has finitely many objects. Then entrywise profinite completion defines a modular $\infty$-operad
\[
\widehat X:\bM^{\op}\longrightarrow \widehat{\Grpd}.
\]
\end{lemma}

\begin{proof}
Since profinite completion sends the trivial groupoid to the trivial profinite groupoid, we have $\widehat X_{\updownarrow}=*$.

It remains to check the Segal condition. For every graph $G$, the Segal map for $X$ is an equivalence of groupoids
\[
X_G \longrightarrow \prod_{v\in V_G} X_{C_{(\epsilon(v),|\nb(v)|)}} .
\]
Applying profinite completion gives an equivalence in $\widehat{\Grpd}$
\[
\widehat X_G \longrightarrow
\left(\prod_{v\in V_G} X_{C_{(\epsilon(v),|\nb(v)|)}}\right)^{\wedge}.
\]
The graph $G$ has finitely many vertices, and each corolla value
$X_{C_{(\epsilon(v),|\nb(v)|)}}$ has finitely many objects. Hence Proposition~\ref{prop: profinite completion preserves products of groupoids} identifies the target naturally with $\prod_{v\in V_G}\widehat X_{C_{(\epsilon(v),|\nb(v)|)}}.$
Thus the completed Segal map
\[
\widehat X_G \longrightarrow
\prod_{v\in V_G}\widehat X_{C_{(\epsilon(v),|\nb(v)|)}}
\]
is a weak equivalence in $\widehat{\Grpd}$. Therefore $\widehat X$ satisfies the Segal condition, and hence is a modular $\infty$-operad in profinite groupoids.
\end{proof}

\subsection{From endomorphisms to homotopy endomorphisms}\label{sec: homotopy endomorphisms}

We now compare strict endomorphisms of the profinite modular operad $\widehat{\bS}$ with homotopy endomorphisms of the completed modular dendroidal nerve. The first point is that, under the relevant finiteness hypotheses, profinite completion commutes with the modular dendroidal nerve.

\begin{lemma}\label{lemma: N commutes with completion for groupoids}
Let $\calP$ be a modular operad in groupoids such that $\calP(g,n+1)$ has finitely many objects for all $g\geq 0$ and $n\geq 0$. Then there is a natural isomorphism of modular $\infty$-operads
\[
\widehat{\nerve\calP}\cong \nerve\widehat{\calP}.
\]
\end{lemma}

\begin{proof}
Let $\Mod(\Grpd)_{\mathrm{fin}}$ denote the full subcategory of $\Mod(\Grpd)$ spanned by modular operads $\calP$ such that each groupoid $\calP(g,n+1)$ has finitely many objects. We show that the diagram
\[
\begin{tikzcd}
\Mod(\Grpd)_{\mathrm{fin}} \arrow[d, "\widehat{(-)}"'] \arrow[r, "\nerve"] 
& \Mod_{\infty}(\Grpd) \arrow[d, "\widehat{(-)}"] \\
\Mod(\widehat{\Grpd}) \arrow[r, "\nerve"] 
& \Mod_{\infty}(\widehat{\Grpd})
\end{tikzcd}
\]
commutes up to natural isomorphism.

Going around the diagram counterclockwise, let $\widehat{\calP}$ be the modular operad in $\widehat{\Grpd}$ obtained from $\calP$ by entrywise profinite completion, as in Definition~\ref{def: profinite completion of modular operads in groupoids}. Its modular dendroidal nerve is the functor
\[
\nerve\widehat{\calP}:\bM^{\op}\longrightarrow \widehat{\Grpd}
\]
defined, for each graph $G\in\bM$, by
\[
(\nerve\widehat{\calP})_G
=
\prod_{v\in V_G}
\widehat{\calP}\bigl(\epsilon(v),|\nb(v)|\bigr).
\]

Going around the diagram in the other direction, profinite completion of the modular dendroidal object
$\nerve\calP:\bM^{\op}\to\Grpd$ is defined entrywise. Hence, for each graph $G\in\bM$,
\[
(\widehat{\nerve\calP})_G
=
\left(
\prod_{v\in V_G}
\calP\bigl(\epsilon(v),|\nb(v)|\bigr)
\right)^{\wedge},
\]
where $(-)^{\wedge}$ denotes profinite completion. Since $G$ has finitely many vertices and each groupoid $\calP(g,n+1)$ has finitely many objects, Proposition~\ref{prop: profinite completion preserves products of groupoids} identifies this groupoid naturally with
\[
\prod_{v\in V_G}
\widehat{\calP}\bigl(\epsilon(v),|\nb(v)|\bigr)
=
(\nerve\widehat{\calP})_G.
\]
These isomorphisms are natural in $G$, and therefore assemble to an isomorphism of modular dendroidal objects
\[
\widehat{\nerve\calP}\cong \nerve\widehat{\calP}.
\]
By Lemma~\ref{lemma: profinite-completion-preserves-segal-groupoids}, the left-hand side is a modular $\infty$-operad in profinite groupoids, and the right-hand side is the nerve of a modular operad in profinite groupoids. Hence this is an isomorphism of modular $\infty$-operads.
\end{proof}


We will use the following elementary path-object notion of homotopy for modular operads in groupoids and profinite groupoids. Let $\mathscr I$ denote the finite groupoid with objects $\{0,1\}$ and a unique morphism between any two objects. If $\calQ$ is a modular operad in groupoids, or in profinite groupoids, we write $\calQ^{\mathscr I}$ for the modular operad defined entrywise by
\[
(\calQ^{\mathscr I})(g,n+1)=\calQ(g,n+1)^{\mathscr I}.
\]
The two objects of $\mathscr I$ induce evaluation maps $d_0,d_1\colon \calQ^{\mathscr I}\to\calQ$. Given morphisms of modular operads $f,g\colon \calP\to\calQ$, a \emph{homotopy} from $f$ to $g$ is a morphism of modular operads $H\colon\calP\to\calQ^{\mathscr I}$ such that $d_0H=f$ and $d_1H=g$.

Because weak equivalences and fibrations of groupoids and profinite groupoids are defined entrywise, $\calQ^{\mathscr I}$ is an entrywise path object for $\calQ$. Thus this agrees with the usual model-categorical right homotopy relation in the corresponding category of modular operads. The relation is compatible with composition. We write $\HoAut(\calP)$ for the group of automorphisms of $\calP$ modulo this homotopy relation, and we write $\Aut_0(\calP)$ for the group of object-fixing automorphisms.



\begin{prop}\label{prop:injection}
    The map 
        \[\Aut_0(\widehat{\bS})\to \HoAut(\widehat{\bS})\]
    is injective.
\end{prop}

Before we prove Proposition~\ref{prop:injection} we require the following Lemma about $\widehat\Gamma_{1,1}\cong \widehat{\Gamma(S_{\contraction})}$. We use as a reference Figure~\ref{fig:S-move}.

\begin{lemma}\label{lemma:centralisers in the profinite completion of Gamma11}
    Let $a$ denote the curve in $S_{\contraction}$ depicted in Figure~\ref{fig:S-move}. Then we have the following relation between centralisers 
        \[C_{\widehat{\Gamma}_{1,1}}(D_a)= \overline{C_{\Gamma_{1,1}}(D_a)}=\overline{\langle D_\partial, D_a, (D_aD_bD_a)^2\rangle}.\]
\end{lemma}

\begin{proof}
    The capping of the boundary by gluing a punctured disc induces a short exact sequence
        \[1\to \langle D_\partial \rangle\to \Gamma_{1,1}\xrightarrow{\pi} \Gamma_1^1\to1\]
    where $\Gamma_1^1$ denotes the mapping class group of the one-punctured torus.

    As profinite completion is right-exact, this gives an exact sequence
        \[
            \widehat{\langle D_\partial \rangle}\to \widehat\Gamma_{1,1}\xrightarrow{\pi} \widehat\Gamma_1^1\to1.
        \]

    Restricting to the normaliser $N_{\widehat\Gamma_{1,1}}(\overline{\langle D_a\rangle})$ of $\overline{\langle D_a\rangle}$ in $\widehat\Gamma_{1,1}$ gives us the exact sequence
        \begin{align}\label{eq:exact-sequence-for-injectivity}
        \widehat{\langle D_\partial \rangle}\to N_{\widehat\Gamma_{1,1}}(\overline{\langle D_a\rangle})\xrightarrow{\pi} \pi(N_{\widehat\Gamma_{1,1}}(\overline{\langle D_a\rangle}))\to1.
        \end{align}

    It is clear that    
        \begin{align}\label{eq:injectivity-proof-1}
        \pi(N_{\widehat\Gamma_{1,1}}(\overline{\langle D_a\rangle}))\subset N_{\widehat\Gamma_{1}^1}(\overline{\langle \pi(D_a)\rangle})
        \end{align}
    but we will show this is in fact an equality. As $\Gamma_1^1\cong SL(2,\Z)$ is virtually free, by \cite[Corollary~2.9]{Ribes-Pavel} we have that
        \[N_{\widehat\Gamma_{1}^1}(\overline{\langle \pi(D_a)\rangle})=\overline{N_{\Gamma_{1}^1}(\langle \pi(D_a)\rangle)}.\]
    We now describe the normaliser of $\langle\pi(D_a)\rangle$ in $\Gamma_1^1$. Note that $\pi(D_a)$ is simply the Dehn twist in the once-punctured torus along the curve corresponding to $a$, so we will abuse notation and denote it also by $D_a$.
    
    The normaliser of $\langle D_a\rangle$ in $\Gamma_1^1$ consists of those mapping classes $f$ such that $f\cdot D_a\cdot f^{-1}$ is a power of $D_a$. However, by Lemma~\ref{lemma: conjugation by mapping classes}, such a conjugation is simply the Dehn twist $D_{f(a)}$. Since a power $D_a^k$ is only a Dehn twist if $k=1$, we conclude that any element in the normaliser of $D_a$ \emph{fixes} $a$, and hence $D_a$. In particular, this means that this normaliser is simply the centraliser of $D_a$.
        \[N_{\Gamma_{1}^1}(\langle \pi(D_a)\rangle)=C_{\Gamma_{1}^1}(D_a)\]
    As discussed above, this centraliser consists precisely of those mapping classes $f$ that fix the curve $a$. This can be seen to be precisely the subgroup generated by $D_a$ and $(D_a D_b D_a)^2$, for instance by thinking of $f$ as a mapping class on the surface obtained by cutting along $a$ which is allowed to permute the two new boundary components corresponding to $a$.

    Putting this into \eqref{eq:injectivity-proof-1}, this gives us 
        \[\pi(N_{\widehat\Gamma_{1,1}}(\overline{\langle D_a\rangle}))\subset \overline{\langle D_a, (D_aD_bD_a)^2\rangle}.\]
    Moreover, both $D_a$ and $(D_aD_bD_a)^2$ in $\Gamma_1^1$ are images of the corresponding elements of $\Gamma_{1,1}$, and these elements clearly lie in
    $N_{\widehat\Gamma_{1,1}}(\overline{\langle D_a\rangle})$.
    Hence
    \[\langle D_a,(D_aD_bD_a)^2\rangle
    \subset
    \pi(N_{\widehat\Gamma_{1,1}}(\overline{\langle D_a\rangle})).\]
    Since normalisers of closed subgroups are closed in profinite groups, the normaliser
    $N_{\widehat\Gamma_{1,1}}(\overline{\langle D_a\rangle})$ is closed, hence compact. Therefore its image under $\pi$  is closed and we obtain
    \[\overline{\langle D_a,(D_aD_bD_a)^2\rangle}
    \subset
    \pi(N_{\widehat\Gamma_{1,1}}(\overline{\langle D_a\rangle})). \]
    Together with the previous inclusion, this proves equality.
    
    % Moreover, we have that both $D_a$ and $(D_aD_bD_a)^2$ in $\Gamma_1^1$ are the image of the corresponding elements in $\Gamma_{1,1}$ via the homomorphism $\pi$. Moreover, since $\Gamma_1^1$ is virtually free, and hence subgroup separable, we get that the subgroup generated by $D_a$ and $(D_aD_bD_a)^2$ is closed in $\Gamma_1^1$. Giving us that the above inclusion is in fact an equality.

    Putting this together with  \eqref{eq:exact-sequence-for-injectivity} gives us a right exact sequence
        \begin{align*}
            \widehat{\langle D_\partial \rangle}\to N_{\widehat\Gamma_{1,1}}(\overline{\langle D_a\rangle})\xrightarrow{\pi} \overline{\langle D_a, (D_aD_bD_a)^2\rangle}.
        \end{align*} 
    As this is right exact and the elements $D_a,(D_aD_bD_a)^2\in \Gamma_{1,1}$ project to the corresponding elements in $\Gamma_1^1$, we get that $N_{\widehat\Gamma_{1,1}}(\overline{\langle D_a\rangle})$ is generated by 
        \[D_\partial, \quad D_a, \quad (D_aD_bD_a)^2.\]
    Note that all these generators commute with $D_a$ in $\Gamma_{1,1}$. In particular, this means that the entire normaliser of $\overline{\langle D_a\rangle}$ is simply just the centraliser of $D_a$ in $\widehat{\Gamma}_{1,1}$ and we have the equalities 
        \[C_{\widehat{\Gamma}_{1,1}}(D_a)= \overline{C_{\Gamma_{1,1}}(D_a)}=\overline{\langle D_\partial, D_a, (D_aD_bD_a)^2\rangle}.\]
\end{proof}

We are now ready to return to Proposition~\ref{prop:injection}:

\begin{proof}[Proof of Proposition~\ref{prop:injection}]
    Suppose that $F:\widehat{\bS}\to\widehat{\bS}$ is an object-fixing automorphism together with a natural transformation $H:F\Rightarrow\id_{\widehat{\bS}}.$ We will show that $F=\id_{\widehat{\bS}}$.

    We start by noting that the truncation $\trun{0}$ induces a commutative diagram 
        \[\begin{tikzcd}
            {\Aut_0(\widehat{\bS})} & {\HoAut(\widehat{\bS})} \\
            {\Aut_0(\trun{0}\widehat{\bS})} & {\HoAut(\trun{0}\widehat{\bS})}
            \arrow[from=1-1, to=1-2]
            \arrow["{\trun{0}}"', from=1-1, to=2-1]
            \arrow["{\trun{0}}", from=1-2, to=2-2]
            \arrow[from=2-1, to=2-2]
        \end{tikzcd}\]
    By Proposition~\ref{prop:PaRB-cyclic-is-S0} and \cite[Proposition 7.8]{Boavida-Horel-Robertson}, the bottom map is injective. Since $H:F\Rightarrow\id_{\widehat{\bS}}$ is a natural transformation, the induced map $\trun{0}F$ is homotopic to $\id_{\trun{0}\widehat{\bS}}$. Injectivity of the bottom map therefore implies
        \[\trun{0}F=\id_{\trun{0}\widehat{\bS}}.\]
    Since we know $F$ restricts to the identity in genus zero, by Theorem~\ref{thm: presentation for maps out of bS}, to prove that the automorphism $F$ is the identity on $\widehat\bS$, it suffices to show $F(\sigma)=\sigma$, where $\sigma$ is the standard $S$-move seen as a morphism of $\widehat\bS(1,1)$.

    Recall that $\widehat\bS(1,1)$ has a single object $\contraction$ and morphism group $\widehat\Gamma_{1,1}$. Let $a$ denote the curve in $S_{\contraction}$ depicted in Figure~\ref{fig:S-move}, and let $\sigma=D_aD_bD_a$
    be the standard $S$-move defined in Definition~\ref{def:standard moves}.
    
    The natural transformation  $H:F\Rightarrow\id_{\widehat{\bS}}$ implies there is a morphism $H_{\contraction}$ in $\widehat{\bS}(1,1)$ making the following squares commute
        \[\begin{tikzcd}
        	{{\contraction}} & {{\contraction}} \\
        	{{\contraction}} & {{\contraction}}
        	\arrow["{F(D_a)}", from=1-1, to=1-2]
        	\arrow["{H_{\contraction}}"', from=1-1, to=2-1]
        	\arrow["{H_{\contraction}}", from=1-2, to=2-2]
        	\arrow["{D_a}"', from=2-1, to=2-2]
        \end{tikzcd}
        \hspace{2cm}
        \begin{tikzcd}
        	{{\contraction}} & {{\contraction}} \\
        	{{\contraction}} & {{\contraction}}
        	\arrow["{F(\sigma)}", from=1-1, to=1-2]
        	\arrow["{H_{\contraction}}"', from=1-1, to=2-1]
        	\arrow["{H_{\contraction}}", from=1-2, to=2-2]
        	\arrow["{\sigma}"', from=2-1, to=2-2]
        \end{tikzcd}\]
    Since $D_a=\xi_{12}(\id_P\circ_1\tau)$, the mapping class $D_a$ is obtained from genus-zero morphisms. Since $F$ restricts to the identity in genus zero, we conclude that $F(D_a)=D_a$.
    
    In particular, then the left-hand commutative square implies that $H_{\contraction}$ is in the centraliser of $D_a$ in $\widehat\Gamma_{1,1}$. 

    By Lemma~\ref{lemma:centralisers in the profinite completion of Gamma11}, we get that $C_{\widehat{\Gamma}_{1,1}}(D_a)$ is generated by $D_\partial$, $(D_aD_bD_a)^2$, and $ D_a,$
    which gives that 
        \[H_{\contraction}=(D_aD_bD_a)^{2\chi}\cdot D_\partial^\mu \cdot D_a^\nu. \qquad \text{ for some } \chi,\mu,\nu\in\widehat{\Z}\]
    and
        \[F(\sigma)=((D_aD_bD_a)^{2\chi}\cdot D_\partial^\mu \cdot D_a^\nu)^{-1}\sigma ((D_aD_bD_a)^{2\chi}\cdot D_\partial^\mu \cdot D_a^\nu)\]

    As $\sigma=D_aD_bD_a$ commutes with $(D_aD_bD_a)^2$ and $D_\partial$ is in the centre of $\Gamma_{1,1}$, we get
        \[F(\sigma)= D_a^{-\nu}\cdot\sigma \cdot D_a^\nu= D_a^{-\nu}\cdot(D_aD_bD_a)\cdot D_a^\nu\]
    
    As $F$ must satisfy relation \eqref{g2_rel_morph_1}, using that $F$ is the identity in genus zero and the above formula for $F(\sigma)$, we get (using the curves and notation in Figure~\ref{fig:curves for g2})
        \begin{align}
            D_{\partial_0}^{\frac{1}{2}}D_{\partial_1}^{\frac{1}{2}} &= (D_{a_3}^{-\nu}\cdot(D_{a_3}D_{a_2}D_{a_3}) \cdot D_{a_3}^\nu)\cdot D_{a_3}^{-1} D_{a_1} \cdot (D_{a_3}^{-\nu}\cdot(D_{a_3}D_{a_2}D_{a_3}) \cdot D_{a_3}^\nu)\\
            &=  D_{a_3}^{-\nu}\cdot((D_{a_3}D_{a_2}D_{a_3}) \cdot D_{a_3}^{-1} D_{a_1} \cdot (D_{a_3}D_{a_2}D_{a_3})) \cdot D_{a_3}^\nu\\
            &=D_{a_3}^{-\nu}\cdot(D_{a_2}D_{a_3}D_{a_1}D_{a_2}D_{a_3}D_{a_2}) \cdot D_{a_3}^\nu.\label{eq:injection proof G2}
        \end{align}
    Now recall that $D_{\partial_0}^{\frac{1}{2}}D_{\partial_1}^{\frac{1}{2}}=(D_{a_2}D_{a_3}D_{a_1}D_{a_2}D_{a_3}D_{a_2})$ in $\Gamma_{1,1}$ and that, by the Lemma~\ref{lemma: conjugation by mapping classes}, we have
        \[(D_{\partial_0}^{\frac{1}{2}}D_{\partial_1}^{\frac{1}{2}})D_{a_3}(D_{\partial_0}^{\frac{1}{2}}D_{\partial_1}^{\frac{1}{2}})^{-1}=D_{a_1}.\]
    Putting these in equation \eqref{eq:injection proof G2}, we get
        \begin{align*}
        D_{a_2}D_{a_3}D_{a_1}D_{a_2}D_{a_3}D_{a_2} &=D_{a_3}^{-\nu}D_{a_1}^\nu\cdot(D_{a_2}D_{a_3}D_{a_1}D_{a_2}D_{a_3}D_{a_2}),
        \end{align*}    
    giving that
        $1= D_{a_3}^{-\nu}D_{a_1}^\nu.$
    Since $a_1$ and $a_3$ are distinct non-intersecting curves, this forces $\nu=0$. Therefore $F(\sigma)=\sigma$, as required.
\end{proof}

We next compare object-fixing automorphisms of $\widehat{\bS}$ with homotopy automorphisms of the modular $\infty$-operad $\widehat{\nerve\bS}$. Following \cite[1.2]{dk}, we write $\mathbb R\Aut(\calP)\subseteq \mathbb R\Map(\calP,\calP)$ for the union of those connected components whose images in $\pi_0\mathbb R\Map(\calP,\calP)$ are invertible. 

\begin{prop}\label{prop: strict-to-homotopy-endomorphisms}
There is an isomorphism of groups
\[
\Aut_0(\widehat{\bS})
\cong
\pi_0\mathbb R\Aut_{\Mod_{\infty}(\widehat{\Grpd})}(\widehat{\nerve\bS}).
\]
\end{prop}


\begin{proof}

We first compare the homotopy automorphisms of $\widehat{\nerve\bS}$ with the homotopy automorphisms of the strict profinite modular operad $\widehat{\bS}$. Since profinite completion is applied entrywise to modular dendroidal objects, the profinite completion adjunction gives a weak equivalence
\[
\mathbb R\Map_{\Mod_{\infty}(\widehat{\Grpd})}
(\widehat{\nerve\bS},\widehat{\nerve\bS})
\simeq
\mathbb R\Map_{\Mod_{\infty}(\Grpd)}
(\nerve\bS,|\widehat{\nerve\bS}|).
\]
By Lemma~\ref{lemma: N commutes with completion for groupoids}, $\widehat{\nerve\bS}\cong \nerve\widehat{\bS}$. Since the limit functor commutes with the modular dendroidal nerve, $|\widehat{\nerve\bS}|\cong \nerve|\widehat{\bS}|$. Hence the right-hand side above is weakly equivalent to $\mathbb R\Map_{\Mod_{\infty}(\Grpd)}
(\nerve\bS,\nerve|\widehat{\bS}|).$  By Theorem~\ref{thm: nerve is homotopically fully faithful}, we have a weak equivalence $\mathbb R\Map_{\Mod_{\infty}(\Grpd)}
(\nerve\bS,\nerve|\widehat{\bS}|)\simeq \mathbb R\Map_{\Mod(\Grpd)}(\bS,|\widehat{\bS}|)$.


Since $\bS$ is cofibrant by Proposition~\ref{prop: cofibrant modular operad}, and every modular operad in groupoids is fibrant, Theorem~\ref{thm: identifying mapping spaces in a simplicial model category} identifies
\[
\pi_0\mathbb R\Map_{\Mod(\Grpd)}(\bS,|\widehat{\bS}|)
\cong
\Hom_{\Ho(\Mod(\Grpd))}(\bS,|\widehat{\bS}|).
\]
Moreover, since entrywise profinite completion is the left adjoint in the Quillen adjunction $(-)^{\wedge}:\Mod(\Grpd)\rightleftarrows \Mod(\widehat{\Grpd}):|-|$, the associated derived adjunction identifies the right-hand side with $\Hom_{\Ho(\Mod(\widehat{\Grpd}))}(\widehat{\bS},\widehat{\bS})$. It follows that
\[
\pi_0\mathbb R\Map_{\Mod_{\infty}(\widehat{\Grpd})}(\widehat{\nerve\bS},\widehat{\nerve\bS}) \cong \Hom_{\Ho(\Mod(\widehat{\Grpd}))}(\widehat{\bS},\widehat{\bS}).
\]
Under the identifications above, composition is preserved. Hence the invertible components of \[\pi_0\mathbb R\Map_{\Mod_{\infty}(\widehat{\Grpd})}(\widehat{\nerve\bS},\widehat{\nerve\bS})\] correspond to the automorphisms of $\widehat{\bS}$ in the homotopy category. Therefore
\[
\pi_0\mathbb R\Aut_{\Mod_{\infty}(\widehat{\Grpd})}(\widehat{\nerve\bS}) \cong \HoAut(\widehat{\bS}).
\]


It remains to show that the map $\Aut_0(\widehat{\bS})\to \HoAut(\widehat{\bS})$ is an isomorphism. Injectivity was proven in Proposition~\ref{prop:injection}. 

To prove this map is surjective, it suffices to show that every automorphism of $\widehat{\bS}$ is homotopic to one inducing the identity on objects. This can be constructed exactly as in \cite[Theorem 7.8]{Horel_profinite_groupoids}:  the modular operad $\ob \widehat{\bS}$ is freely generated by the object $P$ in $\ob \widehat{\bS}(0,3)$. Hence the action of an automorphism $F:\widehat{\bS}\to \widehat{\bS}$ on objects is completely determined by the image of $P$. If $F(P)=P$, then $F$ already fixes all objects. Otherwise, we have $F(P)=(12)^*P$. In this case we construct another automorphism $F':\widehat{\bS}\to \widehat{\bS}$  which is the identity on objects together with a homotopy $H:F'\Rightarrow F$, which in this case is simply a natural transformation.

To construct $F'$, we pick a morphism $w:P\to (12)^*P$ in $\widehat{\bS}(0,3)$ (for instance $\beta$). Define $H_P:P\to (12)^*P$ to be this chosen morphism $w$. This induces a unique map of modular operads $H:\ob\widehat{\bS}\to \ob\widehat{\bS}$ and, in particular maps $H_G:G\to F(G)$ for any object $G$ in $\widehat{\bS}$. Now if $\phi:G\to G'$ is any morphism in $\widehat{\bS}$, we define 
    \[F'(\phi)=H_{G'} \cdot\phi\cdot H_G^{-1}.\]
It is simple to see that $F'$ indeed preserves the modular operadic compositions, since $H$ and $F$ do, and that $H_G$ assemble into a natural transformation $F'\Rightarrow F$.  It follows that we have an isomorphism $\Aut_0(\widehat{\bS})\cong \HoAut(\widehat{\bS})$.

\end{proof}






Theorem~\ref{main theorem NS action} gives an object-fixing action of the monoid $\underline{\NS}$ on $\widehat{\bS}$. Restricting to invertible elements gives an action of the group $\NS$ by object-fixing automorphisms of $\widehat{\bS}$. Proposition~\ref{prop: strict-to-homotopy-endomorphisms} identifies these strict object-fixing automorphisms with homotopy automorphisms of the profinite modular $\infty$-operad $\widehat{\nerve\bS}$.

\begin{cor}\label{cor: homotopy NS action}
The action of $\NS$ on $\widehat{\bS}$ induces a faithful action of $\NS$ on $\widehat{\nerve\bS}$ by homotopy automorphisms. Equivalently, there is an injective homomorphism
\[
\NS
\longrightarrow
\pi_0\mathbb R\Aut_{\Mod_{\infty}(\widehat{\Grpd})}
(\widehat{\nerve\bS}).
\]
\end{cor}

\begin{proof}
By Theorem~\ref{main theorem NS action}, the group $\NS$ acts on $\widehat{\bS}$ by object-fixing automorphisms. Proposition~\ref{prop: strict-to-homotopy-endomorphisms} identifies
\[
\Aut_0(\widehat{\bS})\cong \pi_0\mathbb R\Aut_{\Mod_{\infty}(\widehat{\Grpd})}
(\widehat{\nerve\bS}).
\]
Composing with the action map $\NS\longrightarrow \Aut_0(\widehat{\bS})$ with this identification gives the desired homomorphism.

Restricting the action above to the genus zero truncation recovers an action of the Grothendieck--Teichmüller group $\GT$ through homotopy automorphisms, i.e.  
% \[
% \GT
% \longrightarrow
% \pi_0\mathbb R\Aut_{\Mod_{\infty}(\widehat{\Grpd})}
% (\nerve\trun{0}\widehat{\bS}).
% \]
    \[\begin{tikzcd}
    	{\NS } & {\pi_0\mathbb R\Aut_{\Mod_{\infty}(\widehat{\Grpd})} (\widehat{\nerve\bS})} \\
    	\GT & {\pi_0\mathbb R\Aut_{\Mod_{\infty}(\widehat{\Grpd})} (\widehat{\nerve\trun{0}{\bS}}).}
    	\arrow[from=1-1, to=1-2]
    	\arrow[hook, from=1-1, to=2-1]
    	\arrow[from=1-2, to=2-2]
    	\arrow[from=2-1, to=2-2]
    \end{tikzcd}\]
By \cite[Theorem 8.4]{Boavida-Horel-Robertson}, the corresponding action of $\GT$ on the profinite genus zero surface operad is faithful, and therefore so is the action of $\NS$.
\end{proof}




\subsection{Genus-one truncations of modular dendroidal spaces}
Profinite completion of spaces does not generally preserve products. This creates a difficulty for modular $\infty$-operads in profinite spaces, because the Segal condition is formulated using products indexed by the vertices of a graph. In~\cite[Proposition~3.9]{Boavida-Horel-Robertson}, Boavida, Horel and the second author gave a sufficient criterion for profinite completion to commute with such products: if $X$ and $Y$ are connected spaces whose homotopy groups are good in the sense of Definition~\ref{def: good groups}, then the natural map
\[
\widehat{X\times Y}\longrightarrow \widehat{X}\times \widehat{Y}
\]
is a weak equivalence of profinite spaces.

To extend the $\NS$-action on $\widehat{\bS}$ to an action in the homotopy category of profinite spaces, we need to establish compatibility between profinite completion and the products appearing in the Segal condition. The available product comparison requires goodness of the relevant homotopy groups. Since the mapping class groups are known to be good only in genus at most one, we first work with genus-one truncated modular dendroidal objects.

We begin by introducing genus truncations for modular dendroidal objects. Let $\bM_{\leq g}\subseteq \bM$ denote the full subcategory spanned by genus-graded graphs of total genus at most $g$.


\begin{lemma}
The full subcategory $\bM_{\leq g}$ is a sieve on $\bM$. That is, if $\phi\colon G\to H$ is a morphism in $\bM$ and $H\in\bM_{\leq g}$, then $G\in\bM_{\leq g}$.
\end{lemma}

\begin{proof}
We use the description of graphical maps by graph substitution from~\cite[Theorem~2.7]{hry1}. A morphism $\phi\colon G\to H$ factors through a substituted graph $G\{K_v\}_{v\in V_G}$, where each graph $K_v$ has total genus equal to the genus assigned to the vertex $v$ of $G$, and the remaining map preserves genus.

The total genus of the substituted graph is
\[
\beta_1(G)+\sum_{v\in V_G}g(K_v)
=
\beta_1(G)+\sum_{v\in V_G}\epsilon_G(v),
\]
which is precisely the total genus of $G$. Since the map $G\{K_v\}_{v\in V_G}\to H$ preserves genus, the total genus of $G$ agrees with the total genus of its image in $H$, and in particular is at most the total genus of $H$. Thus, if $H$ has total genus at most $g$, then so does $G$.
\end{proof}

Let $i_g\colon \bM_{\leq g}\hookrightarrow\bM$ denote the inclusion. Restriction along $i_g$ defines a functor
\[
i_g^*\colon \dMod(\bE)\longrightarrow \dMod_{\leq g}(\bE).
\]
Here, we recall that $\dMod(\bE)=\operatorname{Fun}(\bM^{\op},\bE)$ denotes the category of modular dendroidal objects in $\bE$, and $\dMod_{\leq g}(\bE)=\operatorname{Fun}(\bM_{\leq g}^{\op},\bE)$ denotes the category of genus-$g$ truncated modular dendroidal objects in $\bE$. We call $i_g^*X$ the genus-$g$ truncation of $X$ and write it as $\trun{g}X$.

We now turn to the comparison between profinite completion of groupoids and profinite completion of their classifying spaces. There are two natural ways to pass from a groupoid to a profinite space: one can first take the classifying space and then profinitely complete, or one can first profinitely complete the groupoid and then take its classifying space. These two constructions are related by a natural comparison map, represented by the square
\[
\begin{tikzcd}
\Grpd \arrow[r, "\cs"] \arrow[d, "\widehat{(-)}"'] &
\sSet \arrow[d, "\widehat{(-)}"] \\
\widehat{\Grpd} \arrow[r, "\cs"'] &
\widehat{\sSet}.
\end{tikzcd}
\]
At a groupoid $\mathscr E$, this comparison is the canonical map $\widehat{\cs\mathscr E}\to \cs\widehat{\mathscr E}$.

For a general groupoid $\mathscr E$, this map need not be a weak equivalence. A sufficient condition is goodness in the sense of Serre.

\begin{definition}\label{def: good groups}
A discrete group $\mathsf G$ is \emph{good} if, for every finite discrete $\mathsf G$-module $\mathsf M$, the natural map $H^*_{\mathrm{cts}}(\widehat{\mathsf G},\mathsf M)\to H^*(\mathsf G,\mathsf M)$ is an isomorphism. A groupoid with finitely many objects is called good if the automorphism group of each object is good.
\end{definition}

By~\cite[Proposition~5.9]{Horel_profinite_groupoids}, if $\mathscr E$ is a good groupoid with finitely many objects, then the comparison map $\widehat{\cs\mathscr E}\to\cs\widehat{\mathscr E}$ is a weak equivalence of profinite spaces.

\begin{lemma}\label{lemma: B of good groupoids commutes with hat}
For each $g\leq 1$ and $n\geq 0$ satisfying $2g+n\geq 1$, the canonical map
\[
\widesthat{\cs\bS(g,n+1)}
\longrightarrow
\cs(\widehat{\bS}(g,n+1))
\]
is a weak equivalence of profinite spaces.
\end{lemma}

\begin{proof}
The groupoid $\bS(g,n+1)$ has finitely many objects. For every object $A\in\ob(\bS(g,n+1))$, the automorphism group $\Aut_{\bS(g,n+1)}(A)$ is isomorphic to the mapping class group $\Gamma_{g,n+1}$. These groups are good for $g\leq 1$; see~\cite[p.~94]{Oda}. Hence $\bS(g,n+1)$ is a good groupoid and the result therefore follows from~\cite[Proposition~5.9]{Horel_profinite_groupoids}.
\end{proof}

\begin{prop}\label{prop: genus-one B N and completion commute}
There is an equivalence of genus-one truncated modular $\infty$-operads in profinite spaces
\[
\widesthat{\cs\nerve(\trun{1}\bS)}\simeq \cs\nerve(\widehat{\trun{1}\bS}).
\]
\end{prop}

\begin{proof}
Both sides are functors $\bM_{\leq 1}^{\op}\to\widehat{\isSet}$. We first verify that both satisfy the truncated Segal condition.

Let $G\in\bM_{\leq 1}$. Since $\trun{1}\bS$ is a truncated modular operad, its modular dendroidal nerve satisfies the strict Segal formula
\[
(\nerve\trun{1}\bS)_G
\cong
\prod_{v\in V_G}
(\nerve\trun{1}\bS)_{\corolla_{(\epsilon_G(v),|\nb_G(v)|)}}
\cong
\prod_{v\in V_G}
\bS(\epsilon_G(v),|\nb_G(v)|),
\]
where the second isomorphism follows from the definition of the modular dendroidal nerve. Since the classifying space functor preserves finite products, we have
\[
(\cs\nerve\trun{1}\bS)_G
\cong
\prod_{v\in V_G}
\cs\bS(\epsilon_G(v),|\nb_G(v)|).
\]

Thus the Segal map for $\widesthat{\cs\nerve(\trun{1}\bS)}$ at $G$ is the natural map
\[
\left(
\prod_{v\in V_G}
\cs\bS(\epsilon_G(v),|\nb_G(v)|)
\right)^{\wedge}
\longrightarrow
\prod_{v\in V_G}
\bigl(\cs\bS(\epsilon_G(v),|\nb_G(v)|)\bigr)^{\wedge}.
\]
Since $G$ has total genus at most one, every vertex has genus at most one. For each $v\in V_G$, the groupoid $\bS(\epsilon_G(v),|\nb_G(v)|)$ is connected, and the automorphism group of each object is isomorphic to $\Gamma_{\epsilon_G(v),|\nb_G(v)|}$ and these mapping class groups are good under our assumptions. Hence the spaces $\cs\bS(\epsilon_G(v),|\nb_G(v)|)$ satisfy the hypotheses of~\cite[Proposition~3.9]{Boavida-Horel-Robertson}. Applying that product comparison iteratively, the displayed Segal map is a weak equivalence. 

For the right-hand side, the nerve of the truncated profinite modular operad $\widehat{\trun{1}\bS}$ satisfies the Segal condition,
\[
(\nerve\widehat{\trun{1}\bS})_G
\cong
\prod_{v\in V_G}
(\nerve\widehat{\trun{1}\bS})_{\corolla_{(\epsilon_G(v),|\nb_G(v)|)}}
\cong
\prod_{v\in V_G}
\widehat{\bS}(\epsilon_G(v),|\nb_G(v)|).
\] Again noting that the classifying space functor preserves this finite products, we have that $\cs\nerve(\widehat{\trun{1}\bS})$ also satisfies the truncated Segal condition.

On a corolla $\corolla_{(g,n+1)}$ with $g\leq 1$, the comparison is the canonical map
$\widesthat{\cs(\bS(g,n+1))}\to \cs(\widehat{\bS}(g,n+1))$, which is a weak equivalence by Lemma~\ref{lemma: B of good groupoids commutes with hat}. The comparison therefore extends to an equivalence on arbitrary graphs by the truncated Segal condition. Indeed, for every $G\in\bM_{\leq 1}$, it fits into a commutative square
\[
\begin{tikzcd}
(\widesthat{\cs\nerve(\trun{1}\bS)})_G
\arrow[r]
\arrow[d, "\simeq"'] &
(\cs\nerve(\widehat{\trun{1}\bS}))_G
\arrow[d, "\simeq"] \\
\displaystyle\prod_{v\in V_G}
(\widesthat{\cs\nerve(\trun{1}\bS)})_{\corolla_{(\epsilon_G(v),|\nb_G(v)|)}}
\arrow[r] &
\displaystyle\prod_{v\in V_G}
(\cs\nerve(\widehat{\trun{1}\bS}))_{\corolla_{(\epsilon_G(v),|\nb_G(v)|)}} .
\end{tikzcd}
\]
The vertical maps are the Segal equivalences established above, and the bottom horizontal map is the product of the comparison maps on the corresponding corollas. Hence the bottom horizontal map is a weak equivalence, and thus so is the top horizontal map.

It follows that $\widesthat{\cs\nerve(\trun{1}\bS)}\to \cs\nerve(\widehat{\trun{1}\bS})$ is an equivalence of genus-one truncated modular $\infty$-operads in profinite spaces.
\end{proof}

The preceding comparison allows us to transport the genus-one part of the $\underline{\NS}$-action from profinite groupoids to profinite spaces.

\begin{prop}\label{prop: genus-one NS action on spaces}
The action of $\NS$ on $\trun{1}\widehat{\bS}$ induces a faithful action on $\widesthat{\cs\nerve(\trun{1}\bS)}$ by automorphisms in the homotopy category of genus-one truncated modular $\infty$-operads in profinite spaces. More precisely, there is an injective homomorphism
\[\NS \longrightarrow \pi_0\mathbb  \Aut_{\dMod_{\leq 1}(\widehat{\isSet})} \bigl(\widehat{\cs\nerve(\trun{1}\bS)}\bigr). \]
\end{prop}

\begin{proof}
By Theorem~\ref{main theorem NS action}, the group $\NS$ acts on $\trun{1}\widehat{\bS}$ by object-fixing endomorphisms. The genus-one restriction of Proposition~\ref{prop: strict-to-homotopy-endomorphisms} gives a homomorphism
\begin{equation}\label{ns action up to homotopy}
{\NS}\longrightarrow
\pi_0\mathbb R\Aut_{\dMod_{\leq 1}(\widehat{\Grpd})}
\bigl(\widehat{\nerve(\trun{1}\bS)}\bigr).
\end{equation}
By Lemma~\ref{lemma: N commutes with completion for groupoids}, we identify $\widehat{\nerve(\trun{1}\bS)}$ with $\nerve(\widehat{\trun{1}\bS})$.

Applying the classifying space functor entrywise, and using the homotopical full faithfulness of $\cs$ recalled in Section~\ref{sec: classifying space functor}, the homomorphism \eqref{ns action up to homotopy} induces
\[
{\NS}\longrightarrow
\pi_0\mathbb R\Aut_{\dMod_{\leq 1}(\widehat{\isSet})}
\bigl(\cs\nerve(\widehat{\trun{1}\bS})\bigr).
\]
The comparison map of Proposition~\ref{prop: genus-one B N and completion commute} identifies this target with
$\pi_0\mathbb R\Aut_{\dMod_{\leq 1}(\widehat{\isSet})}
\bigl(\widehat{\cs\nerve(\trun{1}\bS)}\bigr)$,
giving the stated homomorphism.
% Restricting to invertible elements gives the corresponding group homomorphism for $\NS$. Since elements of $\NS$ act by automorphisms on $\trun{1}\widehat{\bS}$, their images are homotopy-invertible after applying $\cs$ and passing through the comparison above. Hence the restricted homomorphism lands in
% \[
% \pi_0\mathbb R\Aut_{\dMod_{\leq 1}(\widehat{\isSet})}
% \bigl(\widehat{\cs\nerve(\trun{1}\bS)}\bigr).
% \]
\end{proof}

\subsection{Étale homotopy types and the profinite-space model}\label{sec: tower of etale homotopy types}
We now show that the entrywise profinite completion of $\cs\nerve\bS$ assembles into a modular $\infty$-operad in profinite spaces. This subsection addresses a different issue from the genus-one comparison above: here we need to verify that profinite completion still satisfies the Segal conditions. Since profinite completion of spaces does not generally preserve products, this requires a separate product-comparison argument.

We write $\isSet$ for the $\infty$-category of spaces, and let $\isSet_{\pi}\subset \isSet$ denote the full subcategory of $\pi$-finite spaces. We write $\widehat{\isSet}:=\Pro(\isSet_{\pi})$ for the $\infty$-category of profinite spaces, and
$\widehat{(-)}\colon \Pro(\isSet)\to\widehat{\isSet}$ for profinite completion.

Let $\isSet_{<\infty}\subset\isSet$ denote the full subcategory of truncated spaces. Following~\cite[Notation~2.11]{haine2024profinite}, let $\mathcal H\subset\Pro(\isSet_{<\infty})$ denote the smallest full subcategory containing
$\widehat{\isSet}$ and closed under geometric realizations, retracts, and cofiltered limits. We call a prospace $X\in\Pro(\isSet)$ \emph{admissible} if its protruncation $\tau_{<\infty}X$ lies in $\mathcal H$.

\begin{prop}\label{prop:haine-product}
Let $X_1,\ldots,X_r$ be admissible prospaces. Then the natural map
\[
\left(\prod_{i=1}^r X_i\right)^{\wedge}
\longrightarrow
\prod_{i=1}^r \widehat{X_i}
\]
is an equivalence in profinite spaces.
\end{prop}

\begin{proof}
Profinite completion factors through protruncation, and protruncation preserves finite products. Thus it suffices to prove the corresponding statement for objects of $\mathcal H\subset\Pro(\isSet_{<\infty})$. This is the all-primes finite-product case of~\cite[Theorem~2.13]{haine2024profinite}, applied iteratively.
\end{proof}

For a quasi-compact and quasi-separated (qcqs) scheme $X$, the protruncated étale homotopy type $\Pi^{\acute et}_{<\infty}(X):=\tau_{<\infty}\Pi^{\acute et}_{\infty}(X)$ admits a $\Pro(\isSet_{\pi})$-resolution. Equivalently, $\Pi^{\acute et}_{<\infty}(X)$ belongs to $\mathcal H$. Thus $\Pi^{\acute et}_{\infty}(X)$ is admissible. Haine records the corresponding product comparison for two qcqs schemes in~\cite[Example~2.17]{haine2024profinite}; the finite-product version follows from Proposition~\ref{prop:haine-product}.

\begin{cor}\label{cor:haine-qcqs-schemes}
Let $X_1,\ldots,X_r$ be quasi-compact and quasi-separated schemes. Then the natural map
\[
\left(\prod_{i=1}^r \Pi^{\acute et}_{\infty}(X_i)\right)^{\wedge}
\longrightarrow
\prod_{i=1}^r \widehat{\Pi^{\acute et}_{\infty}(X_i)}
\]
is an equivalence in profinite spaces.
\end{cor}

As in Example~\ref{example: Oda}, the profinite étale homotopy type of the moduli stack $\mathcal M_g^n$ of smooth genus $g$ curves with $n$ ordered marked points identifies with the classifying space of the profinite completed marked-point mapping class group:
\[
\Pi^{\acute et}_{\infty}(\mathcal M_g^n)
\simeq
\widehat{\cs\Gamma_g^n}.
\]
Our modular operad, however, involves the mapping class groups $\Gamma_{g,n}$ of surfaces with $n$ boundary components. To pass from marked points to boundary components, we add a non-zero tangent vector at each marked point.

Let $\mathcal M_{g,n}^{\tan}$ denote the stack of smooth genus $g$ curves with $n$ ordered marked points equipped with non-zero tangent vectors at the marked points. Equivalently, $\mathcal M_{g,n}^{\tan}$ is the $(\mathbb G_m)^n$-torsor over $\mathcal M_g^n$ obtained by deleting the zero sections from the cotangent-line bundles. Topologically, this replaces $\Gamma_g^n$ by the central extension
\[
1\longrightarrow \mathbb Z^n\longrightarrow \Gamma_{g,n}
\longrightarrow \Gamma_g^n\longrightarrow 1,
\]
where the central subgroup is generated by the boundary Dehn twists.

The stack $\mathcal M_{g,n}^{\tan}$ is an algebraic stack of finite type over the base field. Since it is not generally a scheme, we cannot apply Haine's scheme-level product comparison directly. Instead, we first verify that its étale homotopy type is admissible in the sense above.

\begin{lemma}\label{lemma:tangent-moduli-admissible}
For every pair $(g,n)$, the prospace
$\Pi_\infty^{\acute et}(\mathcal M_{g,n}^{\tan})$
is admissible.
\end{lemma}

\begin{proof}
The stack $\mathcal M_{g,n}^{\tan}$ is a quasi-compact algebraic stack of finite type over the base field. Since it is quasi-compact, we can choose a smooth surjective atlas $U\to \mathcal M_{g,n}^{\tan}$ with $U$ an affine scheme \cite[\href{https://stacks.math.columbia.edu/tag/04YC}{04YC}]{stacks-project}. We may choose $U$ to be of finite type over the base field, hence quasi-compact and quasi-separated. Let $U_\bullet$ be the Cech nerve of this atlas.

Since smooth morphisms admit étale local sections, the atlas $U\to \mathcal M_{g,n}^{\tan}$ is an effective epimorphism for the étale stack topology. Thus $\mathcal M_{g,n}^{\tan}$ is the geometric realization of the simplicial stack $U_\bullet$. By~\cite[Lemma~2.26]{Carchedi}, the assignment $X\mapsto \operatorname{Sh}_\infty(X_{\acute et})$ extends from affine schemes to higher stacks as a colimit-preserving functor to $\infty$-topoi. Since the shape functor preserves colimits, and since $\Pi_\infty^{\acute et}$ is defined as the shape of the small étale $\infty$-topos, we obtain
\[
\Pi_\infty^{\acute et}(\mathcal M_{g,n}^{\tan})
\simeq
\left|\,\Pi_\infty^{\acute et}(U_\bullet)\,\right|
\]
in $\Pro(\isSet)$.

It remains to note that the terms $U_k$ are quasi-compact and quasi-separated schemes. The diagonal of $\mathcal M_{g,n}^{\tan}$ is representable by algebraic spaces and locally of finite type \cite[\href{https://stacks.math.columbia.edu/tag/04XS}{04XS}]{stacks-project}. In this moduli problem, the diagonal is in fact represented by the Isom scheme between pointed smooth curves equipped with non-zero tangent vectors. Hence the iterated fibre products
$U_k=U\times_{\mathcal M_{g,n}^{\tan}}\cdots\times_{\mathcal M_{g,n}^{\tan}}U$
are schemes. Since $U$ is of finite type over the base field and the relevant Isom schemes are of finite type, each $U_k$ is of finite type over the base field. In particular, each $U_k$ is quasi-compact and quasi-separated.

For each $k$, the scheme $U_k$ is qcqs, so its étale $\infty$-topos is spectral. Hence $\tau_{<\infty}\Pi_\infty^{\acute et}(U_k)$ admits a natural $\Pro(\isSet_\pi)$-resolution by~\cite[Example~A.9]{haine2024profinite}; equivalently, $\Pi_\infty^{\acute et}(U_k)$ is admissible in the sense above. Since protruncation preserves geometric realizations, we have
\[
\tau_{<\infty}\Pi_\infty^{\acute et}(\mathcal M_{g,n}^{\tan})
\simeq
\left|\,\tau_{<\infty}\Pi_\infty^{\acute et}(U_\bullet)\,\right|.
\]
The right-hand side lies in $\mathcal H$, because $\mathcal H$ is closed under geometric realizations. Therefore $\Pi_\infty^{\acute et}(\mathcal M_{g,n}^{\tan})$ is admissible.
\end{proof}

Now that admissibility is established, we can compare the profinite étale homotopy types of these stacks with the profinite completed classifying spaces of the corresponding mapping class groups. This comparison is the topological input that connects the algebro-geometric side of the construction to the modular operadic one.
\begin{lemma}\label{lemma:tangent-moduli-comparison}
For every $g\geq 0$ and $n\geq 0$ satisfying $2g+n\geq 1$, there is an equivalence in profinite spaces
\[
\widehat{\Pi_\infty^{\acute et}(\mathcal M_{g,n+1}^{\tan})}
\simeq
\widehat{\cs\Gamma}_{g,n+1}.
\]
\end{lemma}

\begin{proof}
By the comparison theorem for algebraic stacks locally of finite type over $\mathbb C$, the profinite étale homotopy type of $\mathcal M_{g,n+1}^{\tan}$ agrees with the profinite completion of the homotopy type of its associated topological stack. The topological stack is presented by the Teichmüller space of genus $g$ surfaces with $n+1$ tangent directions modulo the mapping class group $\Gamma_{g,n+1}$. Since this Teichmüller space is contractible, the weak homotopy type of the topological stack is $\cs\Gamma_{g,n+1}$. Hence
\[
\widehat{\Pi_\infty^{\acute et}(\mathcal M_{g,n+1}^{\tan})}
\simeq
\widehat{\cs\Gamma}_{g,n+1}.
\]
\end{proof}

With the above comparison in hand, we can show that profinite completion of classifying spaces of mapping class groups behaves well on the products that appear in the Segal maps. 

\begin{prop}\label{prop:haine-boundary-product}
Let $g_1,\ldots,g_r\geq 0$ and $n_1,\ldots,n_r\geq 0$. Then the natural map
\[
\left(\prod_{i=1}^r \cs\Gamma_{g_i,n_i+1}\right)^{\wedge}
\longrightarrow
\prod_{i=1}^r(\cs\Gamma_{g_i,n_i+1})^{\wedge}
\]
is an equivalence in profinite spaces.
\end{prop}

\begin{proof}
By Lemma~\ref{lemma:tangent-moduli-admissible}, each prospace
$\Pi_\infty^{\acute et}(\mathcal M_{g_i,n_i+1}^{\tan})$ is admissible. Proposition~\ref{prop:haine-product} therefore gives an equivalence
\[
\left(
\prod_{i=1}^r
\Pi_\infty^{\acute et}(\mathcal M_{g_i,n_i+1}^{\tan})
\right)^{\wedge}
\simeq
\prod_{i=1}^r
\widehat{\Pi_\infty^{\acute et}(\mathcal M_{g_i,n_i+1}^{\tan})}.
\]
Using Lemma~\ref{lemma:tangent-moduli-comparison} on each factor identifies the right-hand side with
\[
\prod_{i=1}^r(\cs\Gamma_{g_i,n_i+1})^{\wedge}.
\]
The left-hand side is identified with
\[
\left(\prod_{i=1}^r \cs\Gamma_{g_i,n_i+1}\right)^{\wedge}
\]
by applying the same comparison to the product of the corresponding tangent-direction stacks. This gives the desired product comparison.
\end{proof}

This product comparison is exactly what is needed to verify the Segal condition for the profinite completion of the modular dendroidal nerve of $\bS$.

\begin{lemma}\label{lemma:teichmuller-tower-is-modular-infty-operad}
The composite
\[
\bM^{\op}\xrightarrow{\nerve\bS}\Grpd\xrightarrow{\cs}\sSet\xrightarrow{\widehat{(-)}}\widehat{\isSet}
\]
defines a reduced modular $\infty$-operad in profinite spaces.
\end{lemma}

\begin{proof}
Let $G$ be a graph with vertices $v_1,\ldots,v_r$, and write
$\epsilon(v_i)=g_i$ and $|\nb(v_i)|=n_i+1$. Since $\cs$ preserves finite products and weak equivalences, the Segal equivalence for the modular dendroidal nerve gives
\[
(\cs\nerve\bS)_G
\simeq
\prod_{i=1}^r \cs\bS(g_i,n_i+1).
\]
The groupoid $\bS(g_i,n_i+1)$ is connected with automorphism group $\Gamma_{g_i,n_i+1}$, so
$\cs\bS(g_i,n_i+1)\simeq \cs\Gamma_{g_i,n_i+1}$. Proposition~\ref{prop:haine-boundary-product} then identifies the profinite completion of the product above with the product of the completed corolla values. Thus the Segal map
\[
(\widehat{\cs\nerve\bS})_G
\longrightarrow
\prod_{i=1}^r
(\widehat{\cs\nerve\bS})_{\corolla_{(g_i,n_i+1)}}
\]
is a weak equivalence for every graph $G$. The exceptional edge is still sent to the point. Hence $\widehat{\cs\nerve\bS}$ is a modular $\infty$-operad in profinite spaces.
\end{proof}

Having constructed the profinite-space analogue of our modular operad in groupoids, we can now revisit the two-level principle in this homotopy-coherent setting. The statement for $\bS$ says that genus-one data controls the entire tower (Corollary~\ref{cor: truncation and endomorphisms of the surface mod operad}). The same phenomenon persists after passing to the modular $\infty$-operad in profinite spaces.


\begin{prop}\label{prop: two-level profinite space model}
Restriction to the genus-one truncation induces an isomorphism
\[
\pi_0\mathbb R\End_{\dMod(\widehat{\isSet})}(\widehat{\cs\nerve\bS})
\longrightarrow
\pi_0\mathbb R\End_{\dMod_{\leq 1}(\widehat{\isSet})}(\trun{1}\widehat{\cs\nerve\bS}).
\]
\end{prop}

\begin{proof}
This is the profinite-space form of the two-level principle. The strict two-level result, Corollary~\ref{cor: maps out of bS}, says that maps out of $\bS$ are determined by their restrictions to the genus-one truncation because all generators and defining relations occur in genus at most one. Passing to the modular dendroidal nerve replaces this strict presentation by the corresponding homotopy-coherent Segal presentation. By Lemma~\ref{lemma:teichmuller-tower-is-modular-infty-operad}, the profinite completion $\widehat{\cs\nerve\bS}$ satisfies the Segal condition, so the higher-genus values are recovered from the same genus-one generators and relations by homotopy-coherent composition. Therefore restriction to $\bM_{\leq 1}$ induces an isomorphism on $\pi_0$ of derived endomorphism monoids.
\end{proof}

The action of $\NS$ on $\widehat{\bS}$ was constructed using the two-level principle for $\bS$, namely the fact that the genus-zero and genus-one generators and relations determine the full modular operad. After applying the modular dendroidal nerve and profinite classifying spaces, the comparison of strict automorphisms with homotopy automorphisms gives an action on the profinite-space model.

\begin{thm}\label{NS action on topological operad}
The action of $\NS$ on $\widehat{\bS}$ induces an action of $\NS$ on $\widehat{\cs\nerve\bS}$ by automorphisms in the homotopy category of modular $\infty$-operads in profinite spaces. More precisely, there is a homomorphism
\[
\NS\longrightarrow
\pi_0\mathbb R\Aut_{\dMod(\widehat{\isSet})}
(\widehat{\cs\nerve\bS}).
\]
\end{thm}

\begin{proof}
By Theorem~\ref{main theorem NS action}, the group $\NS$ acts on $\widehat{\bS}$ by object-fixing automorphisms. Applying the modular dendroidal nerve and profinite classifying spaces gives self-maps of $\widehat{\cs\nerve\bS}$. Lemma~\ref{lemma:teichmuller-tower-is-modular-infty-operad} shows that $\widehat{\cs\nerve\bS}$ is a modular $\infty$-operad in profinite spaces. The comparison between strict object-fixing automorphisms and homotopy automorphisms, together with the classifying-space functor, therefore sends the action map $\NS\longrightarrow \Aut_0(\widehat{\bS})$ to a homomorphism
\[
\NS\longrightarrow
\pi_0\mathbb R\Aut_{\dMod(\widehat{\isSet})}
(\widehat{\cs\nerve\bS}).
\]
\end{proof}
\appendix
\section{Homotopy Theory of Modular Operads}
\label{sec: homotopy theory of mod op and truncated mod op}

Modular operads can be viewed as algebras over a certain coloured operad (cf.~\cite[Appendix~A]{dch_coextension}). This perspective allows us to apply Theorem~2.1 of~\cite{bm_resolutions}, which ensures that, under suitable hypotheses on the base category $\bE$, categories of algebras over coloured operads inherit transferred model structures. More precisely, if $\bE$ is a cofibrantly generated symmetric monoidal model category with cofibrant unit, a symmetric monoidal fibrant replacement functor, and a cocommutative coalgebra interval, then the category of algebras over such a coloured operad admits a cofibrantly generated model structure in which weak equivalences and fibrations are defined entrywise in $\bE$.

Since our arguments rely on several specific features of the coloured operad governing modular operads, we begin with a brief review of coloured operads and
their algebras.

\begin{definition}
Let $\mathfrak{C}$ be a nonempty set of \emph{colours}. A $\mathfrak{C}$-coloured \emph{sequence} in $\bE$ is a collection of objects
\[
\mathcal{P}(c_0;c_1,\ldots,c_k)
\]
indexed by lists $(c_0;c_1,\ldots,c_k)$ of colours in $\mathfrak{C}$, together with right actions of the symmetric groups. Thus, for each $\sigma\in\Sigma_k$, there is an isomorphism
\[
\begin{tikzcd}
\mathcal{P}(c_0;c_1,\ldots,c_k) \arrow[r,"\sigma^*"] & \mathcal{P}(c_0;c_{\sigma(1)},\ldots,c_{\sigma(k)}).
\end{tikzcd}
\]

A $\mathfrak{C}$-coloured \emph{operad} is a $\mathfrak{C}$-coloured sequence $\mathcal{P}$ equipped with:
\begin{enumerate}
\item distinguished unit objects $\iota_c \in \mathcal{P}(c;c)$;
%equivalently maps $*\to \mathcal{P}(c;c)$ from the terminal object of $\bE$, for each $c\in\mathfrak{C}$;

    \item composition maps
    \[
    \begin{tikzcd}[column sep=large]
    \mathcal{P}(c_0;c_1,\ldots,c_k)\times \mathcal{P}(d_0;d_1,\ldots,d_j) \arrow[r,"\circ_i"] & \mathcal{P}(c_0;c_1,\ldots,c_{i-1},d_1,\ldots,d_j,c_{i+1},\ldots,c_k),
    \end{tikzcd}
    \] defined whenever $c_i=d_0$ and $1\leq i\leq k$.
\end{enumerate}
These data are required to satisfy the standard associativity, unitality, and equivariance axioms.
\end{definition}

A map $f:\mathcal{P}\to\mathcal{Q}$ of coloured operads is a map of the underlying coloured sequences that respects the units, compositions, and symmetric group actions. The category of $\mathfrak{C}$-coloured operads in $\bE$ is denoted $\Op_{\mathfrak{C}}(\bE)$. For further details, see \cite[Definition~1.1]{bm_resolutions} or \cite[§2.2]{DCH}.

\begin{definition}
Let $\mathcal{P}$ be a $\mathfrak{C}$-coloured operad. A \emph{$\mathcal{P}$-algebra} in $\bE$ is a collection of objects $\{X(c)\}_{c\in\mathfrak{C}}$ together with structure maps
\[
\begin{tikzcd}[column sep=large] \mathcal{P}(c_0;c_1,\ldots,c_k)\times X(c_1)\times\cdots\times X(c_k) \arrow[r,"\rho"] & X(c_0) \end{tikzcd}
\]
satisfying associativity, unitality, and equivariance with respect to the operad structure. The category of such algebras is denoted $\Alg_{\mathcal{P}}(\bE)$.
\end{definition}

\begin{example}\label{def: operad for modular operads}
The \emph{operad of modular operads} is an $\mathbb{N}\times\mathbb{N}_{\ge 1}$-coloured operad $\mathcal{MOp}$. For a list of pairs $\bigl((g,r);(g_1,r_1),\ldots,(g_n,r_n)\bigr),$ the object
\[
\mathcal{MOp}\bigl((g,r);(g_1,r_1),\ldots,(g_n,r_n)\bigr)
\]
is the set of strict isomorphism classes of connected genus-graded, labelled graphs $G=(G,\lambda,\ell,\epsilon)$ with $n$ vertices and $r$ legs, satisfying:
\begin{itemize}
    \item the vertex labelled by $i$ has valence $|\nb(v_{\lambda(i)})|=r_i+1$ and genus $\epsilon(v_{\lambda(i)})=g_i$; 
    \item the total genus of $G$ is $g$, equivalently
    \[
    g=\beta_1(G)+\sum_{i=1}^n g_i.
    \]
\end{itemize}

The symmetric group $\Sigma_n$ acts on $\mathcal{MOp}\bigl((g,r);(g_1,r_1),\ldots,(g_n,r_n)\bigr)$ by permuting the vertex labelling $\lambda$.

\smallskip
Operadic composition is defined by graph substitution. Given
\[
G\in \mathcal{MOp}\bigl((g,r);(g_1,r_1),\ldots,(g_n,r_n)\bigr) \quad \text{and} \quad  H_{v_{\lambda(i)}}\in \mathcal{MOp}\bigl((g_i,r_i);(h_1,s_1),\ldots,(h_m,s_m)\bigr),
\]
the composite $G\circ_i H_{v_{\lambda(i)}}$ is obtained by substituting the
graph $H_{v_{\lambda(i)}}$ into the vertex $v_{\lambda(i)}$ of $G$:
\[
G\circ_i H_{v_{\lambda(i)}}:=G\{H_{v_{\lambda(i)}}\},
\]
as in~\cite[Construction~1.18]{hry1}. This gives an element of
\[
\mathcal{MOp}\bigl((g,r);
(g_1,r_1),\ldots,(g_{i-1},r_{i-1}),
(h_1,s_1),\ldots,(h_m,s_m),
(g_{i+1},r_{i+1}),\ldots,(g_n,r_n)\bigr).
\]

Associativity, equivariance, and unitality follow from the corresponding properties of graph substitution. For a detailed construction and verification, see~\cite[Appendix~A]{dch_coextension} or~\cite[§13.4]{yau_modular}.
\end{example}


\begin{prop}\label{prop: modular operads are algebras over operads}
The category of algebras over the coloured operad $\mathcal{MOp}$ is naturally isomorphic to the category of modular operads in $\bE$:
\[
\Alg_{\mathcal{MOp}}(\bE)\cong \Mod(\bE).
\]
\end{prop}

This identification allows us to import general results about algebras over coloured operads. In particular, under suitable hypotheses on the base category $\bE$, it gives a transferred model structure on $\Mod(\bE)$.

We will use this when $\bE$ is one of the following Cartesian monoidal model categories:

\begin{itemize}
  \item the category of simplicial sets $\sSet$, equipped with the Kan--Quillen model structure;
  \item the category of groupoids $\Grpd$, equipped with the canonical model structure in which weak equivalences are equivalences of categories~\cite{And};
  \item the category of profinite groupoids $\widehat{\Grpd}$, equipped with the cofibrantly generated model structure described by Horel~\cite[Theorem~4.12]{Horel_profinite_groupoids}.
\end{itemize}

In each case, $\bE$ is a cofibrantly generated Cartesian monoidal model category with cofibrant unit, a symmetric monoidal fibrant replacement functor, and a cocommutative coalgebra interval. Therefore Theorem~2.1 of \cite{bm_resolutions} applies and we have the following.

\begin{thm}\label{model structure on Mod(E)}
For each of the categories $\bE$ listed above, the category $\Mod(\bE)$ admits a cofibrantly generated model category structure in which a map of modular operads
$f:\mathcal{P}\longrightarrow \mathcal{Q}$  is a weak equivalence, respectively a fibration, if and only if each component
\[
f(g,n+1):\mathcal{P}(g,n+1)\longrightarrow \mathcal{Q}(g,n+1)
\]
is a weak equivalence, respectively a fibration, in $\bE$, for all $g\geq 0$ and $n\geq 0$.
\end{thm}

In this model structure, fibrant objects are precisely those modular operads that are entrywise fibrant in $\bE$. In particular, since every groupoid is fibrant in the canonical model structure, every modular operad $\mathcal{P}\in\Mod(\Grpd)$ is fibrant.

\begin{prop}\label{prop: cofibrant modular operad}
Let $\mathcal{P}\in\Mod(\Grpd)$ be a modular operad. If the underlying modular operad in sets $\ob(\mathcal{P})\in\Mod(\Set)$ is free, then $\mathcal{P}$ is cofibrant in $\Mod(\Grpd)$.
\end{prop}

\begin{proof}
This follows by the same argument as for operads in groupoids; see \cite[Proposition~6.8]{Horel_profinite_groupoids}.
\end{proof}

\medskip 

A map of coloured operads $f:\mathcal{P}\longrightarrow \mathcal{Q}$ satisfying the \emph{monoidal extension property} (see~\cite[Definition~1.1]{dch_coextension}) induces an adjunction between categories of algebras:
\begin{equation}\label{adj: algebras over different operads}
\begin{tikzcd}
\Alg_{\mathcal{P}}(\bE) \arrow[r, shift left=1, "f_!"] &
\Alg_{\mathcal{Q}}(\bE) \arrow[l, shift left=1, "f^*"]
\end{tikzcd}
\end{equation}
See, for example,~\cite{bm_resolutions} or \cite[Remark~2.26]{dch_coextension}.

A coloured operad is said to be \emph{$\Sigma$-cofibrant} if it is cofibrant as a coloured sequence, that is, cofibrant in the underlying model category of $\mathfrak{C}$-coloured symmetric sequences in $\bE$ (cf.~\cite[Definition~4.7]{hry_shrink}). When $\mathcal{P}$ is $\Sigma$-cofibrant, any sufficiently well-behaved symmetric monoidal Quillen adjunction
\[
\begin{tikzcd}
\bE \arrow[r, shift left=1] & \mathscr{D} \arrow[l, shift left=1]
\end{tikzcd}
\]
lifts to a Quillen adjunction between categories of algebras:
\begin{equation}\label{adj: algebra change of base}
\begin{tikzcd}
\Alg_{\mathcal{P}}(\bE) \arrow[r, shift left=1] &
\Alg_{\mathcal{P}}(\mathscr{D}) \arrow[l, shift left=1]
\end{tikzcd}
\end{equation}
Moreover, if the base adjunction is a Quillen equivalence, then so is the induced adjunction~\eqref{adj: algebra change of base}; see \cite[Theorem~5.8]{hry_shrink}.

These results make $\Sigma$-cofibrant operads especially well behaved from a homotopical perspective. In particular, the operad governing modular operads has this property.

\begin{lemma}\label{lemma: MO is Sigma cofibrant}
The coloured operad $\mathcal{MOp}$ governing modular operads is $\Sigma$-cofibrant.
\end{lemma}

\begin{proof}
It suffices to show that the symmetric group $\Sigma_n$ acts freely on each set
\[
\mathcal{MOp}((g,p);(g_1,k_1),\ldots,(g_n,k_n)).
\]
This action is by permutation of the vertex labelling. Since the operations of $\mathcal{MOp}$ are represented by strict isomorphism classes of connected, genus-graded graphs with labelled vertices and non-empty boundary, a non-trivial permutation of the vertex labels cannot fix such a class. Thus the $\Sigma_n$-action is free. Hence $\mathcal{MOp}$ is $\Sigma$-cofibrant. For a similar argument, see~\cite[Proposition~4.13]{hry_shrink}.
\end{proof}

\begin{remark}
The restriction to operations with non-empty boundary is important for the $\Sigma$-cofibrancy argument. The boundary labelling, together with the strict isomorphism convention for labelled graphs, prevents a non-trivial permutation of the vertex labels from fixing an operation of $\mathcal{MOp}$. If one allows arity-zero operations, equivalently graphs with empty boundary, this rigidity can fail. Such graphs may have non-trivial automorphisms which compensate for a permutation of the vertex labels, and the corresponding symmetric group action need not be free. Thus the coloured operad governing modular operads with arity-zero operations is generally not $\Sigma$-cofibrant. See \cite[Example~6.3]{hry_shrink} or \cite[Figure~2]{deshmukh2022homotopical} for concrete examples.
\end{remark}

\begin{cor}\label{lemma: B is homotopically fully faithful}
The adjunction
\[
\begin{tikzcd}
\Mod(\sSet) \arrow[r, shift left=1, "\Pi_{1}"] & \Mod(\Grpd) \arrow[l, shift left=1, "\cs"]
\end{tikzcd}
\]
is a Quillen adjunction, and the right adjoint $\cs$ is homotopically fully faithful. In particular, for all $\mathcal{P},\mathcal{Q}\in\Mod(\Grpd)$, the induced map
\[
\mathbb{R}\Map_{\Mod(\Grpd)}(\mathcal{P},\mathcal{Q})
\longrightarrow
\mathbb{R}\Map_{\Mod(\sSet)}(\cs\mathcal{P},\cs\mathcal{Q})
\]
is a weak equivalence of spaces.
\end{cor}

\begin{proof}
By Lemma~\ref{lemma: MO is Sigma cofibrant}, the coloured operad $\mathcal{MOp}$ is $\Sigma$-cofibrant. The adjunction
\[
\begin{tikzcd}
\sSet \arrow[r, shift left=1, "\Pi_{1}"] &
\Grpd \arrow[l, shift left=1, "\cs"]
\end{tikzcd}
\]
is a Cartesian monoidal Quillen adjunction. Hence, by \cite[Theorem~5.8]{hry_shrink}, it lifts to a Quillen adjunction between categories of $\mathcal{MOp}$ algebras, equivalently between $\Mod(\sSet)$ and $\Mod(\Grpd)$.

It remains to check homotopical full faithfulness of the right adjoint. Since weak equivalences are defined entrywise and every groupoid is fibrant, every object of $\Mod(\Grpd)$ is fibrant. For each groupoid $G$, the counit $\Pi_{1}\cs G\longrightarrow G$ is an equivalence of groupoids. Therefore, for every $\mathcal{P}\in\Mod(\Grpd)$, the counit $\Pi_{1}\cs\mathcal{P}\longrightarrow \mathcal{P}$ is an entrywise weak equivalence. Since the derived counit is a weak equivalence, it follows that $\cs$ is homotopically fully faithful. The statement on derived mapping spaces follows.
\end{proof}

\subsection{Truncation of Modular Operads}\label{subsec: truncations}

We now define genus truncations of modular operads by restricting the coloured operad $\mathcal{MOp}$ governing modular operads. Recall that our colours are
pairs $(h,p)$, where $h$ is the genus and $p+1$ is the number of boundary components.

For $g\geq 0$, let $\mathcal{MOp}_{\leq g}$ be the coloured suboperad of $\mathcal{MOp}$ whose colours are the pairs \[(h,p)\qquad 0\leq h\leq g,\quad p\geq 0.
\] Its operations are those operations of $\mathcal{MOp}$ whose output colour has genus at most $g$. Thus
\[
\mathcal{MOp}_{\leq g}\bigl((h,p);(h_1,p_1),\ldots,(h_n,p_n)\bigr)
\]
is the set of strict isomorphism classes of connected genus-graded, labelled graphs $G=(G,\lambda,\ell,\epsilon)$ with $n$ vertices and $p+1$ legs, satisfying:
\begin{itemize}
    \item the vertex labelled by $i$ has valence $|\nb(v_{\lambda(i)})|=p_i+1$ and genus $\epsilon(v_{\lambda(i)})=h_i$;
    \item the total genus of $G$ is $h\leq g$, equivalently
    \[
    h=\beta_1(G)+\sum_{i=1}^n h_i.
    \]
\end{itemize}
Graph substitution preserves the total output genus, and hence $\mathcal{MOp}_{\leq g}$ is indeed a coloured suboperad of $\mathcal{MOp}$.

A \emph{genus-$g$ truncated modular operad} in $\bE$ is an algebra over $\mathcal{MOp}_{\leq g}$. We write
\[
\Mod_{\leq g}(\bE):=\Alg_{\mathcal{MOp}_{\leq g}}(\bE)
\]
for the category of genus-$g$ truncated modular operads. Concretely, an object of $\Mod_{\leq g}(\bE)$ consists of the entries $\mathcal P(h,p+1)$ for $0\leq h\leq g$ and $p\geq 0$, together with all composition and contraction operations whose output has genus at most $g$. Thus compositions with $h_1+h_2\leq g$ are retained, and contractions $\xi_i^j:\mathcal P(h,p+1)\to \mathcal P(h+1,p-1)$ are retained only when $h+1\leq g$.

The inclusion of coloured operads $\iota_g:\mathcal{MOp}_{\leq g}\hookrightarrow \mathcal{MOp}$
induces a restriction functor
\[
\trun{g}^*:=\iota_g^*:\Mod(\bE)\longrightarrow \Mod_{\leq g}(\bE),
\]
which we call the \emph{genus-$g$ truncation functor}. This functor admits both a left adjoint and a right adjoint:
\begin{equation}\label{eq: truncation adjunction}
\begin{tikzcd}
\Mod(\bE) \arrow[r, swap, "\trun{g}^*"] &
\Mod_{\leq g}(\bE)
\arrow[l, bend left=50, "(\trun{g})_*"]
\arrow[l, bend right=50, swap, "(\trun{g})_!"] .
\end{tikzcd}
\end{equation}
The left adjoint $(\trun{g})_!$ generalizes the notion of the \emph{modular envelope}; see~\cite{Ward}, \cite[Appendix~A]{dch_coextension}, and~\cite[§4]{GNPR} for related constructions.


\begin{prop}\label{prop: truncation is Quillen functor}
Let $\bE$ be a cofibrantly generated Cartesian monoidal model category with cofibrant unit, a symmetric monoidal fibrant replacement functor, and a cocommutative coalgebra interval. Then both $\Mod(\bE)$ and $\Mod_{\leq g}(\bE)$ admit cofibrantly generated model category structures in which weak equivalences and fibrations are defined entrywise. More precisely, a map $f:\mathcal P\to \mathcal Q$ in $\Mod(\bE)$ is a weak equivalence, respectively a fibration, if and only if each map
\[
f(h,m+1):\mathcal P(h,m+1)\longrightarrow \mathcal Q(h,m+1)
\]
is a weak equivalence, respectively a fibration, in $\bE$, for all $h\geq 0$ and $m\geq 0$. Similarly, in $\Mod_{\leq g}(\bE)$ this condition is imposed only for $0\leq h\leq g$.

Moreover:
\begin{itemize}
    \item the truncation functor
    $\trun{g}^*:\Mod(\bE)\to\Mod_{\leq g}(\bE)$ preserves weak equivalences and
    fibrations;
    \item both adjoints $(\trun{g})_!$ and $(\trun{g})_*$ are fully faithful;
    \item there are canonical isomorphisms
    \[
    \trun{g}^* \circ (\trun{g})_! \cong \id_{\Mod_{\leq g}(\bE)}
    \quad \text{and} \quad
    \trun{g}^* \circ (\trun{g})_* \cong \id_{\Mod_{\leq g}(\bE)}.
    \]
\end{itemize}
\end{prop}

\begin{proof}
The functor $\trun{g}^*$ is restriction along the inclusion of coloured operads
\[
\iota_g:\mathcal{MOp}_{\leq g}\hookrightarrow \mathcal{MOp}.
\]
Since weak equivalences and fibrations are entrywise, restriction clearly preserves both weak equivalences and fibrations.

The left and right adjoints are the left and right algebraic extensions along $\iota_g$. Because $\mathcal{MOp}_{\leq g}$ is a full coloured suboperad of $\mathcal{MOp}$ on the colours of genus at most $g$, extending a truncated modular operad and then restricting back recovers exactly the original truncated
object. Equivalently, the counits
\[
\trun{g}^*(\trun{g})_! \longrightarrow \id_{\Mod_{\leq g}(\bE)}
\qquad\text{and}\qquad
\trun{g}^*(\trun{g})_* \longrightarrow \id_{\Mod_{\leq g}(\bE)}
\]
are isomorphisms. Hence both adjoints are fully faithful.
\end{proof}

As in the modular case, cyclic operads are algebras over a coloured operad. In fact, the coloured operad governing cyclic operads may be identified with the
genus-zero truncation
\[
\mathcal C \cong \mathcal{MOp}_{\leq 0}.
\]
Its operations are represented by labelled unrooted trees, equivalently by the genus-zero graphs
\[
T\in
\mathcal{MOp}_{\leq 0}\bigl((0,p);(0,k_1),\ldots,(0,k_n)\bigr).
\]
Here, as throughout, the colour $(0,p)$ corresponds to $p+1$ boundary components. The following identification can be found in \cite[Example~A3]{dch_coextension} or \cite[Section~1.6.4]{lukacs2010cyclic}.

\begin{prop}\label{prop: cyc are genus 0 modular operads}
There is a natural isomorphism of categories
\[
\Alg_{\mathcal C}(\bE)\cong \Cyc(\bE).
\]
%Equivalently, genus-zero truncated modular operads are cyclic operads.
\end{prop}

As a consequence of Proposition~\ref{prop: cyc are genus 0 modular operads}, Theorem~2.1 of~\cite{bm_resolutions} applies to cyclic operads as well. Thus $\Cyc(\bE)$ admits a cofibrantly generated model category structure in which a map of cyclic operads $f:\mathcal P\to\mathcal Q$ is a weak equivalence,
respectively a fibration, if and only if each component $f(n):\mathcal P(n)\longrightarrow\mathcal Q(n)$ is a weak equivalence, respectively a fibration, in $\bE$, for all $n\geq 1$. Moreover, under the identification $\Mod_{\leq 0}(\bE)\cong\Cyc(\bE)$, the genus-zero truncation functor becomes
\[
\trun{0}^*:\Mod(\bE)\longrightarrow \Cyc(\bE).
\]
It preserves weak equivalences and fibrations, and hence is the right adjoint in a Quillen adjunction.


\section{Homotopical properties of the nerve}\label{sec: nerve is homotopically fully faithful}

Recall that $\dMod(\bE)$ denotes the category of modular dendroidal objects, that is, functors
\[
X:\bM^{\mathrm{op}}\longrightarrow \bE .
\]
This category carries several useful model structures. We write $\dMod(\bE)_p$ for the projective model structure, in which weak equivalences and fibrations are defined entrywise \cite[Theorem 11.6.1]{hirsch}.

The graphical category $\bM$ is also a generalized Reedy category \cite[Theorem~2.22]{hry1}. Hence, for any suitable model category $\bE$, $\dMod(\bE)$ admits a Reedy model structure; see \cite[Theorem~1.6]{bm_reedy}. We write $\dMod(\bE)_R$ for this Reedy model structure.

The identity functor on the underlying category $\dMod(\bE)$ induces a Quillen equivalence
\begin{equation}\label{identity qe}
\begin{tikzcd}
\dMod(\bE)_p \arrow[r, shift left=1] &
\dMod(\bE)_R \arrow[l, shift left=1].
\end{tikzcd}
\end{equation}
This follows from \cite[Lemma 3.9]{hry1}.

Modular $\infty$-operads are obtained as the fibrant objects in a left Bousfield localisation of the Reedy model structure. We use left Bousfield localisation in the standard sense: if $\mathcal{S}$ is a set of maps in a model category $\mathcal{M}$, then the localised model structure $\mathcal{L}_{\mathcal{S}}\mathcal{M}$ has the same cofibrations as $\mathcal{M}$, and its fibrant objects are precisely the fibrant objects $Z$ of $\mathcal{M}$ such that, for every map $A\to B$ in $\mathcal{S}$, the induced map of derived mapping spaces
\[
\mathbb{R}\Map(B,Z)\longrightarrow \mathbb{R}\Map(A,Z)
\]
is a weak equivalence. Its weak equivalences are the corresponding $\mathcal{S}$-local equivalences. We refer to \cite[Chapters~3--4]{hirsch} for the general theory.

In Definition~\ref{def: infty modular}, a modular dendroidal object $X\in\dMod(\bE)$ is called a modular $\infty$-operad if it is reduced and
satisfies the Segal condition. Here reduced means that $X_{\updownarrow}=*$.

Any modular dendroidal object can be reduced via the pushout
\[
\begin{tikzcd}
X_{\updownarrow}\times \bM[\updownarrow]\arrow[d]\arrow[r] &
\bM[\updownarrow]\arrow[d]\\
X\arrow[r] & X_* .
\end{tikzcd}
\]
This defines a functor $(-)_*:\dMod(\bE)\longrightarrow \dMod(\bE)_*$ which is left adjoint to the inclusion of reduced modular dendroidal objects into all
modular dendroidal objects.

Both the projective and Reedy model structures restrict to $\dMod(\bE)_*$. We write $\dMod(\bE)_{*,p}$ and $\dMod(\bE)_{*,R}$ for the reduced projective and reduced Reedy model structures, respectively.

For each graph $G\in\bM$, let $\Sc[G]_*\longrightarrow \bM[G]_*$ denote the reduced Segal core inclusion. Let $S$ be the set of these maps:
\[
S=\{\Sc[G]_*\longrightarrow \bM[G]_*\mid G\in\bM\}.
\] We denote the left Bousfield localisation of $\dMod(\bE)_{*,R}$ with respect to $S$ by $\mathcal{L}_S\dMod(\bE)_{*,R}$. Its fibrant objects are the Reedy
fibrant reduced modular dendroidal objects $X$ for which, for every graph $G$, the induced map
\[
\mathbb R\Map(\bM[G]_*,X)\longrightarrow \mathbb R\Map(\Sc[G]_*,X)
\]
is a weak equivalence. This is exactly the Segal condition.

We can now make precise what we mean by the $\infty$-category of modular $\infty$-operads.


\begin{definition}\label{def: infinity cat of infinity modular operads}
The $\infty$-category of modular $\infty$-operads, denoted $\Mod_{\infty}(\bE)$, is the $\infty$-category obtained from the full relative subcategory of $\dMod(\bE)_{*,p}$ spanned by the modular $\infty$-operads, with weak equivalences given by the entrywise weak equivalences.
\end{definition}

The main result of this appendix is that the modular dendroidal nerve is homotopically fully faithful. In other words, strict modular operads embed into modular $\infty$-operads without changing their derived mapping spaces.

\begin{thm}\label{thm: nerve is homotopically fully faithful}
Let $\bE=\sSet$ or $\Grpd$. For any two modular operads $\calP,\calQ\in\Mod(\bE)$, the nerve induces a weak equivalence of derived mapping spaces
\[
\mathbb{R}\Map_{\Mod(\bE)}(\calP,\calQ)
\longrightarrow
\mathbb{R}\Map_{\Mod_{\infty}(\bE)}(\nerve\calP,\nerve\calQ).
\]
Equivalently, the modular dendroidal nerve
\[
\nerve:\Mod(\bE)\longrightarrow \Mod_{\infty}(\bE)
\]
is homotopically fully faithful.
\end{thm}

Before proving Theorem~\ref{thm: nerve is homotopically fully faithful}, we record a comparison between decorated representables and nerves of free modular
operads. Let $G=(G,\lambda,\ell,\epsilon)$ be a labelled graph with $k$ vertices, written as $v_{\lambda(1)},\ldots,v_{\lambda(k)}$ according to the vertex labelling. For each $1\leq i\leq k$, write $\epsilon(v_{\lambda(i)})=g_i$ and $|\nb(v_{\lambda(i)})|=n_i+1$.

Define a genus-graded sequence in $\Set$, denoted $X_G$, by:
\[
X_G(h,n+1)
=
\coprod_{\substack{1\leq i\leq k\\ g_i=h,\; n_i=n}}
*,
\]
and denote the generator corresponding to the $i$th labelled vertex by $x_i^{(g_i,n_i+1)}$. Thus $X_G$ has one generator for each labelled vertex of $G$, placed in the entry determined by the genus and valence of that vertex.

We also regard $X_G$ as a discrete decoration of the representable modular dendroidal set $\bM[G]$: a map into $G$ carries the vertex data of $G$, and each vertex $v_{\lambda(i)}$ is decorated by the corresponding generator $x_i^{(g_i,n_i+1)}$. We write
\[
Z_G:=\bM[G]\times X_G
\]
for this decorated representable modular dendroidal object. Concretely, $Z_G$ is the modular dendroidal object whose value on a graph $H$ consists of a map $H\to G$ in $\bM$, together with the induced decoration of the vertices of $H$ by the generators of $X_G$.

Let $\mathcal F(X_G)$ be the free modular operad generated by $X_G$. The graph $G$ determines a distinguished operation in $\mathcal F(X_G)$, obtained by composing the generators $x_i^{(g_i,n_i+1)}$ according to the internal edges of $G$.

\begin{lemma}\label{lemma: nerve of free operad}
Let $\bE=\Grpd$ or $\sSet$, and let $G\in\bM$. With $X_G$ and $Z_G$ as above, the decorated representable $Z_G$ and the nerve of the free modular operad $\nerve(\mathcal F(X_G))$ are weakly equivalent in the localised model structure $\mathcal L_S\dMod(\bE)_{*,R}$.
\end{lemma}

\begin{proof}
The argument is the modular analogue of \cite[Proposition~4.4]{bh_group_actions_segal}.

Let $\Sc_G(X_G)$ denote the Segal core of $G$ decorated by the vertex generators $x_i^{(g_i,n_i+1)}$. Thus $\Sc_G(X_G)$ is obtained from $\Sc[G]_*$ by labelling the corolla corresponding to the vertex $v_{\lambda(i)}$ by the generator $x_i^{(g_i,n_i+1)}$. There is a canonical map $\Sc_G(X_G)\longrightarrow Z_G.$  Since $\Sc[G]_*\to \bM[G]_*$ is one of the Segal core inclusions, this map is a weak equivalence in the localised model structure $\mathcal L_S\dMod(\bE)_{*,R}$.

We now compare $\Sc_G(X_G)$ with $\nerve(\mathcal F(X_G))$. The generators $x_i^{(g_i,n_i+1)}$ define maps from the vertex corollas of $G$ into
$\nerve(\mathcal F(X_G))$, and these assemble to a map
\[
\Sc_G(X_G)\longrightarrow \nerve(\mathcal F(X_G)).
\] Since $\nerve(\mathcal F(X_G))$ is the nerve of a strict modular operad, it satisfies the Segal condition strictly. Therefore maps from a Segal object into $\nerve(\mathcal F(X_G))$ are determined by their restrictions to the vertex corollas. By the universal property of the free modular operad
$\mathcal F(X_G)$, this is exactly the same data as a choice of images of the generators $x_i^{(g_i,n_i+1)}$. It follows that the map $\Sc_G(X_G)\longrightarrow \nerve(\mathcal F(X_G))$ exhibits $\nerve(\mathcal F(X_G))$ as a local replacement of $\Sc_G(X_G)$.

Combining this with the local weak equivalence $\Sc_G(X_G)\to Z_G$, we obtain a weak equivalence
\[
Z_G\longrightarrow \nerve(\mathcal F(X_G))
\]
in $\mathcal L_S\dMod(\bE)_{*,R}$.
\end{proof}

The modular dendroidal nerve $\nerve:\Mod(\bE)\longrightarrow \dMod(\bE)_*$  admits a left adjoint, denoted
\[
\tau:\dMod(\bE)_*\longrightarrow \Mod(\bE).
\]
The left adjoint is determined by its values on reduced representables: $\tau(\bM[G]_*)$ is the free modular operad generated by the graph $G$. More explicitly, if a vertex $v\in V_G$ has genus $\epsilon(v)$ and valence $|\nb(v)|$, then the corresponding generator lies in the entry $(\epsilon(v),|\nb(v)|)$.

\begin{prop}\label{quillen pair projective}
Let $\bE=\Grpd$ or $\sSet$. The adjoint pair
\[
\begin{tikzcd}
\dMod(\bE)_{*,p} \arrow[r, shift left=1, "\tau"] &
\Mod(\bE) \arrow[l, shift left=1, "\nerve"]
\end{tikzcd}
\]
is a Quillen pair.
\end{prop}

\begin{proof}
It is enough to show that the right adjoint $\nerve$ preserves fibrations and
trivial fibrations. For a modular operad $\calP$, the nerve is given on corollas
by
\[
(\nerve\calP)_{\corolla_{(g,n+1)}}=\calP(g,n+1).
\]
Thus, if $f:\calP\to\calQ$ is a fibration, respectively a trivial fibration,
in $\Mod(\bE)$, then each map
\[
f(g,n+1):\calP(g,n+1)\longrightarrow \calQ(g,n+1)
\]
is a fibration, respectively a trivial fibration, in $\bE$.

Now let $G\in\bM$. Since the nerve of a modular operad satisfies the Segal condition strictly, we have
\[
(\nerve\calP)_G = \prod_{v\in V_G} (\nerve\calP)_{\corolla_{(\epsilon(v),|\nb(v)|)}}.
\]
Therefore $(\nerve f)_G$ is the finite product of the maps
\[
f\bigl(\epsilon(v),|\nb(v)|\bigr):
\calP\bigl(\epsilon(v),|\nb(v)|\bigr)
\longrightarrow
\calQ\bigl(\epsilon(v),|\nb(v)|\bigr),
\] where $v$ ranges over the vertices of $G$.  Finite products preserve fibrations and trivial fibrations in both $\sSet$ and $\Grpd$. Hence $(\nerve f)_G$ is a fibration, respectively a trivial fibration, for every graph $G$. 

Since fibrations and weak equivalences in the projective model structure on $\dMod(\bE)_{*,p}$ are defined entrywise, $\nerve f$ is a fibration, respectively a trivial fibration. It follows that $\nerve$ is a right Quillen functor, and the adjunction is a Quillen pair.
\end{proof}

\medskip

We now prove the key input for homotopical full faithfulness of the nerve. The idea of the next lemma is that cofibrant modular dendroidal objects are built, up to homotopy colimit, from reduced representables. On reduced representables, the comparison with the nerve is controlled by Lemma~\ref{lemma: nerve of free operad}: the localisation identifies a decorated representable with the nerve of the free modular operad generated by its vertex decorations. The general case then follows because the relevant constructions preserve the homotopy colimits used to build cofibrant diagrams.

\begin{lemma}\label{derived unit}
Let $\bE=\sSet$ or $\Grpd$. If $X\in\dMod(\bE)_p$ is cofibrant, then the derived unit
\[
X_* \longrightarrow \nerve\bigl(\tau(X_*)_f\bigr)
\]
is a local equivalence in $\mathcal{L}_{S}\dMod(\bE)_{*,p}$.
\end{lemma}


\begin{proof}
We give the argument for $\bE=\sSet$. The proof for $\Grpd$ is analogous.

Since $X$ is cofibrant in the projective model structure, we can use standard simplicial resolution of a cofibrant diagrams to write $X_*$, up to weak equivalence, as a homotopy colimit of coproducts of decorated reduced representables:
\[
X_* \simeq \operatorname*{hocolim}_{\Delta^{op}}\coprod_i Z_{G_i}.
\]
Here each $Z_{G_i}$ is of the form considered in Lemma~\ref{lemma: nerve of free operad}, with vertex-decoration sequence $X_{G_i}$.

Since $\tau$ is a left Quillen functor by Proposition~\ref{quillen pair projective}, its left derived functor preserves homotopy colimits. Thus,
\[
\tau(X_*) \simeq \tau\left( \operatorname*{hocolim}_{\Delta^{op}}\coprod_i Z_{G_i} \right) \simeq \operatorname*{hocolim}_{\Delta^{op}}\coprod_i \tau(Z_{G_i}).
\]
By construction of the left adjoint $\tau$, each $\tau(Z_{G_i})$ is the free modular operad generated by the corresponding vertex decorations: $\tau(Z_{G_i})\cong \mathcal F(X_{G_i}).$ Putting this together, we have \[\tau(X_*) \simeq \operatorname*{hocolim}_{\Delta^{op}}\coprod_i \mathcal F(X_{G_i}).\]

We now compare the nerve of this homotopy colimit with the homotopy colimit of the nerves. After taking a fibrant replacement in $\Mod(\sSet)$ and applying
the nerve, we get, up to local equivalence,
\[
\nerve(\tau(X_*)_f) \simeq \nerve\left( \operatorname*{hocolim}_{\Delta^{op}}\coprod_i \mathcal F(X_{G_i})_{f} \right).
\] The nerve of a strict modular operad satisfies the Segal condition strictly and, thus its value on a graph is a finite product of its values on the vertex corollas. Since homotopy colimits over $\Delta^{op}$ commute with finite homotopy products in $\sSet$, it follows that
\[
\nerve\left(\operatorname*{hocolim}_{\Delta^{op}}\coprod_i \mathcal F(X_{G_i})_{f} \right) \simeq \operatorname*{hocolim}_{\Delta^{op}}\coprod_i \nerve\bigl(\mathcal F(X_{G_i})_{f}\bigr).
\]

On the other hand, by Lemma~\ref{lemma: nerve of free operad}, each decorated representable $Z_{G_i}$ is weakly equivalent, in the localised model structure, to $\nerve\bigl(\mathcal F(X_{G_i})\bigr).$ Fibrant replacement in $\Mod(\sSet)$ does not change the local equivalence class of its nerve, so we may equally use
$\nerve(\mathcal F(X_{G_i})_f)$. Therefore
\[
\operatorname*{hocolim}_{\Delta^{op}}\coprod_i Z_{G_i} \longrightarrow \operatorname*{hocolim}_{\Delta^{op}}\coprod_i \nerve\bigl(\mathcal F(X_{G_i})_{f}\bigr)
\]
is a local equivalence. Combining the displayed equivalences above gives a local equivalence
\[
X_* \longrightarrow \nerve(\tau(X_*)_f)
\]
in $\mathcal{L}_{S}\dMod(\sSet)_{*,p}$, as required.
\end{proof}



\begin{proof}[Proof of Theorem~\ref{thm: nerve is homotopically fully faithful}]
By Proposition~\ref{quillen pair projective}, the adjunction
\[
\begin{tikzcd}
\dMod(\bE)_{*,p} \arrow[r, shift left=1, "\tau"] &
\Mod(\bE) \arrow[l, shift left=1, "\nerve"]
\end{tikzcd}
\]
is a Quillen adjunction. Since the nerve of a strict modular operad satisfies the Segal condition strictly, the right adjoint $\nerve$ sends fibrant modular operads to fibrant objects in the localised model structure $\mathcal{L}_{S}\dMod(\bE)_{*,p}$. Hence the same adjunction descends to a Quillen adjunction
\[
\begin{tikzcd}
\mathcal{L}_{S}\dMod(\bE)_{*,p} \arrow[r, shift left=1, "\tau"] &
\Mod(\bE) \arrow[l, shift left=1, "\nerve"] .
\end{tikzcd}
\]

It remains to check that the right adjoint is homotopically fully faithful. For a strict modular operad $\calP$, the counit $\tau\nerve(\calP)\longrightarrow \calP$ is an isomorphism. Indeed, the nerve records the values of $\calP$ on corollas, and the left adjoint $\tau$ reconstructs the strict modular operad generated by these operations modulo the compositions already encoded by the nerve.

Now let $\calP$ be fibrant in $\Mod(\bE)$. Then the derived counit $\mathbb{L}\tau\,\mathbb{R}\nerve(\calP)\longrightarrow \calP$ is represented by the ordinary counit above, and is therefore a weak equivalence. It follows that the right derived functor of $\nerve$ is fully faithful. Equivalently, for all $\calP,\calQ\in\Mod(\bE)$, the induced map on derived mapping spaces
\[
\mathbb{R}\Map_{\Mod(\bE)}(\calP,\calQ)
\longrightarrow
\mathbb{R}\Map_{\Mod_{\infty}(\bE)}(\nerve\calP,\nerve\calQ)
\]
is a weak equivalence.
\end{proof}














\section{The complex of markings over pants decompositions}\label{app: complex of markings}
In this appendix, we prove that the complex of markings $\calM(\surf)$ introduced in Section~\ref{sec: complex of markings} is connected and simply connected. This is the complex used in the proof of Theorem~\ref{thm: presentation for maps out of bS}.

The $2$-dimensional CW complex $\calM(\surf)$ is closely related to the complex $\Mmax(\surf)$ of \cite{bk_marked_surfaces}. The main difference is that $\calM(\surf)$ only retains markings over pants decompositions: disc and cylinder components are not allowed. Following the strategy of \cite{bk_marked_surfaces,complex_of_pants_decomp}, our proof proceeds by comparing $\calM(\surf)$ with a simply connected complex via a forgetful map and applying two general comparison results.

The first lemma allows us to deduce simple connectivity of the target from that of the source.

\begin{lemma}\label{lemma: simply-connect target from domain}
    Let $A$ be a simply connected CW complex, and let $f:A\to B$ be a continuous map such that every loop in $B$ lifts to a loop in $A$. Then $B$ is simply connected. \qed
\end{lemma}

The following proposition gives a criterion for deducing simple connectivity of the source from that of the target.

\begin{prop}\label{prop: simply-connect domain from target}
    Let $A$ and $B$ be $2$-dimensional CW complexes, and let $\pi:A^{[1]}\to B^{[1]}$ be a surjective map of their $1$-skeleta. Suppose that
    \begin{enumerate}
        \item $B$ is connected and simply connected.
        \item For every vertex $b\in B$, the fiber $\pi^{-1}(b)$ is connected, and every loop contained in $\pi^{-1}(b)$ is contractible in $A$.
        \item Let $b_1\xra{e} b_2$ be an edge in $B$, and let $a_1'\xra{e'}a_2'$ and $a_1''\xra{e''}a_2''$ be two lifts to $A$. Then there exist paths $a_1'\xra{e_1}a_1''$ in $\pi^{-1}(b_1)$ and $a_2'\xra{e_2}a_2''$ in $\pi^{-1}(b_2)$ such that the square
            \[
            \begin{tikzcd}
            	{a_1'} & {a_2'} \\
            	{a_1''} & {a_2''}
            	\arrow["{e'}", from=1-1, to=1-2]
            	\arrow["{e_1}"', from=1-1, to=2-1]
            	\arrow["{e_2}", from=1-2, to=2-2]
            	\arrow["{e''}", from=2-1, to=2-2]
            \end{tikzcd}
            \]
        is contractible in $A$.
        \item For every $2$-cell $X$ in $B$, its boundary $\partial X$ can be lifted to a contractible loop in $A$.
    \end{enumerate}
    Then $A$ is connected and simply connected. \qed
\end{prop}

The proofs of the preceding two results are standard. Proposition~\ref{prop: simply-connect domain from target} is established by a cellular lifting argument: one projects a loop in $A$ to $B$ and lifts a nullhomotopy back to $A$ cell by cell. Conditions~\emph{(2)} and~\emph{(3)} are used to compare different choices of lifts, and condition~\emph{(4)} is applied to lift the $2$-cells of $B$. We omit the details.

We now recall the pants complex of Hatcher--Lochak--Schneps \cite{hls} in a form suited to our setting. Let $S=(S,\delta)$ be a labelled surface of type $(g,n+1)$, with $n\ge 0$. Recall that a \emph{pants decomposition} of $S$ is a cut system $C=\{c_1,\ldots,c_k\}$ of pairwise disjoint simple loops on $S$ such that
    \[\overline{S\setminus \bigcup_{r=1}^k c_r}\;=\;\coprod_j S_j\]
is a disjoint union of subsurfaces $S_j$, each of which is of type $(0,3)$. The \emph{pants complex} of $S$, denoted $\pantscpx(S)$, is the $2$-dimensional CW complex whose vertices are isotopy classes of pants decompositions on $S$, whose edges are generated by $A$-moves and $S$-moves between pants decompositions, and whose $2$-cells are attached along loops of type $(3A)$, $(5A)$, $(3S)$, $(6AS)$, and $(C)$, in the sense of \cite[Section~2]{hls}.

By \cite[Theorem~2]{hls}, the complex $\pantscpx(S)$ is connected and simply connected.

Recall from Definition~\ref{def: complex of markings} that the complex of marked pants decompositions on $S$, denoted $\markS$, has vertices given by marked pants decompositions $(C,[m])$ on $S$. The $1$-cells are generated by $T$-, $A$-, $B$-, and $S$-moves. The $2$-cells are slight modifications of those appearing in \cite{bk_marked_surfaces}. Concretely, we include the following families of $2$-cells:
\begin{enumerate}
    \item \textbf{Commutativity of disjoint moves:} if $E_1$ and $E_2$ are moves supported in disjoint subsurfaces of $S$, then
                \[
                E_1E_2 = E_2E_1
                \]
                \cite[Section~4.7]{bk_marked_surfaces}.
          
    \item \textbf{Boundary twist relations:} if $(C,[m])\overset{E}{\rightsquigarrow}(C',[m'])$ is a move supported on a subsurface $S'$ and $a\in C$ is a curve on the boundary of $S'$, then if $E=B_{a,b}$
                \[E \; T_a= T_b\; E\]
            otherwise, we have
                \[E\; T_a=T_a \; E.\]
    
    \item \textbf{Pentagon relation:} whenever there exists a subsurface $S'\subset S$ of type $(0,5)$ together with a diffeomorphism from $S'$ to the standard surface of Figure~\ref{fig:pentagon} carrying the relevant curves and markings to those represented in the figure, we attach a $2$-cell along the boundary loop of the pentagon diagram.
        \begin{figure}[ht]
            \centering
            \includegraphics[width=0.5\linewidth]{Figures/pentagon.pdf}
            \caption{Pentagon relation}
            \label{fig:pentagon}
        \end{figure}
    % \cite[Figure 30, Appendix~C]{bk_marked_surfaces}.

    \item \textbf{Hexagon relations:} whenever there exists a subsurface $S'\subset S$ of type $(0,4)$ together with a diffeomorphism from $S'$ to the standard surface $S_{0,4}$ of Figure~\ref{fig:hexagons} carrying the relevant curves and markings to those represented in the figure, we attach $2$-cells along the two hexagonal loops of Figure~\ref{fig:hexagons}.

    % \begin{figure}[ht]
    %     \includegraphics[width=0.55\linewidth]{Figures/hexagonI.pdf}\hfill\hfill \\
    %     \hfill\includegraphics[width=0.55\linewidth]{Figures/hexagonII.pdf}
    %     \caption{Hexagon II}
    %     \label{fig:hexagonI}
    % \end{figure}
    \begin{figure}[ht]
        % Top-left
        \makebox[\textwidth][l]{%
            \includegraphics[width=0.735\linewidth]{Figures/hexagonI.pdf}}
        % \vspace{-2.5cm}
        % Bottom-right
        \makebox[\textwidth][r]{%
            \includegraphics[width=0.735\linewidth]{Figures/hexagonII.pdf}}
        \caption{Hexagon relations: Hexagon I (top left) and Hexagon II (bottom right)}
        \label{fig:hexagons}
    \end{figure}
    
    \item \textbf{Self-duality of associativity:}
        \[A_{c,c'}A_{c',c}=\id.\]

    \item \textbf{Twist--braiding relation:}
        \[B_{1,2}B_{2,1}=T_0T_1^{-1}T_2^{-1}.\]

    \item \textbf{Gluing relations for Dehn twists:} 
        \begin{itemize}
            \item if two pairs of pants $S_1$ and $S_2$ are glued along a curve $a$, then
                \[\id_{S_1}\cup_a T_a = T_a\cup_a \id_{S_2}\]
            where the $T$-move is supported on $S_2$ on the left-hand side and on $S_1$ on the right-hand side.
            \item if a pair of pants $S'$ has two boundary components $a$ and $b$ that get identified in $S$, then 
                \[T_a = T_b.\]
        \end{itemize}

    \item \textbf{Genus-one relations for $(g,n)=(1,1)$:} as in Figure~\ref{fig:G2-cells-G1} in Definition~\ref{def: complex of markings}.

    \item \textbf{Genus-one relation for $(g,n)=(1,2)$:} as in Figure~\ref{fig:G2-cell-NS-complex-of-markings} in Definition~\ref{def: complex of markings}.
\end{enumerate}

\medskip

The complex of markings over the pair of pants $P$, will be of particular relevance for us.

\begin{lemma}\label{lemma: marked complex for pants is simply connected}
Let $P$ be a pair of pants. Then $\mathcal{M}(P)$ is connected and simply connected.
\end{lemma}


\begin{proof}
For a pair of pants $P$, a vertex of $\mathcal M(P)$ is simply a marked pants decomposition of the form $(\emptyset,[m]),$ since the only pants decomposition of $P$ is the empty cut system.

Let $\mathcal M^{\max}(P)$ denote the complex of maximal markings on $P$ in the sense of \cite[Section~5]{bk_marked_surfaces}. A vertex of $\mathcal M^{\max}(P)$ is given by an isotopy class of a cut system
\[
C=\{c_1,\ldots,c_k\}
\]
together with a marking, such that cutting $P$ along $C$ decomposes it into pairs of pants, possibly together with cylinder and disc components. In addition, one records a distinguished half-edge at each vertex of the marking graph.

We compare $\mathcal M^{\max}(P)$ with $\mathcal M(P)$ by means of a map
\[
q:\mathcal M^{\max}(P)\longrightarrow \mathcal M(P).
\]
On vertices, the map $q$ forgets the cut system $C$ and the distinguished half-edges, and contracts those edges of the marking graph that are not incident to to a boundary component, thereby producing a vertex $(\emptyset,[m])\in \mathcal{M}(P)$.

On edges, the map $q$ behaves as follows. Moves of type $T$ and $B$ are either contracted to points or taken to the corresponding edges of $\mathcal M(P)$. By contrast, moves of type $F$ forget cuts bounding cylindrical components and therefore become trivial after passing to the empty cut system. Likewise, $A$-moves are supported on a subsurface of type $(0,4)$, and also are contracted to points via $q$. Finally, a $Z$-move changes only the choice of distinguished half-edge, which is not retained in $\mathcal M(P)$. Thus $A$-, $F$-, and $Z$-edges are all sent to vertices.


This defines a cellular map on the $1$-skeleton of $\mathcal M^{\max}(P)$. Moreover, each defining $2$-cell of $\mathcal M^{\max}(P)$ maps either to the corresponding $2$-cell of $\mathcal M(P)$ or to a degenerate loop. Hence $q$ extends to a cellular map
\[
q:\mathcal M^{\max}(P)\to \mathcal M(P).
\]


We next show that every loop in $\mathcal M(P)$ lifts to a loop in $\mathcal M^{\max}(P)$. Given a cellular loop in $\mathcal{M}(P)$, each vertex in the loop is a marking of the form $(\emptyset,[m])$. Any such vertex can be lifted to one in $\mathcal{M}^{\max}(P)$ by choosing a distinguished half-edge for the marking $[m]$. For instance, we may choose the half-edge incident to the boundary component labelled $0$. In this way, all vertices in the loop in $\mathcal{M}(P)$ lift to vertices in $\mathcal{M}^{\max}(P)$.

Since every edge of $\mathcal M(P)$ is a $T$- or $B$-move, each $e_i$ lifts to an edge path in $\mathcal M^{\max}(P)$ consisting of the corresponding $T$- or $B$-move, together with $Z$-moves, if necessary, to adjust the distinguished half-edge before or after performing that move. Concatenating these lifted edge paths yields a loop in $\mathcal M^{\max}(P)$ projecting to the original loop in $\mathcal M(P)$.

By \cite[Theorem~4.9]{bk_marked_surfaces}, the complex $\mathcal M^{\max}(P)$ is connected and simply connected. Lemma~\ref{lemma: simply-connect target from domain} therefore implies that $\mathcal M(P)$ is hence connected and simply connected as well.
\end{proof}


To study the complex of markings $\markS$ for more general surfaces, we compare $\markS$ with $\pantscpx(S)$ via the natural forgetful map $F:\markS\longrightarrow \pantscpx(S)$ defined on vertices by
\[
F(C,[m])=C,
\]
that is, by sending a marked pants decomposition to its underlying pants decomposition.

We use this to prove the main result of this section:

\begin{prop*}[Proposition~\ref{prop:our-complex-of-markings-is-simply-connected}]
    Let $S=(S, \delta, \{x_{i}\}_{i=0}^{n})$ be a parametrised, labelled surface of type $(g,n+1)$, $n\geq 0$. Then $\markS$, the complex of markings over pants decompositions on $S$, is connected and simply connected.
\end{prop*}

\begin{proof}
Let
\[
F:\markS\longrightarrow \pantscpx(S)
\]
be the forgetful map sending a marked pants decomposition $(C,[m])$ to its underlying pants decomposition $C$. By \cite[Theorem~2]{hls}, the pants complex $\pantscpx(S)$ is connected and simply connected. We now verify the hypotheses of Proposition~\ref{prop: simply-connect domain from target} for the map $F$. By construction, $F$ is surjective on vertices and edges.

\emph{Checking condition~$(2)$:} Let $C$ be a vertex of $\pantscpx(S)$, that is, a pants decomposition of $S$. The fibre $F^{-1}(C)$ is the subcomplex of $\markS$ whose vertices are markings with underlying pants decomposition $C$, and whose edges and cells are precisely those built from moves that leave $C$ unchanged. In particular, its $1$-cells are the $T$- and $B$-moves, and its $2$-cells are the relations among these moves that do not alter the underlying pants decomposition, namely commutativity of disjoint moves, the boundary twist relations, the twist--braiding relation, and the Gluing relations for Dehn twists.

We next show that $F^{-1}(C)$ is connected and simply connected. Let $S_1,\dots,S_k$ be the pair-of-pants components obtained as the closures of the connected components of $S\setminus C$. For each curve $c\in C$, choose a basepoint on $c$. These basepoints induce marked points on the corresponding boundary components of the $S_i$, so that each $S_i$ becomes a parametrised, labelled surface of type $(0,3)$ (Definition~\ref{def: labelled surface}). In particular, a marking on each $S_i$ may be viewed as an embedded corolla meeting each boundary component at the chosen marked point.

Using these basepoints, we may glue markings on the $S_i$ along the cut curves: if two boundary components of $S_i$ and $S_j$ arise from the same curve $c\in C$, then we identify them so that the marked points determined by the chosen basepoint of $c$ agree. In this way we obtain a map
\[
\mathcal M(S_1)\times\cdots\times \mathcal M(S_k)\longrightarrow F^{-1}(C),
\]
sending a tuple of markings on the pair-of-pants components to the induced marking on $S$ with underlying pants decomposition $C$. Since each $\mathcal M(S_i)$ is connected and simply connected by Lemma~\ref{lemma: marked complex for pants is simply connected}, the product
    \begin{align}\label{eq:gluing-map-for-markings}
        \mathcal M(S_1)\times\cdots\times \mathcal M(S_k)
    \end{align}
is also connected and simply connected.

The gluing map is surjective on vertices: any marking on $S$ with underlying pants decomposition $C$ can be isotoped so that it intersects each curve $c\in C$ at its chosen marked point. Such a marking then restricts, after cutting along $C$, to markings on the pair-of-pants components $S_1,\dots,S_k$, and gluing these restrictions recovers the original marking. It is also surjective on $1$-cells, since every edge of $F^{-1}(C)$ is a $T$- or $B$-move, hence is supported on a single pair-of-pants component, and therefore is obtained by applying the corresponding edge in one factor $\mathcal M(S_i)$ while leaving the other factors fixed.

The gluing map is not injective on vertices, because isotopies of the glued surface $S$ are not required to preserve the chosen basepoints on the cut curves. In fact, the map \eqref{eq:gluing-map-for-markings} is the quotient of $$\mathcal{M}(S_1)\times \dots\times \mathcal{M}(S_k)$$ by an action of $\Z[C]$, the free abelian group generated by the curves of the cut system $C$. For each $c\in C$, the corresponding generator acts on a tuple
\[
(m_1,\dots,m_k)\in \mathcal M(S_1)\times\cdots\times \mathcal M(S_k)
\]
as follows. If $c$ bounds two distinct components $S_i$ and $S_j$, then
\[
c\cdot(m_1,\dots,m_k)
=
(m_1,\dots,D_c m_i,\dots,D_c^{-1}m_j,\dots,m_k).
\]
If $c$ corresponds to two boundary components $c_1,c_2$ of a single component $S_i$, then
\[
c\cdot(m_1,\dots,m_k)
=
(m_1,\dots,D_{c_1}D_{c_2}^{-1}m_i,\dots,m_k).
\]

It is simple to check that this is a well-defined action on the vertices of $\mathcal M(S_1)\times\cdots\times \mathcal M(S_k)$ which extends to a free and properly discontinuous action on the entire space. Then the quotient has fundamental group $\Z^C$ generated by loops of the type 
\[
(\id_{S_i}\cup_c T_c)\,(T_c^{-1}\cup_c \id_{S_j}),
\]
and similarly in the nonseparating case. Since these loops are precisely the boundaries of the 2-cells (Gluing relations for Dehn twists) in $F^{-1}(C)$, we conclude that $F^{-1}(C)$ is connected and simply connected.

It remains to verify conditions~$(3)$ and~$(4)$ of Proposition~\ref{prop: simply-connect domain from target}.


\emph{Checking condition~$(3)$:} this concerns comparing different lifts of edges in $\pantscpx(S)$. Recall that the edges of $\pantscpx(S)$ are given by $A$- and $S$-moves. We begin with the former.

Let $C\xra{e}C'$ be an $A$-move in $\pantscpx(S)$ exchanging the curve $a\in C$ with $a'\in C'$, and let $e_1$ and $e_2$ be two lifts of $e$. Then $e_1$ and $e_2$ must be edge-paths of the form $\gamma'\circ A_{a,a'}^{\pm 1} \circ \gamma$, where $\gamma$ and $\gamma'$ are edge-paths in $F^{-1}(C)$ and $F^{-1}(C')$, respectively. Since Condition~$(2)$ has already been verified, we may assume that 
    \[e_1:(C,m_1)\xra{}(C',m_1') \quad\text{and}\quad e_2:(C,m_2)\xra{}(C',m_2')\]
are single edges corresponding to $A$-moves. By the $2$-cell (Self-duality of associativity), we may assume without loss of generality that these are $A$-moves, rather than their inverses.  

For a fixed marking over the cut system $C$, there is a unique $A$-move exchanging $a\in C$. We now analyse how the markings $m_1$ and $m_2$ may differ.

By the definition of the $A$-move (Definition~\ref{def: complex of markings}), there exist diffeomorphisms $\phi_1,\phi_2:S'\to S_{\forkone}$ such that $\phi_1(a)=\phi_2(a)=c$, and both $\phi_1\circ m_1$ and $\phi_2\circ m_2$ define the standard marking on $S_{\forkone}$. Moreover, we have $\alpha\phi_1(a')=\alpha\phi_2(a')=c'$, and both $\alpha\circ\phi_1\circ m_1'$ and $\alpha\circ\phi_2\circ m_2'$ define the standard marking on $S_{\forktwo}$. By Lemma~\ref{lemma:aux-for-A-move}, we may further assume that $\phi_1$ and $\phi_2$ send the same boundary components of $S'$ to the components $\{1,2\}$ of $S_{\forkone}$.

It follows that $\phi_1\circ \phi_2^{-1}$ is a diffeomorphism of $S_{\forkone}$ preserving each component of the standard pants decomposition. In particular, it is isotopic to a diffeomorphism fixing the curve $c$ pointwise. Thus, we may regard $\phi_1\circ \phi_2^{-1}$ as a pair of diffeomorphisms of the two pairs of pants in $S'$, that fix the boundary components of the pants corresponding to $c$. From the description of the mapping class group of a pair of pants, it follows that $\phi_1\circ \phi_2^{-1}$ is a composition of Dehn twists along $c$ and the other boundary components, together with possible braiding morphisms that fix $c$.

On the other hand,
    \[ \alpha\phi_1(\alpha\phi_2)^{-1}=\alpha(\phi_1\phi_2^{-1})\alpha^{-1}\]
is a diffeomorphism of $S_{\forktwo}$ fixing $c'$. This implies that $\phi_1\circ \phi_2^{-1}$ also fixes $\alpha^{-1}(c')$.  In particular, by considering intersection numbers, we see that $\phi_1\circ \phi_2^{-1}$ cannot involve Dehn twists about $c$, since such twists would not preserve $\alpha^{-1}(c')$.

It follows that $\phi_1\circ \phi_2^{-1}$ is either a composition of Dehn twists along boundary components, or such a composition followed by the pair of braidings (one on each pair-of-pants component) represented on the right-hand side of Figure~\ref{fig:standard-markings-A-move}. Hence we may assume that $m_1$ and $m_2$ differ from one of the two markings in Figure~\ref{fig:standard-markings-A-move} only by Dehn twists along the boundary components.

% By analysing how braidings and Dehn twists affect the intersection number of $\alpha^{-1}(c')$ with its image, we conclude that applying $\phi_1\circ \phi_2^{-1}$ to the standard marking on $S_{\forkone}$ produces one of the two markings shown in Figure~\ref{fig:standard-markings-A-move}, possibly composed with Dehn twists along the boundary components. Hence we may assume that $m_1$ and $m_2$ differ from one of the two markings in Figure~\ref{fig:standard-markings-A-move} only by Dehn twists along the boundary components.

    \begin{figure}[ht]
        \centering
        \includegraphics[width=0.5\linewidth]{Figures/A-markings.pdf}
        \caption{Two standard markings on $S_{0,4}$}
        \label{fig:standard-markings-A-move}
    \end{figure}

If $m_1$ and $m_2$ are obtained from the same marking in Figure~\ref{fig:standard-markings-A-move} by Dehn twists along the boundary components, then the required commutativity of the square follows from the $2$-cell (Boundary twist relations).

It remains to consider the case where $m_1$ and $m_2$ arise from different markings in Figure~\ref{fig:standard-markings-A-move}. Without loss of generality, assume that $m_1$ and $m_2$ correspond to the left and right markings, respectively. Then Figure~\ref{fig:condition-3} gives the required commutative square.

    \begin{figure}
        \centering
        \includegraphics[width=0.65\linewidth]{Figures/square-for-condition-3.pdf}
        \caption{Commutative diagram checking condition (3)}
        \label{fig:condition-3}
    \end{figure}

A very similar argument works for the lifts of the S-move, and can be found in \cite[Section~7.7]{bk_marked_surfaces}. Here, the contractible square is obtained using the (Commutativity of disjoint moves) and the genus-one relation for $(g,n)=(1,1)$ given by \eqref{eq:g1.1}.

\emph{Checking condition~$(4)$:} now we must show that the boundary of each standard $2$-cell of $\pantscpx(S)$ lifts to a contractible loop in $\markS$. In fact, even more, we have that each such boundary can be lifted to the boundary of a standard $2$-cell in $\markS$.

The correspondence is as follows:
\begin{itemize}
    \item The boundary of the $2$-cell $(3A)$ can be lifted to the boundary of the $2$-cell corresponding to the hexagon relations; see Figure~\ref{fig:hexagons};
    \item The boundary of the $2$-cell $(5A)$ can be lifted to the boundary of the $2$-cell corresponding to the pentagon relation; see Figure~\ref{fig:pentagon};
    \item The boundary of the $2$-cell $(3S)$ can be lifted to the boundary of the $2$-cell corresponding to the Genus-one relation \eqref{eq:g1.2}; see Figure~\ref{fig:G2-cell-G1-2};
    \item The boundary of the $2$-cell $(6AS)$ can be lifted to the boundary of the $2$-cell corresponding to the Genus-one relation for $(g,n)=(1,2)$; see Figure~\ref{fig:G2-cell-NS-complex-of-markings};
    \item The boundary of the $2$-cell $(C)$ can be lifted to the boundary of the $2$-cell corresponding to the relation (Commutativity of disjoint moves).
\end{itemize}

Thus the boundary of each $2$-cell of $\pantscpx(S)$ lifts to a loop bounding a $2$-cell of $\markS$. Hence condition~$(4)$ holds. Proposition~\ref{prop: simply-connect domain from target} now implies that $\markS$ is connected and simply connected.
\end{proof}






\begin{lemma}\label{lemma:aux-for-A-move}
    Let $\phi:S'\to S_{\forkone}$ be a diffeomorphism witnessing an $A$-move, as in Definition~\ref{def: complex of markings}, and let $\{c,c'\}$ denote the boundary components of $S'$ that are mapped to the boundary components $\{1,2\}$ of $S_{\forkone}$. Then there exists a diffeomorphism $\phi':S'\to S_{\forkone}$ witnessing the same $A$-move such that the boundary components $\{c,c'\}$ are mapped to $\{0,3\}$.
\end{lemma}

\begin{proof}
    Let $(C,m)\xra{A_{a,b}}(C',m')$ be an $A$-move on $\mathcal{M}(S)$ and let $m|_{S'}$ and $m'_{S'}$ denote representatives of the markings $m$ and $m'$ restricted to $S'$ that, together with $\phi$, witnesses the $A$-move. 

    Let $z:S_{\forkone}\to S_{\forkone}$ denote a difffeomorphism that fixes the standard marking and on the boundary components gives the permutation $(13)(20)$ (this can be constructed using the standard marking and the diagrammatic representation of mapping classes).
    
    Define the diffeomorphism
        \[\phi'\coloneqq z_S\circ \phi.\]
    By construction,  $\phi'\circ m|_{S'}$ is still the standard marking on $S_{\forkone}$, and $\alpha\circ\phi'\circ m'_{S'}$ is still the standard marking on $S_{\forktwo}$. Hence $\phi'$ also witnesses this $A$-move and, by construction, takes the boundary components $\{c,c'\}$ to the boundaries $\{0,3\}$ on $S_{\forkone}$.
    % 
    % Let $z:S_{\forkone}\to S_{\forkone}$ be the standard diffeomorphism which rotates the surface and swaps the boundary components $\{1,2\}$ with $\{0,3\}$ preserving the standard marking (this can be constructed explicitly using the diagrammatic representation of mapping classes). Define the new diffeomorphism by
    %     \[\phi'=z\circ \phi.\]
    % Since $z$ preserves the standard marking, we still have that
    %     \[\phi'(a)=c \quad \text{and} \quad \phi'\circ m= m_{\forkone}.\]
    % All that remains is to check that
    %     \[\alpha\circ\phi'(b)=c' \quad \text{and} \quad \alpha\circ\phi'\circ m'= m_{\forktwo}.\]
    % Substituting, this becomes
    %     \[\alpha\circ z\circ \phi(b)=c' \quad \text{and}\quad \alpha\circ z\circ\phi\circ m'= m_{\forktwo}.\]
\end{proof}
%\input{Sections/appendix_adjunctions}
%\input{Sections/appendix_model_cat}



\bibliographystyle{plain}
\bibliography{bibliography.bib}
\end{document}
